\section{형식적 스킴}
\label{section:I.10-ko}

\subsection{아핀 형식적 스킴.}
\label{subsection:I.10.1-ko}

\begin{env}[10.1.1]
\label{I.10.1.1-ko}
$A$를 위상환인 \emph{허용환}이라 하자
(\hyperref[0.7.1.2-ko]{0, 7.1.2}). $A$의 모든 정의 아이디얼
$\mathfrak{J}$에 대하여 $\operatorname{Spec}(A/\mathfrak{J})$는
$\operatorname{Spec}(A)$의 닫힌 부분공간 $V(\mathfrak{J})$와
동일시된다(\hyperref[I.1.1.11-ko]{1.1.11}). 이는 $A$의
\emph{열린} 소 아이디얼들의 집합이다. 이 위상공간은 고려한
\oldpage[I]{181}
정의 아이디얼 $\mathfrak{J}$의 선택과 무관하며, 이를
$\mathfrak{X}$로 나타내자. $(\mathfrak{J}_\lambda)$를 정의
아이디얼들로 이루어진 $A$에서 0의 근방들의 기본계라 하고, 모든
$\lambda$에 대하여 $\mathscr{O}_\lambda$를
$\operatorname{Spec}(A/\mathfrak{J}_\lambda)$의 구조층이라 하자. 이 층은
$\widetilde{A}/\widetilde{\mathfrak{J}}_\lambda$에 의해
$\mathfrak{X}$ 위에 유도된다(이 몫층은 $\mathfrak{X}$ 밖에서 영이다).

$\mathfrak{J}_\mu\subset\mathfrak{J}_\lambda$이면 표준 준동형
$A/\mathfrak{J}_\mu\to A/\mathfrak{J}_\lambda$는 환의 층의 준동형
$u_{\lambda\mu}:\mathscr{O}_\mu\to\mathscr{O}_\lambda$를 정의하므로
(\hyperref[I.1.6.1-ko]{1.6.1}), $(\mathscr{O}_\lambda)$는 이 준동형들에
관하여 \emph{환의 층들의 사영계}를 이룬다. $\mathfrak{X}$의 위상에는
준콤팩트 열린집합들로 이루어진 기저가 있으므로, 각
$\mathscr{O}_\lambda$에 그 바탕 환의 층(위상을 잊은 층)이
$\mathscr{O}_\lambda$인 \emph{의사이산 위상환의 층}을 대응시킬 수
있다(\hyperref[0.3.8.1-ko]{0, 3.8.1}). 이 층도 여전히
$\mathscr{O}_\lambda$로 나타내겠다. 또한 $\mathscr{O}_\lambda$들은
여전히 \emph{위상환의 층들의 사영계}를 이룬다
(\hyperref[0.3.8.2-ko]{0, 3.8.2}). 이 계
$(\mathscr{O}_\lambda)$의 사영극한인 $\mathfrak{X}$ 위의
\emph{위상환의 층}을 $\mathscr{O}_{\mathfrak{X}}$로 나타내겠다. 따라서
$\mathfrak{X}$의 모든 \emph{준콤팩트} 열린집합 $U$에 대하여
$\Gamma(U,\mathscr{O}_{\mathfrak{X}})$는 \emph{이산}환들의 계
$\Gamma(U,\mathscr{O}_\lambda)$의 사영극한인 위상환이다
(\hyperref[0.3.2.6-ko]{0, 3.2.6}).
\end{env}

\begin{definition}[10.1.2]
\label{I.10.1.2-ko}
허용환 $A$가 주어졌을 때, $A$의 열린 소 아이디얼들로 이루어진
$\operatorname{Spec}(A)$의 닫힌 부분공간 $\mathfrak{X}$를 $A$의
형식적 스펙트럼이라 하고 $\operatorname{Spf}(A)$로 나타낸다. 위상환 달린
공간이 형식적 스펙트럼 $\operatorname{Spf}(A)=\mathfrak{X}$에 다음의
구조층을 갖춘 것과 동형이면 이를 아핀 형식적 스킴이라고 한다. 여기서
$\mathfrak{J}_\lambda$가 $A$의 정의 아이디얼들의 여과집합을 달릴 때
$\mathscr{O}_{\mathfrak{X}}$는
의사이산 위상환의 층
$(\widetilde{A}/\widetilde{\mathfrak{J}}_\lambda)|\mathfrak{X}$들의
사영극한이다.
\end{definition}

\emph{형식적 스펙트럼} $\mathfrak{X}=\operatorname{Spf}(A)$를 아핀 형식적
스킴으로 말할 때에는 언제나 위에서 정의한
$\mathscr{O}_{\mathfrak{X}}$를 갖는 위상환 달린 공간
$(\mathfrak{X},\mathscr{O}_{\mathfrak{X}})$을 뜻한다.

모든 \emph{아핀 스킴} $X=\operatorname{Spec}(A)$는 $A$를 이산
위상환으로 보아 유일한 방식으로 아핀 형식적 스킴으로 볼 수 있다. 이때
$U$가 준콤팩트이면 위상환 $\Gamma(U,\mathscr{O}_X)$는 이산이지만,
$U$가 $X$의 임의의 열린집합일 때에는 일반적으로 그렇지 않다.

\begin{proposition}[10.1.3]
\label{I.10.1.3-ko}
$A$가 허용환이고 $\mathfrak{X}=\operatorname{Spf}(A)$이면,
$\Gamma(\mathfrak{X},\mathscr{O}_{\mathfrak{X}})$는 $A$와 위상적으로
동형이다.
\end{proposition}

\begin{proof}
실제로 $\mathfrak{X}$는 $\operatorname{Spec}(A)$에서 닫혀 있으므로
준콤팩트이다. 따라서
$\Gamma(\mathfrak{X},\mathscr{O}_{\mathfrak{X}})$는 이산환
$\Gamma(\mathfrak{X},\mathscr{O}_\lambda)$들의 사영극한과 위상적으로
동형이다. 그런데 $\Gamma(\mathfrak{X},\mathscr{O}_\lambda)$는
$A/\mathfrak{J}_\lambda$와 동형이다
(\hyperref[I.1.3.7-ko]{1.3.7}). $A$는 분리이고 완비이므로
$\varprojlim A/\mathfrak{J}_\lambda$와 위상적으로 동형이다
(\hyperref[0.7.2.1-ko]{0, 7.2.1}). 이로부터 명제가 따른다.
\end{proof}

\begin{proposition}[10.1.4]
\label{I.10.1.4-ko}
$A$를 허용환, $\mathfrak{X}=\operatorname{Spf}(A)$라 하고, 모든
$f\in A$에 대하여 $\mathfrak{D}(f)=D(f)\cap\mathfrak{X}$라 하자. 그러면
위상환 달린 공간
$(\mathfrak{D}(f),\mathscr{O}_{\mathfrak{X}}|\mathfrak{D}(f))$는 아핀
형식적 스펙트럼 $\operatorname{Spf}(A_{\{f\}})$와 동형이다
(\hyperref[0.7.6.15-ko]{0, 7.6.15}).
\end{proposition}

\begin{proof}
$A$의 모든 정의 아이디얼 $\mathfrak{J}$에 대하여 이산환
$S_f^{-1}A/S_f^{-1}\mathfrak{J}$는
$A_{\{f\}}/\mathfrak{J}_{\{f\}}$와 표준적으로 동일시된다
(\hyperref[0.7.6.9-ko]{0, 7.6.9}). 따라서
(\hyperref[I.1.2.5-ko]{1.2.5}와 \hyperref[I.1.2.6-ko]{1.2.6})에 따라 위상공간
$\operatorname{Spf}(A_{\{f\}})$는 $\mathfrak{D}(f)$와 표준적으로
동일시된다. 더 나아가 $\mathfrak{D}(f)$에 포함되는
$\mathfrak{X}$의 모든 준콤팩트 열린집합 $U$에 대하여
$\Gamma(U,\mathscr{O}_\lambda)$는
$\operatorname{Spec}(S_f^{-1}A/S_f^{-1}\mathfrak{J}_\lambda)$의 구조층의
$U$ 위의 단면가군과 동일시된다
(\hyperref[I.1.3.6-ko]{1.3.6}). 따라서
$\mathfrak{Y}=\operatorname{Spf}(A_{\{f\}})$로 놓으면
$\Gamma(U,\mathscr{O}_{\mathfrak{X}})$는 단면가군
$\Gamma(U,\mathscr{O}_{\mathfrak{Y}})$와 동일시된다. 이로써 명제가
증명된다.
\end{proof}

\oldpage[I]{182}
\begin{env}[10.1.5]
\label{I.10.1.5-ko}
\emph{위상을 잊은} 환의 층으로서 $\operatorname{Spf}(A)$의 구조층
$\mathscr{O}_{\mathfrak{X}}$는 모든 $x\in\mathfrak{X}$에서
(\hyperref[I.10.1.4-ko]{10.1.4})에 따라 귀납극한
$\varinjlim_{f\notin\mathfrak{j}_x}A_{\{f\}}$와 동일시되는 줄기를 갖는다.
따라서 (\hyperref[0.7.6.17-ko]{0, 7.6.17}과
\hyperref[0.7.6.18-ko]{7.6.18}):
\end{env}

\begin{proposition}[10.1.6]
\label{I.10.1.6-ko}
모든 $x\in\mathfrak{X}=\operatorname{Spf}(A)$에 대하여 줄기
$\mathscr{O}_x$는 국소환이고, 그 잉여체는
$k(x)=A_x/\mathfrak{j}_xA_x$와 동형이다. 더 나아가 $A$가 아딕이고
뇌터이면 $\mathscr{O}_x$는 뇌터 환이다.
\end{proposition}

$k(x)$가 영환이 아니므로, 이 결과로부터 환의 층
$\mathscr{O}_{\mathfrak{X}}$의 \emph{지지집합}은
\emph{$\mathfrak{X}$와 같다}는 결론을 얻는다.

\subsection{아핀 형식적 스킴의 사상.}
\label{subsection:I.10.2-ko}

\begin{env}[10.2.1]
\label{I.10.2.1-ko}
$A$, $B$를 두 허용환이라 하고,
$\varphi:B\to A$를 \emph{연속} 준동형이라 하자. 그러면 연속사상
$ {}^a\varphi:\operatorname{Spec}(A)\to\operatorname{Spec}(B)$
(\hyperref[I.1.2.1-ko]{1.2.1})은
$\mathfrak{X}=\operatorname{Spf}(A)$를
$\mathfrak{Y}=\operatorname{Spf}(B)$ 안으로 보낸다. 실제로 $A$의 열린
소 아이디얼의 $\varphi$에 의한 역상은 $B$의 열린 소 아이디얼이다.
한편 모든 $g\in B$에 대하여 $\varphi$는 연속 준동형
\[
  \Gamma(\mathfrak{D}(g),\mathscr{O}_{\mathfrak{Y}})
  \longrightarrow
  \Gamma(\mathfrak{D}(\varphi(g)),\mathscr{O}_{\mathfrak{X}})
\]
을 정의한다. 이는 (\hyperref[I.10.1.4-ko]{10.1.4}),
(\hyperref[I.10.1.3-ko]{10.1.3}),
(\hyperref[0.7.6.7-ko]{0, 7.6.7})에 따른다. 이 준동형들은 $g$를 $g$의
배수로 바꿀 때의 제한사상들과 양립하고,
$\mathfrak{D}(\varphi(g))={}^a\varphi^{-1}(\mathfrak{D}(g))$이므로,
위상환의 층 사이의 \emph{연속} 준동형
$\mathscr{O}_{\mathfrak{Y}}\to
{}^a\varphi_*(\mathscr{O}_{\mathfrak{X}})$
(\hyperref[0.3.2.5-ko]{0, 3.2.5})을 정의한다. 이를 다시
$\widetilde{\varphi}$로 나타내겠다. 이렇게 하여 위상환 달린 공간의 사상
$\Phi=({}^a\varphi,\widetilde{\varphi})$를
$\mathfrak{X}\to\mathfrak{Y}$로 얻는다. 위상을 잊은 환의 층의
준동형으로서 $\widetilde{\varphi}$는 모든 $x\in\mathfrak{X}$에서 줄기
사이의 준동형
$\widetilde{\varphi}_x^\sharp:
\mathscr{O}_{{}^a\varphi(x)}\to\mathscr{O}_x$를 정의한다.
\end{env}

\begin{proposition}[10.2.2]
\label{I.10.2.2-ko}
$A$, $B$를 두 허용환이라 하고,
$\mathfrak{X}=\operatorname{Spf}(A)$,
$\mathfrak{Y}=\operatorname{Spf}(B)$라 하자. 위상환 달린 공간의 사상
$u=(\psi,\theta):\mathfrak{X}\to\mathfrak{Y}$가 $B\to A$인 어떤 연속 환 준동형
$\varphi$로부터 얻어지는
$({}^a\varphi,\widetilde{\varphi})$ 꼴이기 위한 필요충분조건은 모든
$x\in\mathfrak{X}$에 대하여 $\theta_x^\sharp$가 국소 준동형
$\mathscr{O}_{\psi(x)}\to\mathscr{O}_x$인 것이다.
\end{proposition}

\begin{proof}
이 조건은 필요하다. 실제로
$\mathfrak{p}=\mathfrak{j}_x\in\operatorname{Spf}(A)$,
$\mathfrak{q}=\varphi^{-1}(\mathfrak{j}_x)$라 하자.
$g\notin\mathfrak{q}$이면 $\varphi(g)\notin\mathfrak{p}$이고,
$\varphi$에서 얻어지는 준동형
$B_{\{g\}}\to A_{\{\varphi(g)\}}$
(\hyperref[0.7.6.7-ko]{0, 7.6.7})은
$\mathfrak{q}_{\{g\}}$를
$\mathfrak{p}_{\{\varphi(g)\}}$의 부분집합으로 보낸다. 따라서 귀납극한을
취하고 (\hyperref[I.10.1.5-ko]{10.1.5})와
(\hyperref[0.7.6.17-ko]{0, 7.6.17})을 고려하면
$\widetilde{\varphi}_x^\sharp$가 국소 준동형임을 알 수 있다.

역으로 $(\psi,\theta)$가 명제의 조건을 만족하는 사상이라고 하자.
(\hyperref[I.10.1.3-ko]{10.1.3})에 따라 $\theta$는 다음 연속 환 준동형을
정의한다:
\[
  \varphi=\Gamma(\theta):
  B=\Gamma(\mathfrak{Y},\mathscr{O}_{\mathfrak{Y}})
  \longrightarrow
  \Gamma(\mathfrak{X},\mathscr{O}_{\mathfrak{X}})=A.
\]
$\theta$에 대한 가정에 의하여 $\mathfrak{X}$ 위의
$\mathscr{O}_{\mathfrak{X}}$의 단면 $\varphi(g)$의 $x$에서의 싹이
가역일 필요충분조건은 $g$의 $\psi(x)$에서의 싹이 가역인 것이다. 그런데
(\hyperref[0.7.6.17-ko]{0, 7.6.17})에 따라
$\mathfrak{X}$ 위의 $\mathscr{O}_{\mathfrak{X}}$의 단면(각각
$\mathfrak{Y}$ 위의 $\mathscr{O}_{\mathfrak{Y}}$의 단면) 가운데
$x$에서의 싹(각각 $\psi(x)$에서의 싹)이 가역이 아닌 것들은 정확히
$\mathfrak{j}_x$
\oldpage[I]{183}
(각각 $\mathfrak{j}_{\psi(x)}$)의 원소들이다. 따라서 앞의 관찰로부터
${}^a\varphi=\psi$를 얻는다. 마지막으로 모든 $g\in B$에 대하여 도식
\[
  \xymatrix{
    B=\Gamma(\mathfrak{Y},\mathscr{O}_{\mathfrak{Y}})
      \ar[r]^\varphi\ar[d] &
    \Gamma(\mathfrak{X},\mathscr{O}_{\mathfrak{X}})=A\ar[d]\\
    B_{\{g\}}=\Gamma(\mathfrak{D}(g),\mathscr{O}_{\mathfrak{Y}})
      \ar[r]_{\Gamma(\theta_{\mathfrak{D}(g)})} &
    \Gamma(\mathfrak{D}(\varphi(g)),\mathscr{O}_{\mathfrak{X}})
      =A_{\{\varphi(g)\}}
  }
\]
은 가환이다. 완비 분수환의 보편 성질
(\hyperref[0.7.6.6-ko]{0, 7.6.6})에 따라 모든 $g\in B$에 대하여
$\theta_{\mathfrak{D}(g)}$는
$\widetilde{\varphi}_{\mathfrak{D}(g)}$와 같고, 따라서
(\hyperref[0.3.2.5-ko]{0, 3.2.5})에 따라
$\theta=\widetilde{\varphi}$이다.
\end{proof}

(\hyperref[I.10.2.2-ko]{10.2.2})의 조건을 만족하는 위상환 달린 공간의
사상 $(\psi,\theta)$를 \emph{아핀 형식적 스킴의 사상}이라 하겠다.
$A$에 $\operatorname{Spf}(A)$를 대응시키는 함자와 $\mathfrak{X}$에
$\Gamma(\mathfrak{X},\mathscr{O}_{\mathfrak{X}})$를 대응시키는 함자는
허용환의 범주와 아핀 형식적 스킴의 범주의 반대 범주 사이의
\emph{동치}를 정의한다고 할 수 있다(T, I, 1.2).

\begin{env}[10.2.3]
\label{I.10.2.3-ko}
(\hyperref[I.10.2.2-ko]{10.2.2})의 특수한 경우로서, $f\in A$에 대하여
$\mathfrak{X}$가 $\mathfrak{D}(f)$ 위에 유도하는 아핀 형식적 스킴의
표준 포함 사상은 표준 연속 준동형 $A\to A_{\{f\}}$에 대응한다. 이제
(\hyperref[I.10.2.2-ko]{10.2.2})의 가정 아래에서 $h$를 $B$의 원소,
$g$를 $\varphi(h)$의 배수인 $A$의 원소라 하자. 그러면
$\psi(\mathfrak{D}(g))\subset\mathfrak{D}(h)$이다. $u$를
$\mathfrak{D}(g)$에 제한한 사상을 $\mathfrak{D}(g)$에서
$\mathfrak{D}(h)$로 가는 사상으로 보면, 이는 도식
\[
  \xymatrix{
    \mathfrak{D}(g)\ar[r]^v\ar[d] &
    \mathfrak{D}(h)\ar[d]\\
    \mathfrak{X}\ar[r]_u &
    \mathfrak{Y}
  }
\]
을 가환이게 하는 유일한 사상 $v$이다. 이 사상은 도식
\[
  \xymatrix{
    A\ar[d] &
    B\ar[l]_\varphi\ar[d]\\
    A_{\{g\}} &
    B_{\{h\}}\ar[l]_{\varphi'}
  }
\]
을 가환이게 하는 유일한 연속 준동형
$\varphi':B_{\{h\}}\to A_{\{g\}}$
(\hyperref[0.7.6.7-ko]{0, 7.6.7})에 대응한다.
\end{env}

\subsection{아핀 형식적 스킴의 정의 아이디얼.}
\label{subsection:I.10.3-ko}

\begin{env}[10.3.1]
\label{I.10.3.1-ko}
$A$를 허용환, $\mathfrak{J}$를 $A$의 열린 아이디얼,
$\mathfrak{X}$를 아핀 형식적 스킴 $\operatorname{Spf}(A)$라 하자.
$(\mathfrak{J}_\lambda)$를 $\mathfrak{J}$에 포함되는 $A$의 정의
아이디얼들의 집합이라 하자. 그러면
$\widetilde{\mathfrak{J}}/\widetilde{\mathfrak{J}}_\lambda$는
$\widetilde{A}/\widetilde{\mathfrak{J}}_\lambda$의 아이디얼층이다.
$\mathfrak{X}$ 위에
$\widetilde{\mathfrak{J}}/\widetilde{\mathfrak{J}}_\lambda$가 유도하는
층들의 사영극한을 $\mathfrak{J}^\Delta$로 나타내자. 이는
$\mathscr{O}_{\mathfrak{X}}$의 \emph{아이디얼층}과 동일시된다
(\hyperref[0.3.2.6-ko]{0, 3.2.6}). 모든 $f\in A$에 대하여
$\Gamma(\mathfrak{D}(f),\mathfrak{J}^\Delta)$는
$S_f^{-1}\mathfrak{J}/S_f^{-1}\mathfrak{J}_\lambda$들의 사영극한, 즉
환 $A_{\{f\}}$의 열린 아이디얼 $\mathfrak{J}_{\{f\}}$와 동일시된다
(\hyperref[0.7.6.9-ko]{0, 7.6.9}). 특히
$\Gamma(\mathfrak{X},\mathfrak{J}^\Delta)=\mathfrak{J}$이다.
$\mathfrak{D}(f)$들이 $\mathfrak{X}$의 위상의 기저를 이루므로 이로부터
다음을 얻는다:
\[
  \mathfrak{J}^\Delta|\mathfrak{D}(f)
  =(\mathfrak{J}_{\{f\}})^\Delta.
  \tag{10.3.1.1}\label{I.10.3.1.1-ko}
\]
\end{env}

\oldpage[I]{184}
\begin{env}[10.3.2]
\label{I.10.3.2-ko}
(\hyperref[I.10.3.1-ko]{10.3.1})의 표기 아래에서, 모든 $f\in A$에
대하여
$A_{\{f\}}=\Gamma(\mathfrak{D}(f),\mathscr{O}_{\mathfrak{X}})$에서
$\Gamma\bigl(\mathfrak{D}(f),
(\widetilde{A}/\widetilde{\mathfrak{J}})|\mathfrak{X}\bigr)
=S_f^{-1}A/S_f^{-1}\mathfrak{J}$로 가는 표준 사상은 \emph{전사}이고,
그 핵은
$\Gamma(\mathfrak{D}(f),\mathfrak{J}^\Delta)=\mathfrak{J}_{\{f\}}$이다
(\hyperref[0.7.6.9-ko]{0, 7.6.9}). 따라서 이 사상들은 위상환의 층
$\mathscr{O}_{\mathfrak{X}}$에서 이산환의 층
$(\widetilde{A}/\widetilde{\mathfrak{J}})|\mathfrak{X}$로 가는 연속
\emph{전사} 준동형을 정의하며, 이를 \emph{표준}이라 한다. 그 핵은
$\mathfrak{J}^\Delta$이다. 이 준동형은
$\varphi$가 연속 준동형 $A\to A/\mathfrak{J}$일 때의
$\widetilde{\varphi}$
(\hyperref[I.10.2.1-ko]{10.2.1})와 다르지 않다. 아핀 형식적 스킴의 사상
$({}^a\varphi,\widetilde{\varphi}):
\operatorname{Spec}(A/\mathfrak{J})\to\mathfrak{X}$
(여기서 ${}^a\varphi$는 닫힌 부분공간
$V(\mathfrak{J})\subset\mathfrak{X}$로의 표준 위상동형사상에 그
포함사상을 합성한 것이다)도 \emph{표준}이라 한다.\footnote{\textbf{번역 주.} 인쇄 원문은
${}^a\varphi$가 $\mathfrak{X}$에서 자기 자신으로 가는 항등
위상동형사상이라고 쓰지만, 이는 $\mathfrak{J}$가 임의의 열린 아이디얼일
때 일반적으로 성립하지 않는다. 예를 들어 체 $k$에 대하여 이산 허용환
$A=k\times k$와 열린 아이디얼 $\mathfrak{J}=k\times\{0\}$에 대하여
$\operatorname{Spf}(A)$는 두 점이지만
$\operatorname{Spec}(A/\mathfrak{J})$는 한 점이다. 올바른 상은
$V(\mathfrak{J})\subset\mathfrak{X}$이며, $\mathfrak{J}$가 정의
아이디얼일 때에는 이 상이 $\mathfrak{X}$ 전체이다. 이는 공식 정오표에
실린 항목은 아니며, 위와 같이 형에 맞게 바로잡았다.} 따라서 앞의
결과에 의하여 \emph{표준 동형사상}
\[
  \mathscr{O}_{\mathfrak{X}}/\mathfrak{J}^\Delta
  \xrightarrow{\sim}
  (\widetilde{A}/\widetilde{\mathfrak{J}})|\mathfrak{X}.
  \tag{10.3.2.1}\label{I.10.3.2.1-ko}
\]
을 얻는다.

$\Gamma(\mathfrak{X},\mathfrak{J}^\Delta)=\mathfrak{J}$이므로
$\mathfrak{J}\to\mathfrak{J}^\Delta$가
\emph{포함관계를 엄격히 보존한다}는 것은 명백하다. 앞의 결과에 따라
$\mathfrak{J}\subset\mathfrak{J}'$이면 층
${\mathfrak{J}'}^\Delta/\mathfrak{J}^\Delta$는
$\widetilde{\mathfrak{J}'}/\widetilde{\mathfrak{J}}
=\widetilde{\mathfrak{J}'/\mathfrak{J}}$와 표준적으로 동형이다.
\end{env}

\begin{env}[10.3.3]
\label{I.10.3.3-ko}
가정과 표기는 계속
(\hyperref[I.10.3.1-ko]{10.3.1})의 것이라 하자.
$\mathscr{O}_{\mathfrak{X}}$의 아이디얼층 $\mathscr{J}$가 다음 조건을
만족할 때 이를 \emph{$\mathfrak{X}$의 정의 아이디얼층}(또는
\emph{$\mathfrak{X}$의 정의 아이디얼})이라 하겠다. 모든
$x\in\mathfrak{X}$에 대하여 $x$의 열린 근방 $\mathfrak{D}(f)$가
존재하고, 여기서 $f\in A$이며,
$\mathscr{J}|\mathfrak{D}(f)$는 $\mathfrak{H}^\Delta$ 꼴이다. 이때
$\mathfrak{H}$는 $A_{\{f\}}$의 정의 아이디얼이다.
\end{env}

\begin{proposition}[10.3.4]
\label{I.10.3.4-ko}
모든 $f\in A$에 대하여 $\mathfrak{X}$의 모든 정의 아이디얼은
$\mathfrak{D}(f)$의 정의 아이디얼을 유도한다.
\end{proposition}

\begin{proof}
이는 (\hyperref[I.10.3.1.1-ko]{10.3.1.1})에서 따른다.
\end{proof}

\begin{proposition}[10.3.5]
\label{I.10.3.5-ko}
$A$가 허용환이면 $\mathfrak{X}=\operatorname{Spf}(A)$의 모든 정의
아이디얼은 $A$의 유일하게 결정되는 정의 아이디얼 $\mathfrak{J}$에 대한
$\mathfrak{J}^\Delta$ 꼴이다.
\end{proposition}

\begin{proof}
실제로 $\mathscr{J}$를 $\mathfrak{X}$의 정의 아이디얼이라 하자.
가정과 $\mathfrak{X}$의 준콤팩트성에 따라 유한 개의 원소
$f_i\in A$가 존재하여 $\mathfrak{D}(f_i)$들이 $\mathfrak{X}$를 덮고,
$\mathscr{J}|\mathfrak{D}(f_i)=\mathfrak{H}_i^\Delta$이다. 여기서
$\mathfrak{H}_i$는 $A_{\{f_i\}}$의 정의 아이디얼이다. 모든 $i$에
대하여 $(\mathfrak{K}_i)_{\{f_i\}}=\mathfrak{H}_i$인 $A$의 열린
아이디얼 $\mathfrak{K}_i$가 존재한다
(\hyperref[0.7.6.9-ko]{0, 7.6.9}). 모든 $\mathfrak{K}_i$에 포함되는
$A$의 정의 아이디얼을 $\mathfrak{K}$라 하자.
$\mathscr{J}/\mathfrak{K}^\Delta$의
$\operatorname{Spec}(A/\mathfrak{K})$의 구조층
$\widetilde{A/\mathfrak{K}}$ 안에서의 표준 상
(\hyperref[I.10.3.2-ko]{10.3.2})은 $\mathfrak{D}(f_i)$ 위로 제한하면
$\widetilde{\mathfrak{K}_i/\mathfrak{K}}$의 제한과 같다. 따라서 이
표준 상은 $\operatorname{Spec}(A/\mathfrak{K})$ 위의 준연접층이고,
$\mathfrak{K}$를 포함하는 $A$의 아이디얼 $\mathfrak{J}$에 대하여
$\widetilde{\mathfrak{J}/\mathfrak{K}}$ 꼴이다
(\hyperref[I.1.4.1-ko]{1.4.1}). 그러므로
$\mathscr{J}=\mathfrak{J}^\Delta$이다
(\hyperref[I.10.3.2-ko]{10.3.2}). 또한 모든 $i$에 대하여
$\mathfrak{H}_i^{n_i}\subset\mathfrak{K}_{\{f_i\}}$인 정수 $n_i$가
존재한다. $n$을 $n_i$들의 최댓값이라 하면
$(\mathscr{J}/\mathfrak{K}^\Delta)^n=0$이고, 따라서
(\hyperref[I.10.3.2-ko]{10.3.2})에 의하여
$(\widetilde{\mathfrak{J}/\mathfrak{K}})^n=0$이다. 그러므로 마침내
$(\mathfrak{J}/\mathfrak{K})^n=0$이고
(\hyperref[I.1.3.13-ko]{1.3.13}), 이는 $\mathfrak{J}$가 $A$의 정의
아이디얼임을 보인다 (\hyperref[0.7.1.4-ko]{0, 7.1.4}).
\end{proof}

\begin{proposition}[10.3.6]
\label{I.10.3.6-ko}
$A$를 아딕환, $\mathfrak{J}$를 $A$의 정의 아이디얼이라 하고,
$\mathfrak{J}/\mathfrak{J}^2$가 유한 생성
$A/\mathfrak{J}$-가군이라고 하자. 그러면 모든 정수 $n>0$에 대하여
$(\mathfrak{J}^\Delta)^n=(\mathfrak{J}^n)^\Delta$이다.
\end{proposition}

\begin{proof}
실제로 모든 $f\in A$에 대하여 $\mathfrak{J}^n$은 열린 아이디얼이므로
\[
  \bigl(\Gamma(\mathfrak{D}(f),\mathfrak{J}^\Delta)\bigr)^n
  =(\mathfrak{J}_{\{f\}})^n
  =(\mathfrak{J}^n)_{\{f\}}
  =\Gamma\bigl(\mathfrak{D}(f^n),(\mathfrak{J}^n)^\Delta\bigr)
\]
\oldpage[I]{185}
이다. 이는 (\hyperref[I.10.3.1.1-ko]{10.3.1.1})과
(\hyperref[0.7.6.12-ko]{0, 7.6.12})에서 따른다.
$(\mathfrak{J}^\Delta)^n$은 준층
$U\mapsto(\Gamma(U,\mathfrak{J}^\Delta))^n$의 층화이므로
(\hyperref[0.4.1.6-ko]{0, 4.1.6}), $\mathfrak{D}(f)$들이
$\mathfrak{X}$의 위상의 기저를 이룬다는 사실로부터 결론이
따른다.\footnote{\textbf{번역 주.} 인쇄 원문은 이 명제의 증명에서
“따름정리가 따른다”라고 쓰지만, 이곳에는 따름정리가 없으므로 “결론이
따른다”로 바로잡았다.}
\end{proof}

\begin{env}[10.3.7]
\label{I.10.3.7-ko}
$\mathfrak{X}$의 정의 아이디얼들의 족 $(\mathscr{J}_\lambda)$가 다음
조건을 만족할 때 이를 \emph{정의 아이디얼들의 기본계}라고 한다.
$\mathfrak{X}$의 모든 정의 아이디얼은 어떤 $\mathscr{J}_\lambda$를
포함한다. $\mathscr{J}_\lambda=\mathfrak{J}_\lambda^\Delta$이므로 이는
$\mathfrak{J}_\lambda$들이 $A$에서 \emph{0의 근방 기본계}를 이룬다는
것과 같다. $\mathfrak{D}(f_\alpha)$들이 $\mathfrak{X}$를 덮도록
$A$의 원소들의 족 $(f_\alpha)$를 잡자.
$(\mathscr{J}_\lambda)$가 $\mathscr{O}_{\mathfrak{X}}$의 아이디얼들의
감소 여과족이고, 모든 $\alpha$에 대하여
$(\mathscr{J}_\lambda|\mathfrak{D}(f_\alpha))$가
$\mathfrak{D}(f_\alpha)$의 정의 아이디얼들의 기본계를 이룬다면,
$(\mathscr{J}_\lambda)$는 $\mathfrak{X}$의 정의 아이디얼들의 기본계이다.
실제로 $\mathfrak{X}$의 모든 정의 아이디얼 $\mathscr{J}$에 대하여
$\mathfrak{D}(f_i)$들로 이루어진 $\mathfrak{X}$의 유한 덮개가 존재하여,
모든 $i$에 대하여
$\mathscr{J}_{\lambda_i}|\mathfrak{D}(f_i)$는
$\mathfrak{D}(f_i)$의 정의 아이디얼이고
$\mathscr{J}|\mathfrak{D}(f_i)$에 포함된다. 모든 $i$에 대하여
$\mathscr{J}_\mu\subset\mathscr{J}_{\lambda_i}$인 지표 $\mu$를 잡으면,
(\hyperref[I.10.3.3-ko]{10.3.3})에 따라 $\mathscr{J}_\mu$는
$\mathfrak{X}$의 정의 아이디얼이고, 명백히 $\mathscr{J}$에 포함된다.
이로부터 주장이 따른다.
\end{env}

\subsection{형식적 준스킴과 형식적 준스킴의 사상.}
\label{subsection:I.10.4-ko}

\begin{env}[10.4.1]
\label{I.10.4.1-ko}
위상환 달린 공간 $\mathfrak{X}$가 주어졌다고 하자. 열린집합
$U\subset\mathfrak{X}$ 위에 $\mathfrak{X}$가 유도하는 위상환 달린 공간이
아핀 형식적 스킴(각각 그 환이 아딕인 아핀 형식적 스킴, 아딕이고 뇌터인
아핀 형식적 스킴)일 때, $U$를 \emph{아핀 형식적 열린집합}(각각
\emph{아딕 아핀 형식적 열린집합}, \emph{뇌터 아핀 형식적 열린집합})이라
한다.
\end{env}

\begin{definition}[10.4.2]
\label{I.10.4.2-ko}
모든 점이 아핀 형식적 열린 근방을 갖는 위상환 달린 공간
$\mathfrak{X}$를 형식적 준스킴이라 한다. 형식적 준스킴
$\mathfrak{X}$가 아딕(각각 국소 뇌터)이라는 것은 $\mathfrak{X}$의 모든
점이 아딕 아핀 형식적 열린 근방(각각 뇌터 아핀 형식적 열린 근방)을
갖는다는 뜻이다.
국소 뇌터인 $\mathfrak{X}$의 바탕공간이 준콤팩트일 때(따라서 그
바탕공간은 뇌터이다) 이를 뇌터라 한다.
\end{definition}

\begin{proposition}[10.4.3]
\label{I.10.4.3-ko}
$\mathfrak{X}$가 형식적 준스킴(각각 국소 뇌터 형식적 준스킴)이면,
아핀 형식적 열린집합들(각각 뇌터 아핀 형식적 열린집합들)은
$\mathfrak{X}$의 위상의 기저를 이룬다.
\end{proposition}

\begin{proof}
이는 (\hyperref[I.10.4.2-ko]{10.4.2})와
(\hyperref[I.10.1.4-ko]{10.1.4})에서 따른다. 여기서 $A$가 뇌터
아딕환이면 $A_{\{f\}}$도 모든 $f\in A$에 대하여 그러하다는 사실
(\hyperref[0.7.6.11-ko]{0, 7.6.11})을 사용한다.
\end{proof}

\begin{corollary}[10.4.4]
\label{I.10.4.4-ko}
$\mathfrak{X}$가 형식적 준스킴(각각 국소 뇌터 형식적 준스킴, 뇌터
형식적 준스킴)이면, $\mathfrak{X}$의 모든 열린집합 위에 유도된
위상환 달린 공간도 형식적 준스킴(각각 국소 뇌터 형식적 준스킴,
뇌터 형식적 준스킴)이다.
\end{corollary}

\begin{definition}[10.4.5]
\label{I.10.4.5-ko}
두 형식적 준스킴 $\mathfrak{X}$, $\mathfrak{Y}$가 주어졌다고 하자.
$\mathfrak{X}$에서 $\mathfrak{Y}$로 가는 위상환 달린 공간의 사상
$(\psi,\theta)$ 가운데 모든 $x\in\mathfrak{X}$에 대하여
$\theta_x^\sharp$가 국소 준동형
$\mathscr{O}_{\psi(x)}\to\mathscr{O}_x$인 것을 (형식적 준스킴의)
사상이라 한다.
\end{definition}

두 형식적 준스킴의 사상의 합성도 형식적 준스킴의 사상임은 자명하다.
따라서 형식적 준스킴들은 하나의 \emph{범주}를 이룬다.
$\operatorname{Hom}(\mathfrak{X},\mathfrak{Y})$로 형식적 준스킴
$\mathfrak{X}$에서 형식적 준스킴 $\mathfrak{Y}$로 가는 사상들의 집합을
나타낸다.

\oldpage[I]{186}
$U$가 $\mathfrak{X}$의 열린부분이면, $U$ 위에 $\mathfrak{X}$가 유도하는
형식적 준스킴에서 $\mathfrak{X}$로 가는 표준 포함 사상은 형식적
준스킴의 사상이다(더 나아가 위상환 달린 공간의 \emph{단사사상}이다
(\hyperref[0.4.1.1-ko]{0, 4.1.1})).

\begin{proposition}[10.4.6]
\label{I.10.4.6-ko}
$\mathfrak{X}$를 형식적 준스킴이라 하고,
$\mathfrak{S}=\operatorname{Spf}(A)$를 아핀 형식적 스킴이라 하자.
형식적 준스킴 $\mathfrak{X}$에서 형식적 준스킴 $\mathfrak{S}$로 가는
사상들과 환 $A$에서 위상환
$\Gamma(\mathfrak{X},\mathscr{O}_{\mathfrak{X}})$로 가는 연속
준동형들 사이에는 표준 전단사 대응이 존재한다.
\end{proposition}

증명은 (\hyperref[I.2.2.4-ko]{2.2.4})의 증명과 같다. 다만
“준동형”을 “연속 준동형”으로, “아핀 열린집합”을 “아핀 형식적
열린집합”으로 바꾸고, (\hyperref[I.10.2.2-ko]{10.2.2})를
(\hyperref[I.1.7.3-ko]{1.7.3}) 대신 사용한다. 자세한 내용은 독자에게
맡긴다.

\begin{env}[10.4.7]
\label{I.10.4.7-ko}
형식적 준스킴 $\mathfrak{S}$가 주어졌다고 하자. 형식적 준스킴
$\mathfrak{X}$와 사상 $\varphi:\mathfrak{X}\to\mathfrak{S}$의 자료가
주어졌을 때, 형식적 준스킴 $\mathfrak{X}$를
\emph{$\mathfrak{S}$ 위의 형식적 준스킴} 또는
\emph{$\mathfrak{S}$-형식적 준스킴}이라 하고,
$\varphi$를 $\mathfrak{S}$-형식적 준스킴 $\mathfrak{X}$의
\emph{구조 사상}이라 한다. $\mathfrak{S}=\operatorname{Spf}(A)$이고
$A$가 허용환이면, $\mathfrak{S}$-형식적 준스킴 $\mathfrak{X}$를
\emph{$A$-형식적 준스킴} 또는 \emph{$A$ 위의 형식적 준스킴}이라고도
한다. 모든 형식적 준스킴은 (이산 위상을 갖춘) $\mathbf{Z}$ 위의
형식적 준스킴으로 볼 수 있다.

$\mathfrak{X}$, $\mathfrak{Y}$를 두 $\mathfrak{S}$-형식적 준스킴이라
하자. 사상 $u:\mathfrak{X}\to\mathfrak{Y}$에 대하여 도식
\[
  \xymatrix{
    \mathfrak{X}\ar[rr]^u\ar[rd] & &
    \mathfrak{Y}\ar[ld]\\
    & \mathfrak{S}
  }
\]
(빗금 화살표들은 구조 사상이다)이 가환일 때, 이 사상을
\emph{$\mathfrak{S}$-사상}이라 한다. 이 정의에 따라 (고정된
$\mathfrak{S}$에 대한) $\mathfrak{S}$-형식적 준스킴들은 하나의
\emph{범주}를 이룬다.
$\operatorname{Hom}_{\mathfrak{S}}(\mathfrak{X},\mathfrak{Y})$로
$\mathfrak{S}$-형식적 준스킴 $\mathfrak{X}$에서
$\mathfrak{S}$-형식적 준스킴 $\mathfrak{Y}$로 가는
$\mathfrak{S}$-사상들의 집합을 나타낸다.
$\mathfrak{S}=\operatorname{Spf}(A)$이면 \emph{$A$-사상}이라는 말도
\emph{$\mathfrak{S}$-사상} 대신 사용한다.
\end{env}

\begin{env}[10.4.8]
\label{I.10.4.8-ko}
모든 아핀 스킴은 아핀 형식적 스킴으로 볼 수 있으므로
(\hyperref[I.10.1.2-ko]{10.1.2}), 모든 (일반적인) 준스킴은 형식적
준스킴으로 볼 수 있다. 또한 (\hyperref[I.10.4.5-ko]{10.4.5})에 따라
\emph{일반적인} 준스킴에 대해서는 \emph{형식적} 준스킴의 사상(각각
$S$-사상)이 §~\hyperref[section:I.2-ko]{2}에서 정의한 사상(각각
$S$-사상)과 일치한다.
\end{env}

\subsection{형식적 준스킴의 정의 아이디얼.}
\label{subsection:I.10.5-ko}

\begin{env}[10.5.1]
\label{I.10.5.1-ko}
형식적 준스킴 $\mathfrak{X}$가 주어졌다고 하자.
$\mathscr{O}_{\mathfrak{X}}$-아이디얼 $\mathscr{J}$가 다음 조건을 만족할
때 이를 $\mathfrak{X}$의 \emph{정의 아이디얼층}(또는
\emph{정의 아이디얼})이라 한다. 모든 $x\in\mathfrak{X}$에 대하여 아핀
형식적 열린 근방 $U$가 존재하여, $\mathscr{J}|U$는 $U$ 위에
$\mathfrak{X}$가 유도하는 아핀 형식적 스킴의 정의 아이디얼이다
(\hyperref[I.10.3.3-ko]{10.3.3}). 이때
(\hyperref[I.10.3.1.1-ko]{10.3.1.1})과
(\hyperref[I.10.4.3-ko]{10.4.3})에 따라 모든 열린집합
$V\subset\mathfrak{X}$에 대하여 $\mathscr{J}|V$는 $V$ 위에
$\mathfrak{X}$가 유도하는 형식적 준스킴의 정의 아이디얼이다.

$\mathfrak{X}$의 정의 아이디얼들로 이루어진 족
$(\mathscr{J}_\lambda)$에 대하여, $\mathfrak{X}$의 아핀 형식적
열린집합들로 이루어진 덮개 $(U_\alpha)$가 존재하여 모든 $\alpha$에
대해 $\mathscr{J}_\lambda|U_\alpha$들로 이루어진 족이 $U_\alpha$ 위에
$\mathfrak{X}$가 유도하는 아핀 형식적 스킴의 정의 아이디얼들의
기본계를 이룰 때(\hyperref[I.10.3.6-ko]{10.3.6}), 이 족을
\emph{정의 아이디얼들의}
\oldpage[I]{187}
\emph{기본계}라고 한다. (\hyperref[I.10.3.7-ko]{10.3.7})의 마지막
주에 따라, $\mathfrak{X}$가 아핀 형식적 스킴이면 이 정의는
(\hyperref[I.10.3.7-ko]{10.3.7})에서 주어진 정의와 일치한다. $V$를
$\mathfrak{X}$의 임의의 열린집합이라 하면, $\mathscr{J}_\lambda|V$들은
$V$ 위에 유도된 형식적 준스킴의 정의 아이디얼들의 기본계를 이룬다. 이는
(\hyperref[I.10.3.1.1-ko]{10.3.1.1})에서 따른다. $\mathfrak{X}$가
\emph{국소 뇌터} 형식적 준스킴이고 $\mathscr{J}$가 $\mathfrak{X}$의
정의 아이디얼이면, (\hyperref[I.10.3.6-ko]{10.3.6})에 따라 거듭제곱들
$\mathscr{J}^n$은 $\mathfrak{X}$의 정의 아이디얼들의 기본계를 이룬다.
\end{env}

\begin{env}[10.5.2]
\label{I.10.5.2-ko}
$\mathfrak{X}$를 형식적 준스킴이라 하고, $\mathscr{J}$를
$\mathfrak{X}$의 정의 아이디얼이라 하자. 그러면 환 달린 공간
$(\mathfrak{X},\mathscr{O}_{\mathfrak{X}}/\mathscr{J})$는
\emph{준스킴}(일반적인 의미에서)이다. $\mathfrak{X}$가 아핀 형식적
스킴(각각 국소 뇌터 형식적 준스킴, 뇌터 형식적 준스킴)이면 이 준스킴은
아핀(각각 국소 뇌터, 뇌터)이다. 실제로 곧바로 아핀인 경우로 환원되며,
그 경우에는 (\hyperref[I.10.3.2-ko]{10.3.2})에서 이미 이 명제를
증명했다. 또한
$\theta:\mathscr{O}_{\mathfrak{X}}\to
\mathscr{O}_{\mathfrak{X}}/\mathscr{J}$가 표준 준동형이면,
$u=(1_{\mathfrak{X}},\theta)$는 형식적 준스킴의 \emph{사상}
(\emph{표준 사상}이라 한다)
$(\mathfrak{X},\mathscr{O}_{\mathfrak{X}}/\mathscr{J})\to
(\mathfrak{X},\mathscr{O}_{\mathfrak{X}})$이다. 이 사실도 곧바로 아핀인
경우로 환원되며, 그 경우에는
(\hyperref[I.10.3.2-ko]{10.3.2})에서 이미 확인했다.
\end{env}

\begin{proposition}[10.5.3]
\label{I.10.5.3-ko}
$\mathfrak{X}$를 형식적 준스킴이라 하고, $(\mathscr{J}_\lambda)$를
$\mathfrak{X}$의 정의 아이디얼들의 기본계라 하자. 그러면 위상환의 층
$\mathscr{O}_{\mathfrak{X}}$는 의사이산 위상환의 층들
(\hyperref[0.3.8.1-ko]{0, 3.8.1})
$\mathscr{O}_{\mathfrak{X}}/\mathscr{J}_\lambda$의 사영극한이다.
\end{proposition}

\begin{proof}
$\mathfrak{X}$의 위상에는 준콤팩트 아핀 형식적 열린집합들로 이루어진
기저가 있으므로(\hyperref[I.10.4.3-ko]{10.4.3}), 아핀인 경우로
환원된다. 이 경우 명제는 (\hyperref[I.10.3.5-ko]{10.3.5}),
(\hyperref[I.10.3.2-ko]{10.3.2})와
(\hyperref[I.10.1.1-ko]{10.1.1})의 정의에서 따른다.
\end{proof}

모든 형식적 준스킴이 정의 아이디얼을 갖는지는 확실하지 않다. 그러나
다음은 성립한다.

\begin{proposition}[10.5.4]
\label{I.10.5.4-ko}
$\mathfrak{X}$를 국소 뇌터 형식적 준스킴이라 하자. 그러면
정의 아이디얼 $\mathscr{T}$가 $\mathfrak{X}$의 가장 큰 정의
아이디얼로서 존재한다. 이는 정의 아이디얼 $\mathscr{J}$ 가운데 준스킴
$(\mathfrak{X},\mathscr{O}_{\mathfrak{X}}/\mathscr{J})$가 축소인
유일한 것이다. $\mathscr{J}$가
$\mathfrak{X}$의 정의 아이디얼이면, $\mathscr{T}$는 표준 준동형
$\mathscr{O}_{\mathfrak{X}}\to
\mathscr{O}_{\mathfrak{X}}/\mathscr{J}$에 의한
$\mathscr{O}_{\mathfrak{X}}/\mathscr{J}$의 멱영근기의 역상이다.
\end{proposition}

\begin{proof}
먼저 $\mathfrak{X}=\operatorname{Spf}(A)$이고 $A$가 뇌터 아딕환이라고
하자. $\mathscr{T}$의 존재성과 성질들은
(\hyperref[I.10.3.5-ko]{10.3.5})\footnote{\textbf{번역 주.} 인쇄 원문은
이곳에서 (10.3.4)를 참조하지만, 출판된 EGA II 정오표는 이를
(10.3.5)로 고치도록 지시한다. 그 공식 교정을 적용했다.}와
(\hyperref[I.5.1.1-ko]{5.1.1})에서 바로 따른다. 이때 $A$의 가장 큰
정의 아이디얼의 존재성과 성질들
(\hyperref[0.7.1.6-ko]{0, 7.1.6}과
\hyperref[0.7.1.7-ko]{7.1.7})을 이용한다.

일반적인 경우에 $\mathscr{T}$의 존재성과 성질들을 증명하려면,
$U\supset V$를 $\mathfrak{X}$의 두 뇌터 아핀 형식적 열린집합이라 할 때
$U$의 가장 큰 정의 아이디얼 $\mathscr{T}_U$가 $V$의 가장 큰 정의
아이디얼 $\mathscr{T}_V$를 유도함을 보이면 충분하다. 그런데
$(V,(\mathscr{O}_{\mathfrak{X}}|V)/(\mathscr{T}_U|V))$는 축소이므로,
이는 앞에서 보인 결과에서 따른다.
\end{proof}

$\mathfrak{X}_{\mathrm{red}}$는 축소 준스킴(일반적인 의미에서)
$(\mathfrak{X},\mathscr{O}_{\mathfrak{X}}/\mathscr{T})$를 나타낸다.

\begin{corollary}[10.5.5]
\label{I.10.5.5-ko}
$\mathfrak{X}$를 국소 뇌터 형식적 준스킴이라 하고, $\mathscr{T}$를
$\mathfrak{X}$의 가장 큰 정의 아이디얼이라 하자. $V$가
$\mathfrak{X}$의 임의의 열린집합이면, $\mathscr{T}|V$는
$\mathfrak{X}$가 $V$ 위에 유도하는 형식적 준스킴의 가장 큰 정의
아이디얼이다.
\end{corollary}

\oldpage[I]{188}

\begin{proposition}[10.5.6]
\label{I.10.5.6-ko}
$\mathfrak{X}$, $\mathfrak{Y}$를 두 형식적 준스킴이라 하고,
$\mathscr{J}$와 $\mathscr{K}$를 각각 $\mathfrak{X}$와
$\mathfrak{Y}$의 정의 아이디얼이라 하며,
$f:\mathfrak{X}\to\mathfrak{Y}$를 형식적 준스킴의 사상이라 하자.
\begin{enumerate}
  \item[{\upshape(i)}]
    $f^*(\mathscr{K})\mathscr{O}_{\mathfrak{X}}\subset\mathscr{J}$이면,
    일반적인 의미에서 준스킴의 사상인
    $f':(\mathfrak{X},\mathscr{O}_{\mathfrak{X}}/\mathscr{J})\to
    (\mathfrak{Y},\mathscr{O}_{\mathfrak{Y}}/\mathscr{K})$가 유일하게
    존재하고 다음 도식을 가환이게 한다.
    \[
      \label{I.10.5.6.1-ko}
      \xymatrix{
        (\mathfrak{X},\mathscr{O}_{\mathfrak{X}})\ar[r]^f &
        (\mathfrak{Y},\mathscr{O}_{\mathfrak{Y}})\\
        (\mathfrak{X},\mathscr{O}_{\mathfrak{X}}/\mathscr{J})
          \ar[r]_{f'}\ar[u] &
        (\mathfrak{Y},\mathscr{O}_{\mathfrak{Y}}/\mathscr{K})\ar[u]
      }
      \tag{10.5.6.1}
    \]
    여기서 세로 화살표들은 표준 사상이다.
  \item[{\upshape(ii)}]
    $\mathfrak{X}=\operatorname{Spf}(A)$,
    $\mathfrak{Y}=\operatorname{Spf}(B)$가 아핀 형식적 스킴이고,
    $\mathscr{J}=\mathfrak{J}^{\Delta}$,
    $\mathscr{K}=\mathfrak{K}^{\Delta}$라 하자. 여기서
    $\mathfrak{J}$와 $\mathfrak{K}$는 각각 $A$와 $B$의 정의
    아이디얼이다. 또한
    $f=({}^a\varphi,\widetilde{\varphi})$이고
    $\varphi:B\to A$가 연속 준동형이라 하자. 이때
    $f^*(\mathscr{K})\mathscr{O}_{\mathfrak{X}}\subset\mathscr{J}$일
    필요충분조건은
    $\varphi(\mathfrak{K})\subset\mathfrak{J}$인 것이다. 이 경우
    $f'$은 사상 $({}^a\varphi',\widetilde{\varphi'})$이며, 여기서
    $\varphi':B/\mathfrak{K}\to A/\mathfrak{J}$은 $\varphi$에서
    몫으로 내려가 얻은 준동형이다.
\end{enumerate}
\end{proposition}

\begin{proof}
\medskip\noindent
\begin{enumerate}
  \item[{\upshape(i)}] $f=(\psi,\theta)$라 하면, 가정에 따라 준동형
    $\theta^\sharp:\psi^*(\mathscr{O}_{\mathfrak{Y}})\to
    \mathscr{O}_{\mathfrak{X}}$에 의한
    $\psi^*(\mathscr{O}_{\mathfrak{Y}})$의 아이디얼층
    $\psi^*(\mathscr{K})$의 상은 $\mathscr{J}$에 포함된다
    (\hyperref[0.4.3.5-ko]{0, 4.3.5}). 따라서 몫으로 내려가면
    $\theta^\sharp$에서 환의 층의 준동형
    \[
      \omega:\psi^*(\mathscr{O}_{\mathfrak{Y}}/\mathscr{K})
      =\psi^*(\mathscr{O}_{\mathfrak{Y}})/\psi^*(\mathscr{K})
      \to\mathscr{O}_{\mathfrak{X}}/\mathscr{J}~;
    \]
    을 얻는다. 또한 모든 $x\in\mathfrak{X}$에 대하여
    $\theta_x^\sharp$가 \emph{국소} 준동형이므로 $\omega_x$도 그러하다.
    따라서 환 달린 공간의 사상 $(\psi,\omega^b)$는
    (\hyperref[I.2.2.1-ko]{2.2.1}) 문제의 요구를 만족하는 환 달린 공간의
    유일한 사상 $f'$이다.
  \item[{\upshape(ii)}] 아핀 형식적 스킴의 사상들과 환의 연속 준동형들
    사이의 표준적이고 함자적인 대응
    (\hyperref[I.10.2.2-ko]{10.2.2})에 따르면, 지금의 경우 관계
    $f^*(\mathscr{K})\mathscr{O}_{\mathfrak{X}}\subset\mathscr{J}$로부터
    $f'=({}^a\varphi',\widetilde{\varphi'})$를 얻는다. 여기서
    $\varphi':B/\mathfrak{K}\to A/\mathfrak{J}$은 다음 도식을 가환이게
    하는 유일한 준동형이다.
    \[
      \label{I.10.5.6.2-ko}
      \xymatrix{
        B\ar[r]^\varphi\ar[d] & A\ar[d]\\
        B/\mathfrak{K}\ar[r]_{\varphi'} & A/\mathfrak{J}
      }
      \tag{10.5.6.2}
    \]
    따라서 $\varphi'$이 존재하면
    $\varphi(\mathfrak{K})\subset\mathfrak{J}$이다. 역으로 이 조건이
    성립한다고 하자. $\varphi'$을 도식
    (\hyperref[I.10.5.6.2-ko]{10.5.6.2})을 가환이게 하는 유일한
    준동형이라 하고 $f'=({}^a\varphi',\widetilde{\varphi'})$로 놓으면,
    도식 (\hyperref[I.10.5.6.1-ko]{10.5.6.1})이 가환임은 명백하다. 이제
    각각 $f$와 $f'$에 대응하는 준동형
    ${}^a\varphi^*(\mathscr{O}_{\mathfrak{Y}})\to
    \mathscr{O}_{\mathfrak{X}}$와
    ${}^a{\varphi'}^*(\mathscr{O}_{\mathfrak{Y}}/\mathscr{K})\to
    \mathscr{O}_{\mathfrak{X}}/\mathscr{J}$을 고려하면, 이 가환성으로부터
    $f^*(\mathscr{K})\mathscr{O}_{\mathfrak{X}}\subset\mathscr{J}$가
    따름을 알 수 있다.
\end{enumerate}
\end{proof}

위에서 정의한 대응 $f\mapsto f'$이 \emph{함자적}임은 명백하다.

\subsection{준스킴들의 귀납극한으로서의 형식적 준스킴.}
\label{subsection:I.10.6-ko}

\begin{env}[10.6.1]
\label{I.10.6.1-ko}
$\mathfrak{X}$를 형식적 준스킴이라 하고,
$(\mathscr{J}_\lambda)$를 $\mathfrak{X}$의 정의 아이디얼들의 기본계라
하자. 각 $\lambda$에 대하여 $f_\lambda$를 표준 사상
$(\mathfrak{X},\mathscr{O}_{\mathfrak{X}}/\mathscr{J}_\lambda)
\to\mathfrak{X}$라 하자(\hyperref[I.10.5.2-ko]{10.5.2}).
$\mathscr{J}_\mu\subset\mathscr{J}_\lambda$이면 표준 준동형
$\mathscr{O}_{\mathfrak{X}}/\mathscr{J}_\mu\to
\mathscr{O}_{\mathfrak{X}}/\mathscr{J}_\lambda$는 보통 의미의 준스킴
사이의 표준 사상

\oldpage[I]{189}

$f_{\mu\lambda}:(\mathfrak{X},
\mathscr{O}_{\mathfrak{X}}/\mathscr{J}_\lambda)\to
(\mathfrak{X},\mathscr{O}_{\mathfrak{X}}/\mathscr{J}_\mu)$를 정의하며,
$f_\lambda=f_\mu\circ f_{\mu\lambda}$가 성립한다. 따라서 준스킴들
$X_\lambda=(\mathfrak{X},
\mathscr{O}_{\mathfrak{X}}/\mathscr{J}_\lambda)$와 사상들
$f_{\mu\lambda}$는 (\hyperref[I.10.4.8-ko]{10.4.8})에 따라 형식적
준스킴의 범주에서 하나의 \emph{귀납계}를 이룬다.
\end{env}

\begin{proposition}[10.6.2]
\label{I.10.6.2-ko}
(\hyperref[I.10.6.1-ko]{10.6.1})의 표기 아래에서 형식적 준스킴
$\mathfrak{X}$와 사상들 $f_\lambda$는 형식적 준스킴의 범주에서 계
$(X_\lambda,f_{\mu\lambda})$의 귀납극한(T, I, 1.8)을 이룬다.
\end{proposition}

\begin{proof}
$\mathfrak{Y}$를 형식적 준스킴이라 하고, 각 지표 $\lambda$에 대하여
사상
\[
  g_\lambda=(\psi_\lambda,\theta_\lambda):X_\lambda\to\mathfrak{Y}
\]
가 주어졌으며, $\mathscr{J}_\mu\subset\mathscr{J}_\lambda$이면
$g_\lambda=g_\mu\circ f_{\mu\lambda}$가 성립한다고 하자. 이 조건과
$X_\lambda$의 정의로부터 우선 모든 $\psi_\lambda$가 바탕공간 사이의
동일한 연속사상 $\psi:\mathfrak{X}\to\mathfrak{Y}$와 일치함을 알 수 있다.
또한 준동형
$\theta_\lambda^\sharp:\psi^*(\mathscr{O}_{\mathfrak{Y}})\to
\mathscr{O}_{X_\lambda}=
\mathscr{O}_{\mathfrak{X}}/\mathscr{J}_\lambda$들은 환의 층의 준동형들의
\emph{사영계}를 이룬다. 따라서 사영극한을 취하면 준동형
\[
  \omega:\psi^*(\mathscr{O}_{\mathfrak{Y}})\to
  \varprojlim\mathscr{O}_{\mathfrak{X}}/\mathscr{J}_\lambda
  =\mathscr{O}_{\mathfrak{X}}
\]
을 얻는다. \emph{환 달린 공간}의 사상 $g=(\psi,\omega^b)$가 도식들
\[
  \label{I.10.6.2.1-ko}
  \xymatrix{
    X_\lambda\ar[rr]^{g_\lambda}\ar[rd]_{f_\lambda} &&
    \mathfrak{Y}\\
    & \mathfrak{X}\ar[ru]_g
  }
  \tag{10.6.2.1}
\]
을 가환이게 하는 \emph{유일한} 사상임은 명백하다. 그러므로 이제 $g$가
\emph{형식적 준스킴}의 사상임을 보이면 된다. 이 문제는
$\mathfrak{X}$와 $\mathfrak{Y}$ 위에서 국소적이므로
$\mathfrak{X}=\operatorname{Spf}(A)$,
$\mathfrak{Y}=\operatorname{Spf}(B)$라 해도 된다. 여기서 $A$와 $B$는
허용환이고 $\mathscr{J}_\lambda=\mathfrak{J}_\lambda^{\Delta}$이며,
$(\mathfrak{J}_\lambda)$는 $A$의 정의 아이디얼들의 기본계이다
(\hyperref[I.10.3.5-ko]{10.3.5}). 이제
$A=\varprojlim A/\mathfrak{J}_\lambda$이므로, 도식
(\hyperref[I.10.6.2.1-ko]{10.6.2.1})을 가환이게 하는 아핀 형식적
스킴의 사상 $g$의 존재는 아핀 형식적 스킴의 사상과 환의 연속
준동형 사이의 전단사 대응 (\hyperref[I.10.2.2-ko]{10.2.2}) 및
사영극한의 정의에서 따른다. 한편 환 달린 공간의 사상으로서 $g$가
유일하므로, 이 사상은 증명 앞부분에서 같은 기호로 나타낸 사상과
일치한다.
\end{proof}

다음 명제는 몇 가지 추가 조건 아래, 보통 의미의 준스킴들로 주어진
귀납계의 귀납극한이 형식적 준스킴의 범주에서 존재함을 보인다.

\begin{proposition}[10.6.3]
\label{I.10.6.3-ko}
$\mathfrak{X}$를 위상공간이라 하고,
$(\mathscr{O}_i,u_{ji})$를 $\mathfrak{X}$ 위의 환의 층들의 사영계라
하자. 그 지표집합은 $\mathbf{N}$이다. $\mathscr{J}_i$를
$u_{0i}:\mathscr{O}_i\to\mathscr{O}_0$의 핵이라 하자. 다음을 가정한다.
\begin{enumerate}
  \item[a)] 환 달린 공간 $(\mathfrak{X},\mathscr{O}_i)$는 준스킴
    $X_i$이다.
  \item[b)] 모든 $x\in\mathfrak{X}$와 모든 $i$에 대하여,
    $\mathfrak{X}$ 안에서 $x$의 열린 근방 $U_i$가 존재하여 제한
    $\mathscr{J}_i|U_i$는 멱영이다.\footnote{\textbf{번역 주.} 인쇄
    원문은 여기서 로만체 $X$를 쓰지만, 출판된 EGA II 정오표는 이를
    프락투어체 $\mathfrak{X}$로 고치도록 지시한다. 그 공식 교정을
    적용했다.}
  \item[c)] 준동형들 $u_{ji}$는 전사이다.
\end{enumerate}

\oldpage[I]{190}

$\mathscr{O}_{\mathfrak{X}}$를 의사이산 환의 층들 $\mathscr{O}_i$의
사영극한인 위상환의 층이라 하고,
$u_i:\mathscr{O}_{\mathfrak{X}}\to\mathscr{O}_i$를 표준 준동형이라
하자. 그러면 위상환 달린 공간
$(\mathfrak{X},\mathscr{O}_{\mathfrak{X}})$는 형식적 준스킴이다.
준동형들 $u_i$는 전사이고, 그 핵들 $\mathscr{J}^{(i)}$는
$\mathfrak{X}$의 정의 아이디얼들의 기본계를 이루며,
$\mathscr{J}^{(0)}$은 아이디얼층들 $\mathscr{J}_i$의 사영극한이다.
\end{proposition}

\begin{proof}
먼저 각 줄기에서 $u_{ji}$는 전사 준동형이며
\emph{따라서 특히} 국소 준동형임에 유의하자. 그러므로
$v_{ij}=(1_{\mathfrak{X}},u_{ji})$는 준스킴의 사상
$X_j\to X_i$ ($i\geq j$)이다
(\hyperref[I.2.2.1-ko]{2.2.1}). 우선 각 $X_i$가
환 $A_i$의 아핀 스킴이라고 가정하자.
\emph{환 준동형}
$\varphi_{ji}:A_i\to A_j$가 존재하여
$u_{ji}=\widetilde{\varphi_{ji}}$를 만족한다
(\hyperref[I.1.7.3-ko]{1.7.3}). 따라서
(\hyperref[I.1.6.3-ko]{1.6.3}), 층 $\mathscr{O}_j$는
$\mathscr{O}_i$-가군으로서 $X_i$ 위에서 준연접이다(그 스칼라
작용은 $u_{ji}$로 정의된다). 이 층은 $A_j$를 $A_i$-가군으로 보아
얻는 가군에 대응하며, 이때 그 가군 구조는 $\varphi_{ji}$로 정한다.
모든 $f\in A_i$에 대하여
$f'=\varphi_{ji}(f)$라 하자. 가정에 따라 열린집합 $D(f)$와
$D(f')$는 $\mathfrak{X}$에서 동일하며,
$\Gamma(D(f),\mathscr{O}_i)=(A_i)_f$에서
$\Gamma(D(f),\mathscr{O}_j)=(A_j)_{f'}$로 가는 준동형 가운데
$u_{ji}$에 대응하는 것은 바로
$(\varphi_{ji})_f$이다
(\hyperref[I.1.6.1-ko]{1.6.1}). 한편 $A_j$를 $A_i$-가군으로 보면,
$(A_j)_{f'}$는 $(A_i)_f$-가군 $(A_j)_f$이다. 따라서
$u_{ji}=\widetilde{\varphi_{ji}}$가 여전히 성립하는데, 이번에는
$\varphi_{ji}$를 $A_i$-가군 준동형으로 본다. 이제 $u_{ji}$가
전사이므로 $\varphi_{ji}$도 전사임을 알 수 있다
(\hyperref[I.1.3.9-ko]{1.3.9}). $\mathfrak{J}_{ji}$가
$\varphi_{ji}$의 핵이면, $u_{ji}$의 핵은 준연접
$\mathscr{O}_i$-가군이고 $\widetilde{\mathfrak{J}_{ji}}$와 같다. 특히
$\mathscr{J}_i=\widetilde{\mathfrak{J}_i}$이고, 여기서
$\mathfrak{J}_i$는 $\varphi_{0i}:A_i\to A_0$의 핵이다. 가정 b)로부터
$\mathscr{J}_i$가 \emph{멱영}임이 따른다. 실제로
$\mathfrak{X}$가 준콤팩트이므로, $\mathfrak{X}$를 유한 개의
열린집합 $U_k$로 덮을 수 있고, 각각
$(\mathscr{J}_i|U_k)^{n_k}=0$을 만족한다.
$n$을 $n_k$들 가운데 가장 큰 수로 택하면
$\mathscr{J}_i^n=0$이다. 따라서 $\mathfrak{J}_i$는 멱영이다
(\hyperref[I.1.3.13-ko]{1.3.13}). 이제 환
$A=\varprojlim A_i$는 허용환이다
(\hyperref[0.7.2.2-ko]{0, 7.2.2}). 또한 표준 준동형
$\varphi_i:A\to A_i$는 전사이다. 그 핵 $\mathfrak{J}^{(i)}$는
$\mathfrak{J}_{ik}$들($k\geq i$)의 사영극한이며,
$\mathfrak{J}^{(i)}$들은 $A$에서 0의 근방들의 기본계를 이룬다.
이 경우 (\hyperref[I.10.6.3-ko]{10.6.3})의 주장들은
(\hyperref[I.10.1.1-ko]{10.1.1})과
(\hyperref[I.10.3.2-ko]{10.3.2})에서 따른다. 실제로
$(\mathfrak{X},\mathscr{O}_{\mathfrak{X}})$는 다름 아닌
$\operatorname{Spf}(A)$이다.

여전히 이 특수한 경우에, $f=(f_i)$가 사영극한
$A=\varprojlim A_i$의 원소라고 하자. 모든 열린집합 $D(f_i)$는
$X_i$에서의 아핀 열린집합이며, $\mathfrak{D}(f)$라는
$\mathfrak{X}$의 열린집합과 동일시된다. 따라서 $X_i$가
$\mathfrak{D}(f)$에 유도하는 준스킴은 아핀 스킴
$\operatorname{Spec}((A_i)_{f_i})$와 동일시된다.

일반적인 경우에는 먼저, 임의의 준콤팩트 열린집합 $U$를
$\mathfrak{X}$에서 택하면 각 $\mathscr{J}_i|U$가 멱영임에
유의하자. 이는 위의
논법으로 알 수 있다. 이제 모든 $x\in\mathfrak{X}$에 대하여,
어떤 열린집합 $U$가 존재하여 $x$를 포함하고 $\mathfrak{X}$에서 열린 근방이며,
\emph{모든 $X_i$에 대하여 아핀 열린집합}임을 보이자. 실제로
$U$를 $X_0$의 아핀 열린집합으로 택하고
$\mathscr{O}_{X_0}=\mathscr{O}_{X_i}/\mathscr{J}_i$임에 유의하자.
앞에서 본 바와 같이 $\mathscr{J}_i|U$가 멱영이므로,
(\hyperref[I.5.1.9-ko]{5.1.9})에 따라 $U$는 각 $X_i$에 대하여도
아핀 열린집합이다. 이제 앞의 조건을 만족하는 모든 $U$에 대하여,
위에서 다룬 아핀 경우에 의해 $(U,\mathscr{O}_{\mathfrak{X}}|U)$는
형식적 준스킴이고, $\mathscr{J}^{(i)}|U$들은 정의 아이디얼들의
기본계를 이루며, $\mathscr{J}^{(0)}|U$는
$\mathscr{J}_i|U$들의 사영극한이다. 이로써 결론을 얻는다.
\end{proof}

\begin{corollary}[10.6.4]
\label{I.10.6.4-ko}
$i\geq j$일 때 $u_{ji}$의 핵이
$\mathscr{J}_i^{j+1}$이고 $\mathscr{J}_1/\mathscr{J}_1^2$가

\oldpage[I]{191}

$\mathscr{O}_0=\mathscr{O}_1/\mathscr{J}_1$-가군으로서 유한
생성이라고 가정하자. 그러면 $\mathfrak{X}$는 아딕 형식적 준스킴이다.
또한 $\mathscr{J}^{(n)}$을
$\mathscr{O}_{\mathfrak{X}}\to\mathscr{O}_n$의 핵이라 하면
$\mathscr{J}^{(n)}=\mathscr{J}^{n+1}$이고
$\mathscr{J}/\mathscr{J}^2$는 $\mathscr{J}_1$과 동형이다. 더욱이
$X_0$가 국소 뇌터(각각 뇌터)이면 $\mathfrak{X}$도 국소 뇌터(각각
뇌터)이다.
\end{corollary}

\begin{proof}
$\mathfrak{X}$와 $X_0$의 바탕공간은 같으므로 문제는 국소적이고, 모든
$X_i$가 아핀이라 해도 된다. 관계
$\mathscr{J}_{ij}=\widetilde{\mathfrak{J}_{ji}}$
((\hyperref[I.10.6.3-ko]{10.6.3})의 표기)를 고려하면
(\hyperref[0.7.2.7-ko]{0, 7.2.7}과
\hyperref[0.7.2.8-ko]{7.2.8})의 대응하는 주장들로 곧바로 환원된다.
여기서 $\mathfrak{J}_1/\mathfrak{J}_1^2$는 유한 생성
$A_0$-가군임에 유의한다
(\hyperref[I.1.3.9-ko]{1.3.9}).
\end{proof}

특히 \emph{모든 국소 뇌터 형식적 준스킴
$\mathfrak{X}$}는
(\hyperref[I.10.6.3-ko]{10.6.3})과
(\hyperref[I.10.6.4-ko]{10.6.4})의 조건을 만족하는 (보통 의미의)
국소 뇌터 준스킴들의 수열 $(X_n)$의 귀납극한이다. 실제로
$\mathscr{J}$를 $\mathfrak{X}$의 정의 아이디얼로 택하고
(\hyperref[I.10.5.4-ko]{10.5.4})
$X_n=(\mathfrak{X},
\mathscr{O}_{\mathfrak{X}}/\mathscr{J}^{n+1})$
로 두면 충분하다
((\hyperref[I.10.5.1-ko]{10.5.1})과
(\hyperref[I.10.6.2-ko]{10.6.2})).

\begin{corollary}[10.6.5]
\label{I.10.6.5-ko}
$A$를 허용환이라 하자. 아핀 형식적 스킴
$\mathfrak{X}=\operatorname{Spf}(A)$가 뇌터이기 위한 필요충분조건은
$A$가 아딕이고 뇌터인 것이다.
\end{corollary}

\begin{proof}
그 조건이 충분함은 명백하다. 역으로 $\mathfrak{X}$가 뇌터라고
가정하고, $\mathfrak{J}$를 $A$의 정의 아이디얼로,
$\mathscr{J}=\mathfrak{J}^\Delta$를 이에 대응하는
$\mathfrak{X}$의 정의 아이디얼로 택하자. 그러면 (보통 의미의) 준스킴
$X_n=(\mathfrak{X},
\mathscr{O}_{\mathfrak{X}}/\mathscr{J}^{n+1})$은 아핀이고 뇌터이다.
따라서 환 $A_n=A/\mathfrak{J}^{n+1}$도 뇌터이고
(\hyperref[I.6.1.3-ko]{6.1.3}), 이로부터
$\mathfrak{J}/\mathfrak{J}^2$가 유한 생성
$A/\mathfrak{J}$-가군임을 알 수 있다. $\mathscr{J}^n$들이
$\mathfrak{X}$의 정의 아이디얼들의 기본계를 이루므로
(\hyperref[I.10.5.1-ko]{10.5.1}),
$\mathscr{O}_{\mathfrak{X}}
=\varprojlim(\mathscr{O}_{\mathfrak{X}}/\mathscr{J}^n)$이다
(\hyperref[I.10.5.3-ko]{10.5.3}). 따라서
(\hyperref[I.10.1.3-ko]{10.1.3})에 의하여 $A$는
$\varprojlim A/\mathfrak{J}^n$과 위상적으로 동형이고, 그러므로
아딕이고 뇌터이다
(\hyperref[0.7.2.8-ko]{0, 7.2.8}).
\end{proof}

\begin{remark}[10.6.6]
\label{I.10.6.6-ko}
(\hyperref[I.10.6.3-ko]{10.6.3})의 표기에서
$\mathscr{F}_i$를 $\mathscr{O}_i$-가군이라 하자. 또한
$i\geq j$에 대하여 $v_{ij}$-사상
$\theta_{ji}:\mathscr{F}_i\to\mathscr{F}_j$가 주어졌고,
$\theta_{kj}\circ\theta_{ji}=\theta_{ki}$가
$k\leq j\leq i$에 대하여 성립한다고 하자. $v_{ij}$의 바탕 연속사상은
항등사상이므로, $\theta_{ji}$는 공간 $\mathfrak{X}$ 위의 아벨군의
층의 준동형이다. 또한 $\mathscr{F}$가 아벨군의 층들의 사영계
$(\mathscr{F}_i)$의 사영극한이면, $\theta_{ji}$들이
$v_{ij}$-사상이라는 사실로부터 사영극한을 취하여 $\mathscr{F}$에
$\mathscr{O}_{\mathfrak{X}}$-가군 구조를 정의할 수 있다. 이 구조를
갖춘 $\mathscr{F}$를, $\theta_{ji}$들에 관한
$\mathscr{O}_i$-가군들의 계 $(\mathscr{F}_i)$의
\emph{사영극한}이라 하겠다. 특별히
$v_{ij}^*(\mathscr{F}_i)=\mathscr{F}_j$이고 $\theta_{ji}$가
\emph{항등사상}인 경우에는, 줄여서 $\mathscr{F}$를 계
$(\mathscr{F}_i)$의 사영극한이라 하겠다. 이때
$v_{ij}^*(\mathscr{F}_i)=\mathscr{F}_j$는 $j\leq i$에 대하여
성립하며, $\theta_{ji}$들은 언급하지 않는다.
\end{remark}

\begin{env}[10.6.7]
\label{I.10.6.7-ko}
$\mathfrak{X}$와 $\mathfrak{Y}$를 두 형식적 준스킴이라 하고,
$\mathscr{J}$와 $\mathscr{K}$를 각각 $\mathfrak{X}$와
$\mathfrak{Y}$의 정의 아이디얼이라 하자. 또한
$f:\mathfrak{X}\to\mathfrak{Y}$를
$f^*(\mathscr{K})\mathscr{O}_{\mathfrak{X}}\subset\mathscr{J}$를
만족하는 형식적 준스킴의 사상이라 하자. 그러면 모든 정수 $n>0$에 대하여
$f^*(\mathscr{K}^n)\mathscr{O}_{\mathfrak{X}}
=(f^*(\mathscr{K})\mathscr{O}_{\mathfrak{X}})^n
\subset\mathscr{J}^n$이다. 따라서
(\hyperref[I.10.5.6-ko]{10.5.6})에 의하여 $f$로부터 (보통 의미의)
준스킴의 사상 $f_n:X_n\to Y_n$을 유도할 수 있다. 여기서
$X_n=(\mathfrak{X},
\mathscr{O}_{\mathfrak{X}}/\mathscr{J}^{n+1})$,
$Y_n=(\mathfrak{Y},
\mathscr{O}_{\mathfrak{Y}}/\mathscr{K}^{n+1})$로 둔다. 정의로부터
곧바로 다음 도식들이
\[
  \label{I.10.6.7.1-ko}
  \xymatrix{
    X_m\ar[r]^{f_m}\ar[d] &
    Y_m\ar[d]\\
    X_n\ar[r]_{f_n} &
    Y_n
  }
  \tag{10.6.7.1}
\]

\oldpage[I]{192}

$m\leq n$일 때 가환함을 알 수 있다. 다시 말해, $(f_n)$은 사상들의
\emph{귀납계}이다.
\end{env}

\begin{env}[10.6.8]
\label{I.10.6.8-ko}
역으로, $(X_n)$과 $(Y_n)$을 각각
(\hyperref[I.10.6.3-ko]{10.6.3})의 조건 b)와 c)를 만족하는 (보통
의미의) 준스킴들의 귀납계라 하고, $\mathfrak{X}$와
$\mathfrak{Y}$를 각각 그 귀납극한이라 하자. 귀납극한의 정의에 따라,
임의의 수열 $(f_n)$이 $X_n\to Y_n$인 사상들로 이루어진 귀납계이면 그
\emph{귀납극한} $f:\mathfrak{X}\to\mathfrak{Y}$가 존재한다. 이 사상은
다음 도식들을 가환이게 하는 유일한 형식적 준스킴의 사상이다.
\[
  \xymatrix{
    X_n\ar[r]^{f_n}\ar[d] &
    Y_n\ar[d]\\
    \mathfrak{X}\ar[r]_f &
    \mathfrak{Y}
  }
\]
\end{env}

\begin{proposition}[10.6.9]
\label{I.10.6.9-ko}
$\mathfrak{X}$와 $\mathfrak{Y}$를 두 국소 뇌터 형식적 준스킴이라
하고, $\mathscr{J}$와 $\mathscr{K}$를 각각 $\mathfrak{X}$와
$\mathfrak{Y}$의 정의 아이디얼이라 하자.
(\hyperref[I.10.6.7-ko]{10.6.7})에서 정의한 대응
$f\mapsto(f_n)$은
사상 $f:\mathfrak{X}\to\mathfrak{Y}$ 가운데
$f^*(\mathscr{K})\mathscr{O}_{\mathfrak{X}}\subset\mathscr{J}$를
만족하는 것들의 집합과, 도식들
(\hyperref[I.10.6.7.1-ko]{10.6.7.1})을 가환이게 하는 사상들의 수열
$(f_n)$의 집합 사이의 전단사 대응이다.
\end{proposition}

\begin{proof}
$f$가 그러한 사상 수열의 귀납극한이면
$f^*(\mathscr{K})\mathscr{O}_{\mathfrak{X}}\subset\mathscr{J}$임을
보여야 한다. 문제는 $\mathfrak{X}$와 $\mathfrak{Y}$ 위에서
국소적이므로, $\mathfrak{X}=\operatorname{Spf}(A)$,
$\mathfrak{Y}=\operatorname{Spf}(B)$인 아핀 경우로 한정해도 된다.
이때 $A$와 $B$는 아딕이고 뇌터이다. 또한
$\mathscr{J}=\mathfrak{J}^\Delta$,
$\mathscr{K}=\mathfrak{K}^\Delta$이다. 여기서 $\mathfrak{J}$와
$\mathfrak{K}$는 각각 $A$와 $B$의 정의 아이디얼이다. 그러면
$X_n=\operatorname{Spec}(A_n)$,
$Y_n=\operatorname{Spec}(B_n)$이고, 여기서
$A_n=A/\mathfrak{J}^{n+1}$,
$B_n=B/\mathfrak{K}^{n+1}$이며, 이는
(\hyperref[I.10.3.6-ko]{10.3.6})과
(\hyperref[I.10.3.2-ko]{10.3.2})에 따른다. 또한
$f_n=({}^a\varphi_n,\widetilde{\varphi_n})$이며, 환 준동형
$\varphi_n:B_n\to A_n$들은 사영계를 이룬다. 따라서
$f=({}^a\varphi,\widetilde{\varphi})$이고,
$\varphi=\varprojlim\varphi_n$이다. 이제 $m=0$일 때 도식
(\hyperref[I.10.6.7.1-ko]{10.6.7.1})의 가환성으로부터 조건
$\varphi_n(\mathfrak{K}/\mathfrak{K}^{n+1})
\subset\mathfrak{J}/\mathfrak{J}^{n+1}$이 모든 $n$에 대하여
성립한다. 따라서 사영극한을 취하면
$\varphi(\mathfrak{K})\subset
\mathfrak{J}$이고, 이로부터
$f^*(\mathscr{K})\mathscr{O}_{\mathfrak{X}}\subset\mathscr{J}$가
따른다
(\hyperref[I.10.5.6-ko]{10.5.6, (ii)}).
\end{proof}

\begin{corollary}[10.6.10]
\label{I.10.6.10-ko}
$\mathfrak{X}$와 $\mathfrak{Y}$를 두 국소 뇌터 형식적 준스킴이라 하고,
$\mathscr{T}$를 $\mathfrak{X}$의 가장 큰 정의 아이디얼이라 하자
(\hyperref[I.10.5.4-ko]{10.5.4}).
\begin{enumerate}
  \item[{\upshape(i)}] $\mathfrak{Y}$의 임의의 정의 아이디얼
    $\mathscr{K}$와 임의의 사상
    $f:\mathfrak{X}\to\mathfrak{Y}$에 대하여,
    $f^*(\mathscr{K})\mathscr{O}_{\mathfrak{X}}\subset\mathscr{T}$이다.
  \item[{\upshape(ii)}]
    $\operatorname{Hom}(\mathfrak{X},\mathfrak{Y})$과 사상 수열
    $(f_n)$ 가운데 도식들
    (\hyperref[I.10.6.7.1-ko]{10.6.7.1})을 가환이게 하는 것들의 집합
    사이에는 표준 전단사 대응이 있다. 여기서
    $X_n=(\mathfrak{X},
    \mathscr{O}_{\mathfrak{X}}/\mathscr{T}^{n+1})$,
    $Y_n=(\mathfrak{Y},
    \mathscr{O}_{\mathfrak{Y}}/\mathscr{K}^{n+1})$이다.
\end{enumerate}
\end{corollary}

\begin{proof}
(ii)는 (i)와 (\hyperref[I.10.6.9-ko]{10.6.9})에서 곧바로 따른다.
(i)를 보이기 위해서는 $\mathfrak{X}=\operatorname{Spf}(A)$,
$\mathfrak{Y}=\operatorname{Spf}(B)$이고 $A$와 $B$가 뇌터이며,
$\mathscr{T}=\mathfrak{T}^\Delta$,
$\mathscr{K}=\mathfrak{K}^\Delta$인 경우로 한정해도 된다. 여기서
$\mathfrak{T}$는 $A$의 가장 큰 정의 아이디얼이고
$\mathfrak{K}$는 $B$의 정의 아이디얼이다.
$f=({}^a\varphi,\widetilde{\varphi})$라 하자. 여기서
$\varphi:B\to A$는 연속 준동형이다. $\mathfrak{K}$의 원소들은
위상적으로 멱영이고
(\hyperref[0.7.1.4-ko]{0, 7.1.4, (ii)}), 따라서
$\varphi(\mathfrak{K})$의 원소들도 위상적으로 멱영이다. 그러므로
$\mathfrak{T}$가 $A$의 위상적으로 멱영인 원소들의 집합이므로
(\hyperref[0.7.1.6-ko]{0, 7.1.6}),
$\varphi(\mathfrak{K})\subset\mathfrak{T}$이다. 따라서
(\hyperref[I.10.5.6-ko]{10.5.6, (ii)})에 의해 결론이 따른다.
\end{proof}

\begin{corollary}[10.6.11]
\label{I.10.6.11-ko}
$\mathfrak{S}$, $\mathfrak{X}$, $\mathfrak{Y}$를 세 국소 뇌터
형식적 준스킴이라 하고,
$f:\mathfrak{X}\to\mathfrak{S}$,
$g:\mathfrak{Y}\to\mathfrak{S}$를 $\mathfrak{X}$와
$\mathfrak{Y}$를 $\mathfrak{S}$-형식적 준스킴으로 만드는 사상이라 하자.
$\mathscr{J}$ (각각 $\mathscr{K}$, $\mathscr{L}$)를
$\mathfrak{S}$ (각각 $\mathfrak{X}$, $\mathfrak{Y}$)의 정의 아이디얼이라 하고,
$f^*(\mathscr{J})\mathscr{O}_{\mathfrak{X}}\subset\mathscr{K}$,
$g^*(\mathscr{J})\mathscr{O}_{\mathfrak{Y}}=\mathscr{L}$이라 가정하자.
다음과 같이 둔다.
$S_n=(\mathfrak{S},
\mathscr{O}_{\mathfrak{S}}/\mathscr{J}^{n+1})$,
$X_n=(\mathfrak{X},
\mathscr{O}_{\mathfrak{X}}/\mathscr{K}^{n+1})$,
$Y_n=(\mathfrak{Y},
\mathscr{O}_{\mathfrak{Y}}/\mathscr{L}^{n+1})$.
그러면

\oldpage[I]{193}

$\operatorname{Hom}_{\mathfrak{S}}(\mathfrak{X},\mathfrak{Y})$과 수열
$(u_n)$ 가운데 각 항이 $S_n$-사상
$u_n:X_n\to Y_n$이고 도식들
(\hyperref[I.10.6.7.1-ko]{10.6.7.1})을 가환이게 하는 것들의 집합
사이에는 표준 전단사 대응이 있다.
\end{corollary}

\begin{proof}
임의의 $\mathfrak{S}$-사상
$u:\mathfrak{X}\to\mathfrak{Y}$에 대하여, 정의에 의해
$f=g\circ u$이므로
\[
  u^*(\mathscr{L})\mathscr{O}_{\mathfrak{X}}
  =u^*(g^*(\mathscr{J})\mathscr{O}_{\mathfrak{Y}})
    \mathscr{O}_{\mathfrak{X}}
  =f^*(\mathscr{J})\mathscr{O}_{\mathfrak{X}}
  \subset\mathscr{K}
\]
이다. 따라서 이 따름정리는
(\hyperref[I.10.6.9-ko]{10.6.9})에서 따른다.
\end{proof}

$m\leq n$일 때, 사상 $f_n:X_n\to Y_n$ 하나가 주어지면
도식 (\hyperref[I.10.6.7.1-ko]{10.6.7.1})을 가환이게 하는 사상
$f_m:X_m\to Y_m$가 유일하게 정해진다. 이는 아핀 경우로 환원하면
곧바로 알 수 있다. 이로써 사상
$\varphi_{mn}:\operatorname{Hom}_{S_n}(X_n,Y_n)
\to\operatorname{Hom}_{S_m}(X_m,Y_m)$가 정의되고,
$\operatorname{Hom}_{S_n}(X_n,Y_n)$들은 $\varphi_{mn}$들에 관하여
\emph{집합들의 사영계}를 이룬다.
(\hyperref[I.10.6.11-ko]{10.6.11})은 또한 다음과 같은 표준 전단사 사상이
존재한다고 진술할 수도 있다.
\[
  \operatorname{Hom}_{\mathfrak{S}}(\mathfrak{X},\mathfrak{Y})
  \xrightarrow{\sim}
  \varprojlim_n\operatorname{Hom}_{S_n}(X_n,Y_n)
\]

\subsection{형식적 준스킴의 곱.}
\label{subsection:I.10.7-ko}

\begin{env}[10.7.1]
\label{I.10.7.1-ko}
$\mathfrak{S}$를 형식적 준스킴이라 하자.
$\mathfrak{S}$-형식적 준스킴들은 하나의 범주를 이루므로,
$\mathfrak{S}$-형식적 준스킴의 \emph{곱}이라는 개념을 정의할 수 있다.
\end{env}

\begin{proposition}[10.7.2]
\label{I.10.7.2-ko}
$\mathfrak{X}=\operatorname{Spf}(B)$,
$\mathfrak{Y}=\operatorname{Spf}(C)$를 아핀 형식적 스킴
$\mathfrak{S}=\operatorname{Spf}(A)$ 위의 두 아핀 형식적 스킴이라 하자.
$\mathfrak{Z}=\operatorname{Spf}(B\widehat{\otimes}_A C)$라 하고,
$p_1$, $p_2$를 $B$와 $C$에서
$B\widehat{\otimes}_A C$로 가는 표준 연속 $A$-준동형
$\rho$, $\sigma$에 각각 대응하는
(\hyperref[I.10.2.2-ko]{10.2.2}) $\mathfrak{S}$-사상이라 하자.
그러면 $(\mathfrak{Z},p_1,p_2)$는 $\mathfrak{S}$ 위의 아핀 형식적 스킴
$\mathfrak{X}$와 $\mathfrak{Y}$의 곱이다.
\end{proposition}

\begin{proof}
(\hyperref[I.10.4.6-ko]{10.4.6})에 따라, 위상 $A$-대수인 허용환 $D$와
임의의 연속 $A$-준동형
$\varphi:B\widehat{\otimes}_A C\to D$에 대하여, 쌍
$(\varphi\circ\rho,\varphi\circ\sigma)$를 대응시키는 사상
\[
  \operatorname{Hom}_A(B\widehat{\otimes}_A C,D)
  \xrightarrow{\sim}
  \operatorname{Hom}_A(B,D)\times\operatorname{Hom}_A(C,D)
\]
이 전단사임을 확인하면 충분하다. 이는 다름 아닌 완비 텐서곱의 보편 성질
(\hyperref[0.7.7.6-ko]{0, 7.7.6})이다.
\end{proof}

\begin{proposition}[10.7.3]
\label{I.10.7.3-ko}
두 $\mathfrak{S}$-형식적 준스킴 $\mathfrak{X}$,
$\mathfrak{Y}$가 주어지면, 곱
$\mathfrak{X}\times_{\mathfrak{S}}\mathfrak{Y}$가 존재한다.
\end{proposition}

\begin{proof}
증명은 (\hyperref[I.3.2.6-ko]{3.2.6})의 증명에서 아핀 스킴(각각
아핀 열린집합)을 아핀 형식적 스킴(각각 아핀 형식적 열린집합)으로
바꾸고, 명제 (\hyperref[I.3.2.2-ko]{3.2.2})를
(\hyperref[I.10.7.2-ko]{10.7.2})로 바꾸면 그대로 얻어진다.
\end{proof}

준스킴의 곱에 관한 모든 형식적 성질
(\hyperref[I.3.2.7-ko]{3.2.7}과
\hyperref[I.3.2.8-ko]{3.2.8},
\hyperref[I.3.3.1-ko]{3.3.1}부터
\hyperref[I.3.3.12-ko]{3.3.12}까지)은
형식적 준스킴의 곱에도 아무 변경 없이 성립한다.

\begin{env}[10.7.4]
\label{I.10.7.4-ko}
$\mathfrak{S}$, $\mathfrak{X}$, $\mathfrak{Y}$를 세 형식적 준스킴이라
하고, $f:\mathfrak{X}\to\mathfrak{S}$,
$g:\mathfrak{Y}\to\mathfrak{S}$를 두 사상이라 하자.
$\mathfrak{S}$, $\mathfrak{X}$, $\mathfrak{Y}$ 안에 각각 정의
아이디얼들의 기본계 $(\mathscr{J}_\lambda)$,
$(\mathscr{K}_\lambda)$, $(\mathscr{L}_\lambda)$가 존재한다고 가정하자.
이 세 기본계의 지표집합은 모두 $I$이며,
$f^*(\mathscr{J}_\lambda)\mathscr{O}_{\mathfrak{X}}
\subset\mathscr{K}_\lambda$ 및
$g^*(\mathscr{J}_\lambda)\mathscr{O}_{\mathfrak{Y}}
\subset\mathscr{L}_\lambda$가 모든 $\lambda$에 대하여 성립한다고 하자.
다음과 같이 놓자.
$S_\lambda=(\mathfrak{S},
\mathscr{O}_{\mathfrak{S}}/\mathscr{J}_\lambda)$,
$X_\lambda=(\mathfrak{X},
\mathscr{O}_{\mathfrak{X}}/\mathscr{K}_\lambda)$,
$Y_\lambda=(\mathfrak{Y},
\mathscr{O}_{\mathfrak{Y}}/\mathscr{L}_\lambda)$.
$\mathscr{J}_\mu\subset\mathscr{J}_\lambda$,
$\mathscr{K}_\mu\subset\mathscr{K}_\lambda$,
$\mathscr{L}_\mu\subset\mathscr{L}_\lambda$인 경우,
$S_\lambda$, $X_\lambda$, $Y_\lambda$는 각각
$S_\mu$, $X_\mu$, $Y_\mu$의 닫힌 부분준스킴이고, 각각은 대응하는
준스킴과 같은

\oldpage[I]{194}

바탕공간을 갖는다(\hyperref[I.10.6.1-ko]{10.6.1}).
$S_\lambda\to S_\mu$가 준스킴의 단사사상이므로, 먼저 두 곱
$X_\lambda\times_{S_\lambda}Y_\lambda$와
$X_\lambda\times_{S_\mu}Y_\lambda$가 동일함을 알 수 있고
(\hyperref[I.3.2.4-ko]{3.2.4}), 이어서
$X_\lambda\times_{S_\mu}Y_\lambda$는
$X_\mu\times_{S_\mu}Y_\mu$의 닫힌 부분준스킴으로 동일시되며, 이
부분준스킴은 후자와 같은 바탕공간을 갖는다
(\hyperref[I.4.3.1-ko]{4.3.1}). 이제 곱
$\mathfrak{X}\times_{\mathfrak{S}}\mathfrak{Y}$는 (보통 의미의) 준스킴
$X_\lambda\times_{S_\lambda}Y_\lambda$들의 \emph{귀납극한}이다. 실제로
(\hyperref[I.10.6.2-ko]{10.6.2})에서와 같이
$\mathfrak{S}$, $\mathfrak{X}$, $\mathfrak{Y}$가 아핀 형식적 스킴인
경우로 환원할 수 있다. (\hyperref[I.10.5.6-ko]{10.5.6, (ii)})와
$\mathfrak{S}$, $\mathfrak{X}$, $\mathfrak{Y}$의 정의 아이디얼들의
기본계에 관한 가정을 고려하면, 이 주장은 두 대수의 완비 텐서곱의
정의(\hyperref[0.7.7.1-ko]{0, 7.7.1})에서 곧바로 따른다.

또한 $\mathfrak{Z}$를 $\mathfrak{S}$-형식적 준스킴이라 하고,
$(\mathscr{M}_\lambda)$를 $\mathfrak{Z}$의 정의 아이디얼들의 기본계라
하되 그 지표집합은 $I$라 하자. 또
$u:\mathfrak{Z}\to\mathfrak{X}$,
$v:\mathfrak{Z}\to\mathfrak{Y}$를 두 $\mathfrak{S}$-사상이라 하고,
$u^*(\mathscr{K}_\lambda)\mathscr{O}_{\mathfrak{Z}}
\subset\mathscr{M}_\lambda$ 및
$v^*(\mathscr{L}_\lambda)\mathscr{O}_{\mathfrak{Z}}
\subset\mathscr{M}_\lambda$라고 하자.
$Z_\lambda=(\mathfrak{Z},
\mathscr{O}_{\mathfrak{Z}}/\mathscr{M}_\lambda)$라 놓고,
$u_\lambda:Z_\lambda\to X_\lambda$와
$v_\lambda:Z_\lambda\to Y_\lambda$를 $S_\lambda$-사상 가운데 각각
$u$와 $v$에 대응하는 것이라 하면
(\hyperref[I.10.5.6-ko]{10.5.6}), $(u,v)_{\mathfrak{S}}$가
$S_\lambda$-사상 $(u_\lambda,v_\lambda)_{S_\lambda}$들의
귀납극한임을 곧바로 확인할 수 있다.

이 번호의 고찰은 특히 $\mathfrak{S}$, $\mathfrak{X}$,
$\mathfrak{Y}$가 국소 뇌터인 경우, 각각 하나의 정의 아이디얼의
거듭제곱들로 이루어진 계를 정의 아이디얼들의 기본계로 택하면 적용된다
(\hyperref[I.10.5.1-ko]{10.5.1}). 그러나
$\mathfrak{X}\times_{\mathfrak{S}}\mathfrak{Y}$가 반드시 국소 뇌터인
것은 아니다(다만 (\agref{I.10.13.5-ko}{10.13.5})를 보라).
\end{env}

\subsection{닫힌 부분집합을 따른 준스킴의 형식적 완비화.}
\label{subsection:I.10.8-ko}

\begin{env}[10.8.1]
\label{I.10.8.1-ko}
$X$를 (보통 의미의) 국소 뇌터 준스킴이라 하고, $X'$을 $X$의
바탕공간의 닫힌 부분집합이라 하자. $\Phi$로 다음 조건을 만족하는
아이디얼층들의 집합을 나타내자. $\mathscr{J}$는
$\mathscr{O}_X$ 안의 연접 아이디얼층이고, $\mathscr{O}_X/\mathscr{J}$의
지지집합은 $X'$이다. 집합 $\Phi$는 비어 있지
않으며(\hyperref[I.5.2.1-ko]{5.2.1},
\hyperref[I.4.1.4-ko]{4.1.4}와
\hyperref[I.6.1.1-ko]{6.1.1}), 관계 $\supset$로 순서를 주기로 한다.
\end{env}

\begin{lemma}[10.8.2]
\label{I.10.8.2-ko}
순서집합 $\Phi$는 여과집합이다. $X$가 뇌터이면, 모든
$\mathscr{J}_0\in\Phi$에 대하여 거듭제곱 $\mathscr{J}_0^n$
($n>0$)으로 이루어진 집합은 $\Phi$에서 공종이다.
\end{lemma}

\begin{proof}
실제로 $\mathscr{J}_1$과 $\mathscr{J}_2$가 $\Phi$에 속하고
$\mathscr{J}=\mathscr{J}_1\cap\mathscr{J}_2$라 놓으면,
$\mathscr{O}_X$가 연접이므로 $\mathscr{J}$도 연접이고
(\hyperref[I.6.1.1-ko]{6.1.1}과
\hyperref[0.5.3.4-ko]{0, 5.3.4}), 모든 $x\in X$에 대하여
$\mathscr{J}_x=(\mathscr{J}_1)_x\cap(\mathscr{J}_2)_x$이다. 따라서
$x\notin X'$이면 $\mathscr{J}_x=\mathscr{O}_x$이고
$x\in X'$이면 $\mathscr{J}_x\neq\mathscr{O}_x$이므로
$\mathscr{J}\in\Phi$이다. 한편 $X$가 뇌터이고
$\mathscr{J}_0$와 $\mathscr{J}$가 $\Phi$에 속하면,
$\mathscr{J}_0^n(\mathscr{O}_X/\mathscr{J})=0$이 되는 정수
$n>0$이 존재한다(\hyperref[I.9.3.4-ko]{9.3.4}). 이는
$\mathscr{J}_0^n\subset\mathscr{J}$를 뜻한다.
\end{proof}

\begin{env}[10.8.3]
\label{I.10.8.3-ko}
이제 $\mathscr{F}$를 연접 $\mathscr{O}_X$-가군이라 하자. 모든
$\mathscr{J}\in\Phi$에 대하여
$\mathscr{F}\otimes_{\mathscr{O}_X}(\mathscr{O}_X/\mathscr{J})$는
연접 $\mathscr{O}_X$-가군이고
(\hyperref[I.9.1.1-ko]{9.1.1}), 그 지지집합은 $X'$에 포함된다. 우리는
이 층을 대개 $X'$으로 제한한 층과 동일시하겠다. $\mathscr{J}$가
$\Phi$를 달릴 때 이 층들은 아벨군의 층들의
\emph{사영계}를 이룬다.
\end{env}

\begin{definition}[10.8.4]
\label{I.10.8.4-ko}
국소 뇌터 준스킴 $X$의 닫힌 부분집합 $X'$과 연접
$\mathscr{O}_X$-가군 $\mathscr{F}$가 주어졌을 때, 층

\oldpage[I]{195}

$\varprojlim_{\Phi}
(\mathscr{F}\otimes_{\mathscr{O}_X}
(\mathscr{O}_X/\mathscr{J}))$의 $X'$으로의 제한을
\emph{$X'$을 따른 $\mathscr{F}$의 완비화}라 하고
$\mathscr{F}_{/X'}$로 나타내며, 혼동의 우려가 없으면
$\widehat{\mathscr{F}}$로도 나타낸다. 이 완비화의 $X'$ 위 단면들을
\emph{$X'$을 따른 $\mathscr{F}$의 형식적 단면}이라 한다.
\end{definition}

모든 열린집합 $U\subset X$에 대하여 다음이 곧바로 성립한다.
$(\mathscr{F}|U)_{/(U\cap X')}=(\mathscr{F}_{/X'})|(U\cap X')$.

사영극한을 취하면 $(\mathscr{O}_X)_{/X'}$가 환의 층이고
$\mathscr{F}_{/X'}$를 $(\mathscr{O}_X)_{/X'}$-가군으로 볼 수 있음은
명백하다. 또한 $X'$의 위상에는 준콤팩트 열린집합들로 이루어진 기저가
있으므로, $(\mathscr{O}_X)_{/X'}$ (각각 $\mathscr{F}_{/X'}$)를
\emph{의사이산} 환의 층들 (각각 의사이산 군의 층들)
$\mathscr{O}_X/\mathscr{J}$ (각각
$\mathscr{F}\otimes_{\mathscr{O}_X}
(\mathscr{O}_X/\mathscr{J})=\mathscr{F}/\mathscr{J}\mathscr{F}$)의
사영극한인 \emph{위상환의 층} (각각 \emph{위상군의 층})으로 볼 수
있다. 이때 사영극한을 취하면 $\mathscr{F}_{/X'}$는
\emph{위상 $(\mathscr{O}_X)_{/X'}$-가군}이 된다
(\hyperref[0.3.8.1-ko]{0, 3.8.1}과
\hyperref[0.3.8.2-ko]{3.8.2}). 모든 \emph{준콤팩트} 열린집합
$U\subset X$에 대하여
$\Gamma(U\cap X',(\mathscr{O}_X)_{/X'})$ (각각
$\Gamma(U\cap X',\mathscr{F}_{/X'})$)는 이산환들 (각각 이산군들)
$\Gamma(U,\mathscr{O}_X/\mathscr{J})$ (각각
$\Gamma(U,\mathscr{F}/\mathscr{J}\mathscr{F})$)의 사영극한임을
상기하자.

이제 $u:\mathscr{F}\to\mathscr{G}$가 $\mathscr{O}_X$-가군 사이의
준동형이면, 모든 $\mathscr{J}\in\Phi$에 대하여 표준적으로 준동형
$u_{\mathscr{J}}:
\mathscr{F}\otimes_{\mathscr{O}_X}(\mathscr{O}_X/\mathscr{J})
\to
\mathscr{G}\otimes_{\mathscr{O}_X}(\mathscr{O}_X/\mathscr{J})$를 얻으며,
이 준동형들은 사영계를 이룬다. 사영극한을 취하고 $X'$에 제한하면
연속 $(\mathscr{O}_X)_{/X'}$-준동형
$\mathscr{F}_{/X'}\to\mathscr{G}_{/X'}$를 얻는다. 이를
$u_{/X'}$ 또는 $\widehat u$로 나타내며,
\emph{$X'$을 따른 준동형 $u$의 완비화}라 한다. 또한
$v:\mathscr{G}\to\mathscr{H}$가 $\mathscr{O}_X$-가군 사이의
또 하나의 준동형이면
$(v\circ u)_{/X'}=(v_{/X'})\circ(u_{/X'})$임은 명백하다. 따라서
$\mathscr{F}_{/X'}$는 $\mathscr{F}$에 관한
\emph{가법적 공변 함자}이며, 연접 $\mathscr{O}_X$-가군의 범주에서
위상 $(\mathscr{O}_X)_{/X'}$-가군의 범주에 값을 갖는다.

\begin{proposition}[10.8.5]
\label{I.10.8.5-ko}
$(\mathscr{O}_X)_{/X'}$의 지지집합은 $X'$이다. 위상환 달린 공간
$(X',(\mathscr{O}_X)_{/X'})$는 국소 뇌터 형식적 준스킴이며, 모든
$\mathscr{J}\in\Phi$에 대하여 $\mathscr{J}_{/X'}$는 이 형식적
준스킴의 정의 아이디얼이다. $X=\operatorname{Spec}(A)$가 뇌터 환을
갖는 아핀 스킴이고, $\mathscr{J}=\widetilde{\mathfrak{J}}$이며,
$\mathfrak{J}$가 $A$의 아이디얼이고 $X'=V(\mathfrak{J})$이면,
$(X',(\mathscr{O}_X)_{/X'})$는 표준적으로
$\operatorname{Spf}(\widehat A)$와 동일시된다. 여기서 $\widehat A$는
$A$에 $\mathfrak{J}$-준아딕 위상을 주어 얻는 분리 완비화이다.
\end{proposition}

\begin{proof}
마지막 주장만 증명하면 충분함은 명백하다.
(\hyperref[0.7.3.3-ko]{0, 7.3.3})에 의해, $\mathfrak{J}$의
$\mathfrak{J}$-준아딕 위상에 관한 분리 완비화
$\widehat{\mathfrak{J}}$는 $\widehat A$의 아이디얼
$\mathfrak{J}\widehat A$와 동일시된다. 또한 $\widehat A$는
뇌터 $\widehat{\mathfrak{J}}$-아딕환이며
$\widehat A/\widehat{\mathfrak{J}}^n=A/\mathfrak{J}^n$임을 안다
(\hyperref[0.7.2.6-ko]{0, 7.2.6}). 이 마지막 관계에서
$\widehat A$의 열린 소 아이디얼들은
$\widehat{\mathfrak p}=\mathfrak p\widehat A$ 꼴의 아이디얼들임을 알 수
있다. 여기서 $\mathfrak p$는 $A$의 소 아이디얼이고
$\mathfrak{J}$를 포함하며, $\widehat{\mathfrak p}\cap A=\mathfrak p$이다. 따라서
바탕공간으로서 $\operatorname{Spf}(\widehat A)=X'$이다. 또한
$\mathscr{O}_X/\mathscr{J}^n=(A/\mathfrak{J}^n)^{\sim}$이므로, 명제는
정의들에서 곧바로 따른다.
\end{proof}

이렇게 정의한 형식적 준스킴을 \emph{$X'$을 따른 $X$의 완비화}라 하고
$X_{/X'}$로 나타내며, 혼동의 우려가 없으면 $\widehat X$로도 나타낸다.
$X'=X$로 잡을 때에는 $\mathscr{J}=0$을 택할 수 있으므로
$X_{/X}=X$이다.

$U$가 $X$의 열린집합 위에 유도된 부분준스킴이면,
$U_{/(U\cap X')}$는 $X_{/X'}$가 $X'$의 열린집합 $U\cap X'$ 위에
유도하는 형식적 부분준스킴과 표준적으로 동일시됨은 명백하다.

\begin{corollary}[10.8.6]
\label{I.10.8.6-ko}
(보통 의미의) 준스킴 $\widehat X_{\mathrm{red}}$는 $X$의 유일한 축소
부분준스킴으로서, 그 바탕공간은 $X'$이다
(\hyperref[I.5.2.1-ko]{5.2.1}). $\widehat X$가 뇌터이기 위한
필요충분조건은 $\widehat X_{\mathrm{red}}$가 뇌터인 것이며, 더 나아가
$X$가 뇌터이면 전자도 뇌터이다.
\end{corollary}

\oldpage[I]{196}

$\widehat X_{\mathrm{red}}$의 결정은 국소적이므로
(\hyperref[I.10.5.4-ko]{10.5.4}), 다시 $X$가 뇌터 환을 갖는 아핀 스킴인
경우로 한정할 수 있다. (\hyperref[I.10.8.5-ko]{10.8.5})의 표기를
쓰면, 아이디얼 $\mathfrak{T}$는 $\widehat A$의 위상적으로 멱영인
원소들로 이루어지며, 표준 사상
$\widehat A\to\widehat A/\widehat{\mathfrak{J}}=A/\mathfrak{J}$에 의한
$A/\mathfrak{J}$의 멱영근기의 역상이다
(\hyperref[0.7.1.3-ko]{0, 7.1.3}). 따라서
$\widehat A/\mathfrak{T}$는 $A/\mathfrak{J}$를 그 멱영근기로 나눈
몫과 동형이다. 그러므로 첫 번째 주장은
(\hyperref[I.10.5.4-ko]{10.5.4})와
(\hyperref[I.5.1.1-ko]{5.1.1})에서 따른다.
$\widehat X_{\mathrm{red}}$가 뇌터이면 그 바탕공간 $X'$도 뇌터이므로,
$X'_n=\operatorname{Spec}(\mathscr{O}_X/\mathscr{J}^n)$은 뇌터이고
(\hyperref[I.6.1.2-ko]{6.1.2}), $\widehat X$도 뇌터이다
(\hyperref[I.10.6.4-ko]{10.6.4}). 역은
(\hyperref[I.6.1.2-ko]{6.1.2})에 의해 명백하다.

\begin{env}[10.8.7]
\label{I.10.8.7-ko}
표준 준동형 $\mathscr{O}_X\to\mathscr{O}_X/\mathscr{J}$
($\mathscr{J}\in\Phi$)들은 사영계를 이루므로, 사영극한을 취하면 환의
층의 준동형
$\theta:\mathscr{O}_X\to\psi_*((\mathscr{O}_X)_{/X'})
=\varprojlim_\Phi(\mathscr{O}_X/\mathscr{J})$를 얻는다. 여기서 $\psi$는
바탕공간 사이의 표준 포함 사상 $X'\to X$이다. 환 달린 공간의 사상
\[
  (\psi,\theta):X_{/X'}\to X
\]
을 $i$ (또는 $i_X$)로 나타내고 표준 사상이라 하겠다.

텐서곱을 취하면, 모든 연접 $\mathscr{O}_X$-가군 $\mathscr{F}$에 대하여
표준 준동형 $\mathscr{O}_X\to\mathscr{O}_X/\mathscr{J}$들로부터
$\mathscr{O}_X$-가군 준동형
$\mathscr{F}\to
\mathscr{F}\otimes_{\mathscr{O}_X}(\mathscr{O}_X/\mathscr{J})$들을 얻으며,
이들도 사영계를 이룬다. 따라서 사영극한을 취하면 함자적인 표준
$\mathscr{O}_X$-가군 준동형
$\gamma:\mathscr{F}\to\psi_*(\mathscr{F}_{/X'})$를 얻는다.
\end{env}

\begin{proposition}[10.8.8]
\label{I.10.8.8-ko}
\medskip\noindent
\begin{enumerate}
  \item[\rm{(i)}] 함자 $\mathscr{F}_{/X'}$는
    ($\mathscr{F}$에 관하여) 완전하다.
  \item[\rm{(ii)}] 함자적인 준동형
    $\gamma^\sharp:i^*(\mathscr{F})\to\mathscr{F}_{/X'}$는
    $(\mathscr{O}_X)_{/X'}$-가군의 동형이다.
\end{enumerate}
\end{proposition}

\medskip\noindent
\begin{enumerate}
  \item[\rm{(i)}] 다음을 증명하면 충분하다.
    $0\to\mathscr{F}'\to\mathscr{F}\to\mathscr{F}''\to0$이 연접
    $\mathscr{O}_X$-가군의 완전열이고, $U$가 $X$의 아핀 열린집합이며
    그 환이 뇌터 환 $A$이면, 다음 열은 완전하다.
    \[
      0\to\Gamma(U\cap X',\mathscr{F}'_{/X'})
      \to\Gamma(U\cap X',\mathscr{F}_{/X'})
      \to\Gamma(U\cap X',\mathscr{F}''_{/X'})\to0
    \]
    이때 $\mathscr{F}|U=\widetilde M$,
    $\mathscr{F}'|U=\widetilde{M'}$, $\mathscr{F}''|U=\widetilde{M''}$이고,
    여기서 $M$, $M'$, $M''$는 세 유한 생성 $A$-가군이며 열
    $0\to M'\to M\to M''\to0$은 완전하다
    (\hyperref[I.1.5.1-ko]{1.5.1}과
    \hyperref[I.1.3.11-ko]{1.3.11})~; $\mathscr{J}\in\Phi$라 하고,
    $\mathfrak{J}$를 $A$의 아이디얼로서
    $\mathscr{J}|U=\widetilde{\mathfrak{J}}$인 것으로 택하자. 그러면
    \[
      \Gamma(U\cap X',
      \mathscr{F}\otimes_{\mathscr{O}_X}(\mathscr{O}_X/\mathscr{J}^n))
      =M\otimes_A(A/\mathfrak{J}^n)
    \]
    이다(\hyperref[I.1.3.12-ko]{1.3.12})~; 따라서 사영극한의 정의에 의해
    \[
      \Gamma(U\cap X',\mathscr{F}_{/X'})
      =\varprojlim_n(M\otimes_A(A/\mathfrak{J}^n))=\widehat M
    \]
    이고, 이는 $M$의 $\mathfrak{J}$-준아딕 위상에 관한 분리
    완비화이다. 마찬가지로
    \[
      \Gamma(U\cap X',\mathscr{F}'_{/X'})=\widehat{M'},\qquad
      \Gamma(U\cap X',\mathscr{F}''_{/X'})=\widehat{M''}~;
    \]
    이제 $A$가 뇌터이면 함자 $\widehat M$이 $M$에 관하여 유한 생성
    $A$-가군의 범주에서 완전하다는 사실에서 우리의 주장이 따른다
    (\hyperref[0.7.3.3-ko]{0, 7.3.3}).

\oldpage[I]{197}
  \item[\rm{(ii)}] 문제는 국소적이므로 완전열
    $\mathscr{O}_X^m\to\mathscr{O}_X^n\to\mathscr{F}\to0$이
    있다고 가정할 수 있다
    (\hyperref[0.5.3.2-ko]{0, 5.3.2})~; $\gamma^\sharp$가 함자적이고,
    함자 $i^*(\mathscr{F})$와 $\mathscr{F}_{/X'}$가 각각
    (\hyperref[0.4.3.1-ko]{0, 4.3.1})과 (i)에 의해 오른쪽 완전하므로
    다음 가환도를 얻는다.
    \[
      \xymatrix{
        i^*(\mathscr{O}_X^m)\ar[r]\ar[d]_{\gamma^\sharp} &
        i^*(\mathscr{O}_X^n)\ar[r]\ar[d]_{\gamma^\sharp} &
        i^*(\mathscr{F})\ar[r]\ar[d]_{\gamma^\sharp} &
        0\\
        (\mathscr{O}_X^m)_{/X'}\ar[r] &
        (\mathscr{O}_X^n)_{/X'}\ar[r] &
        \mathscr{F}_{/X'}\ar[r] &
        0
      }
      \tag{10.8.8.1}\label{I.10.8.8.1-ko}
    \]
    이 가환도의 두 행은 완전하다. 또한 두 함자
    $i^*(\mathscr{F})$와 $\mathscr{F}_{/X'}$는 유한 직합과 가환하므로
    (\hyperref[0.3.2.6-ko]{0, 3.2.6}과
    \hyperref[0.4.3.2-ko]{4.3.2}), 우리의 주장을
    $\mathscr{F}=\mathscr{O}_X$인 경우에 증명하는 것으로 환원된다.
    이때
    $i^*(\mathscr{O}_X)=(\mathscr{O}_X)_{/X'}=
    \mathscr{O}_{\widehat X}$
    이고(\hyperref[0.4.3.4-ko]{0, 4.3.4}), $\gamma^\sharp$는
    $\mathscr{O}_{\widehat X}$-가군 준동형이다. 따라서
    $\gamma^\sharp$가 $\mathscr{O}_{\widehat X}$의 단위원 단면을
    $X'$의 한 열린집합 위에서 자기 자신으로 보내는지 확인하면 충분하다.
    이는 명백하며, 따라서 이 경우 $\gamma^\sharp$는 항등준동형이다.
\end{enumerate}

\begin{corollary}[10.8.9]
\label{I.10.8.9-ko}
환 달린 공간의 사상 $i:X_{/X'}\to X$는 평탄하다.
\end{corollary}

이는 실제로 (\hyperref[0.6.7.3-ko]{0, 6.7.3})과
(\hyperref[I.10.8.8-ko]{10.8.8, (i)})에서 따른다.

\begin{corollary}[10.8.10]
\label{I.10.8.10-ko}
$\mathscr{F}$와 $\mathscr{G}$가 연접 $\mathscr{O}_X$-가군이면, 다음
표준 동형사상들이 존재하며 $\mathscr{F}$와 $\mathscr{G}$에 관하여 함자적이다.
\[
  (\mathscr{F}_{/X'})\otimes_{(\mathscr{O}_X)_{/X'}}
  (\mathscr{G}_{/X'})
  \xrightarrow{\sim}
  (\mathscr{F}\otimes_{\mathscr{O}_X}\mathscr{G})_{/X'}
  \tag{10.8.10.1}\label{I.10.8.10.1-ko}
\]
\[
  (\mathscr{H}\!om_{\mathscr{O}_X}(\mathscr{F},\mathscr{G}))_{/X'}
  \xrightarrow{\sim}
  \mathscr{H}\!om_{(\mathscr{O}_X)_{/X'}}
  (\mathscr{F}_{/X'},\mathscr{G}_{/X'})
  \tag{10.8.10.2}\label{I.10.8.10.2-ko}
\]
\end{corollary}

이는 $i^*(\mathscr{F})$와 $\mathscr{F}_{/X'}$의 표준 동일시에서 따른다~;
이때 첫 번째 동형의 존재는 모든 환 달린 공간의 사상에 대하여 성립하는
결과이고 (\hyperref[0.4.3.3.1-ko]{0, 4.3.3.1}), 두 번째 동형의 존재는
모든 평탄 사상에 대하여 성립하는 결과이므로
(\hyperref[0.6.7.6-ko]{0, 6.7.6}),
(\hyperref[I.10.8.9-ko]{10.8.9})에서 따른다.

\begin{proposition}[10.8.11]
\label{I.10.8.11-ko}
모든 연접 $\mathscr{O}_X$-가군 $\mathscr{F}$에 대하여, 표준 준동형
$\Gamma(X,\mathscr{F})\to\Gamma(X',\mathscr{F}_{/X'})$는
$\mathscr{F}\to\mathscr{F}_{/X'}$에서 유도되며, 그 핵은 $X'$의 어떤
근방에서 영인 단면들로 이루어진다.
\end{proposition}

그러한 단면의 표준 상이 영임은 $\mathscr{F}_{/X'}$의 정의에서 따른다.
역으로 $s\in\Gamma(X,\mathscr{F})$의 상이
$\Gamma(X',\mathscr{F}_{/X'})$에서 영이라고 하자. 모든 $x\in X'$가
$X$ 안에 하나의 근방을 가져 그곳에서 $s$가 영임을 보이면 충분하고, 따라서
$X=\operatorname{Spec}(A)$가 아핀 스킴이고 $A$가 뇌터 환이며,
$X'=V(\mathfrak{J})$이고 $\mathfrak{J}$가 $A$의 아이디얼이며,
$\mathscr{F}=\widetilde M$이고 $M$이 유한 생성 $A$-가군인 경우로
환원할 수 있다. 이때 $\Gamma(X',\mathscr{F}_{/X'})$는
분리 완비화 $\widehat M$이고, 이는 $M$을 $\mathfrak{J}$-준아딕 위상에
관하여 분리 완비화한 것이다. 또한
준동형 $\Gamma(X,\mathscr{F})\to\Gamma(X',\mathscr{F}_{/X'})$는
표준 준동형 $M\to\widehat M$이다. 이 준동형의 핵은
$z\in M$ 가운데 $1+\mathfrak{J}$의 한 원소에 의해 소멸되는 것들의 집합임이
알려져 있다 (\hyperref[0.7.3.7-ko]{0, 7.3.7}). 따라서
$(1+f)s=0$인 어떤 $f\in\mathfrak{J}$가 존재한다~; 모든 $x\in X'$에
대하여 이로부터 $(1_x+f_x)s_x=0$을 얻고,
$1_x+f_x$는 $\mathscr{O}_x$에서 가역인데(실제로
$\mathfrak{J}_x\mathscr{O}_x$는 $\mathscr{O}_x$의 극대 아이디얼에
포함된다). 따라서
$s_x=0$이고, 이는 명제를 증명한다.

\begin{corollary}[10.8.12]
\label{I.10.8.12-ko}
$\mathscr{F}_{/X'}$의 지지집합은
$\operatorname{Supp}(\mathscr{F})\cap X'$과 같다.
\end{corollary}

$\mathscr{F}_{/X'}$가 유한 생성 $(\mathscr{O}_X)_{/X'}$-가군임은
(\hyperref[I.10.8.8-ko]{10.8.8, (ii)})와
(\hyperref[0.5.2.4-ko]{0, 5.2.4})에서 명백하므로,

\oldpage[I]{198}

그 지지집합은 닫혀 있고
(\hyperref[0.5.2.2-ko]{0, 5.2.2}), 명백히
$\operatorname{Supp}(\mathscr{F})\cap X'$에 포함된다. 이 마지막
집합과 같음을 보이기 위해서는 등식
$\Gamma(X',\mathscr{F}_{/X'})=0$이
$\operatorname{Supp}(\mathscr{F})\cap X'=\varnothing$을 함의함을
증명하는 것으로 즉시 환원된다~; 그런데 이는
(\hyperref[I.10.8.11-ko]{10.8.11})과
(\hyperref[I.1.4.1-ko]{1.4.1})에서 따른다.

\begin{corollary}[10.8.13]
\label{I.10.8.13-ko}
$u:\mathscr{F}\to\mathscr{G}$를 연접 $\mathscr{O}_X$-가군 사이의 준동형이라
하자. $u_{/X'}:\mathscr{F}_{/X'}\to\mathscr{G}_{/X'}$가 영이기 위한
필요충분조건은 $u$가 $X'$의 어떤 근방에서 영 준동형인 것이다.
\end{corollary}

실제로 (\hyperref[I.10.8.8-ko]{10.8.8, (ii)})에 따라 $u_{/X'}$는
$i^*(u)$와 동일시된다. 따라서 $u$를 $X$ 위에서 층
$\mathscr{H}=\mathscr{H}\!om_{\mathscr{O}_X}(\mathscr{F},\mathscr{G})$의
단면으로 보면, $u_{/X'}$는
$i^*(\mathscr{H})=\mathscr{H}_{/X'}$의 $X'$ 위의 단면이며 앞서 본
단면에 표준적으로 대응한다
((\hyperref[I.10.8.10.2-ko]{10.8.10.2})와
(\hyperref[0.4.4.6-ko]{0, 4.4.6})). 그러므로 연접
$\mathscr{O}_X$-가군 $\mathscr{H}$에
(\hyperref[I.10.8.11-ko]{10.8.11})을 적용하면 충분하다.

\begin{corollary}[10.8.14]
\label{I.10.8.14-ko}
$u:\mathscr{F}\to\mathscr{G}$를 연접 $\mathscr{O}_X$-가군 사이의
준동형이라 하자. $u_{/X'}$가 단사(각각 전사)이기 위한 필요충분조건은
$u$가 $X'$의 어떤 근방에서 단사(각각 전사)인 것이다.
\end{corollary}

$\mathscr{P}$와 $\mathscr{N}$을 각각 $u$의 여핵과 핵이라 하자.
그러면 다음 완전열이 성립한다.
\[
  0\to\mathscr{N}\xrightarrow{v}\mathscr{F}\xrightarrow{u}
  \mathscr{G}\xrightarrow{w}\mathscr{P}\to0,
\]
따라서 (\hyperref[I.10.8.8-ko]{10.8.8, (i)})에 의해 다음 완전열이
성립한다.
\[
  0\to\mathscr{N}_{/X'}\xrightarrow{v_{/X'}}
  \mathscr{F}_{/X'}\xrightarrow{u_{/X'}}
  \mathscr{G}_{/X'}\xrightarrow{w_{/X'}}
  \mathscr{P}_{/X'}\to0.
\]
$u_{/X'}$가 단사(각각 전사)이면 $v_{/X'}=0$ (각각
$w_{/X'}=0$)이다. 따라서 (\hyperref[I.10.8.13-ko]{10.8.13})에 의해
$X'$의 어떤 근방에서는 $v=0$ (각각 $w=0$)이다.

\subsection{완비화로의 사상의 연장.}
\label{subsection:I.10.9-ko}

\begin{env}[10.9.1]
\label{I.10.9.1-ko}
$X$, $Y$를 (보통 의미의) 국소 뇌터 준스킴이라 하고, $f:X\to Y$를
사상이라 하며, $X'$ (각각 $Y'$)를 $X$ (각각 $Y$)의 바탕공간의 닫힌
부분집합이라 하자. $f(X')\subset Y'$라고 가정한다. $\mathscr{J}$
(각각 $\mathscr{K}$)를 $\mathscr{O}_X$ (각각 $\mathscr{O}_Y$)의
아이디얼층이라 하고, $\mathscr{O}_X/\mathscr{J}$ (각각
$\mathscr{O}_Y/\mathscr{K}$)의 지지집합은 $X'$ (각각 $Y'$)이며
$f^*(\mathscr{K})\mathscr{O}_X\subset\mathscr{J}$라고 하자. 이러한
아이디얼층은 항상 존재한다는 점을 지적해 두자. 예를 들어 $\mathscr{J}$로
$\mathscr{O}_X$의 아이디얼층 가운데 $X$의 부분준스킴을 정의하고
$X'$을 바탕공간으로 갖는 가장 큰 것을 택할 수 있고
(\hyperref[I.5.2.1-ko]{5.2.1}),
$f(X')\subset Y'$라는 가정에서
$f^*(\mathscr{K})\mathscr{O}_X\subset\mathscr{J}$가 따른다
(\hyperref[I.5.2.4-ko]{5.2.4}). 따라서 모든 정수 $n>0$에 대하여
$f^*(\mathscr{K}^n)\mathscr{O}_X\subset\mathscr{J}^n$
(\hyperref[0.4.3.5-ko]{0, 4.3.5})이고, 따라서
(\hyperref[I.4.4.6-ko]{4.4.6})에 의해 다음과 같이 놓으면
\[
  X'_n=(X',\mathscr{O}_X/\mathscr{J}^{n+1}),\qquad
  Y'_n=(Y',\mathscr{O}_Y/\mathscr{K}^{n+1}),
\]
$f$로부터 사상 $f_n:X'_n\to Y'_n$을 얻는다. 또한 $f_n$들은
귀납계를 이룬다. 그 귀납극한
(\hyperref[I.10.6.8-ko]{10.6.8})을
$\widehat f:X_{/X'}\to Y_{/Y'}$로 나타내고, (표현을 다소 남용하여)
$\widehat f$를 $f$의, $X$와 $Y$를 각각 $X'$와 $Y'$를 따라 완비화한
것들 사이로의 연장이라 부르겠다. 이 사상이 위 조건을
만족하는 아이디얼층 $\mathscr{J}$, $\mathscr{K}$의 선택에 의존하지
않음은 즉시 확인할 수 있다. 실제로 이를 확인하기 위해서는 $X$, $Y$가
각각 환 $A$, $B$의 아핀 뇌터 준스킴인 경우만 보면 충분하다. 이때
$\mathscr{J}=\widetilde{\mathfrak{J}}$,
$\mathscr{K}=\widetilde{\mathfrak{K}}$이고, $\mathfrak{J}$ (각각
$\mathfrak{K}$)는 $A$ (각각 $B$)의 아이디얼이며, $f$는
$\varphi:B\to A$라는 환 준동형에 대응하고
$\varphi(\mathfrak{K})\subset\mathfrak{J}$이다
(\hyperref[I.4.4.6-ko]{4.4.6} 및
\hyperref[I.1.7.4-ko]{1.7.4}). 이때
\emph{$\widehat f$는 (\hyperref[I.10.2.2-ko]{10.2.2})에 따라 연속
준동형 $\widehat\varphi:\widehat B\to\widehat A$에 대응하는
사상이다}. 여기서 $\widehat A$ (각각 $\widehat B$)는 $A$ (각각
$B$)에 $\mathfrak{J}$-준아딕 (각각 $\mathfrak{K}$-준아딕) 위상을
주어 얻는 분리 완비화이다
(\hyperref[I.10.6.8-ko]{10.6.8}). 또한
$\mathscr{J}$를 다른
\oldpage[I]{199}
아이디얼층 $\mathscr{J}'=\widetilde{\mathfrak{J}'}$로 바꾸어도,
$\mathscr{O}_X/\mathscr{J}'$의 지지집합이 여전히 $X'$이면,
$\mathfrak{J}$-준아딕 위상과 $\mathfrak{J}'$-준아딕 위상은 $A$ 위에서
같음을 알고 있다
(\hyperref[I.10.8.2-ko]{10.8.2}).
이 정의에 따르면, $X_{/X'}$와 $Y_{/Y'}$의 바탕공간 사이의
연속사상 $X'\to Y'$는 $\widehat f$에 대응하며, 다름 아닌 $X'$에 대한
$f$의 제한임을 지적하자.
\end{env}

\begin{env}[10.9.2]
\label{I.10.9.2-ko}
앞의 정의로부터 환 달린 공간의 사상들로 이루어진 도식
\[
  \label{I.10.9.2.1-ko}
  \xymatrix{
    \widehat X\ar[r]^{\widehat f}\ar[d]_{i_X} &
    \widehat Y\ar[d]^{i_Y}\\
    X\ar[r]_f &
    Y
  }
\]
이 가환이고, 세로 화살표들은 표준 사상
(\hyperref[I.10.8.7-ko]{10.8.7})임이 곧바로 따른다.
\end{env}

\begin{env}[10.9.3]
\label{I.10.9.3-ko}
$Z$를 세 번째 준스킴, $g:Y\to Z$를 사상, $Z'$을
$g(Y')\subset Z'$인 $Z$의 닫힌 부분집합이라 하자. $Y'$과 $Z'$을
따른 사상 $g$의 완비화를 $\widehat g$로 나타내면,
(\hyperref[I.10.9.1-ko]{10.9.1})로부터
$\widehat{(g\circ f)}=\widehat g\circ\widehat f$임이 곧바로 따른다.
\end{env}

\begin{proposition}[10.9.4]
\label{I.10.9.4-ko}
$X$, $Y$를 두 국소 뇌터 $S$-준스킴이라 하고, $Y$가 $S$ 위의
유한형이라고 하자. $f$, $g$를 $X$에서 $Y$로 가는 두 $S$-사상이라
하고, $f(X')\subset Y'$, $g(X')\subset Y'$라고 하자.
$\widehat f=\widehat g$일 필요충분조건은 $f$와 $g$가 $X'$의 어떤
근방에서 일치하는 것이다.
\end{proposition}

$Y$에 대한 유한성 가정 없이도 이 조건은 자명하게 충분하다.
필요성을 보기 위해, 먼저 가정 $\widehat f=\widehat g$에서 모든
$x\in X'$에 대하여 $f(x)=g(x)$가 따름을 지적하자. 한편 문제는
국소적이므로, $X$와 $Y$가 각각 $x$와 $y=f(x)=g(x)$의 아핀 열린
근방이고 그 좌표환들이 뇌터이며, $S$가 아핀이고
$\Gamma(Y,\mathscr{O}_Y)$가 $\Gamma(S,\mathscr{O}_S)$ 위의 유한 생성
대수라고 가정할 수 있다
(\hyperref[I.6.3.3-ko]{6.3.3}). 이때 $f$와 $g$는
$\Gamma(Y,\mathscr{O}_Y)$에서 $\Gamma(X,\mathscr{O}_X)$로 가는 두
$\Gamma(S,\mathscr{O}_S)$-대수 준동형 $\rho$, $\sigma$에 대응한다
(\hyperref[I.1.7.3-ko]{1.7.3}). 또한 가정에 따라 이 준동형들을
$\Gamma(Y,\mathscr{O}_Y)$의 분리 완비화에 연속적으로 연장한
준동형들은 서로 같다. 따라서 (\hyperref[I.10.8.11-ko]{10.8.11})에서
모든 단면 $s\in\Gamma(Y,\mathscr{O}_Y)$에 대하여 단면 $\rho(s)$와
$\sigma(s)$가 $X'$의 어떤 근방에서 일치한다는 결론을 얻는다(이
근방은 $s$에 의존한다)~; $\Gamma(Y,\mathscr{O}_Y)$가
$\Gamma(S,\mathscr{O}_S)$ 위의 유한 생성 대수이므로, $X'$의 한 근방
$V$가 존재하여 \emph{모든} 단면 $s\in\Gamma(Y,\mathscr{O}_Y)$에
대하여 $\rho(s)$와 $\sigma(s)$가 $V$에서 일치함이 곧바로 따른다.
$D(h)$가 $V$에 포함되는 $x$의 근방이 되도록
$h\in\Gamma(X,\mathscr{O}_X)$를 택하면, 앞의 결론과
(\hyperref[I.1.4.1-ko]{1.4.1, d)})에서 $f$와 $g$가 $D(h)$에서
일치함을 얻는다.\footnote{\textbf{번역 주.} 인쇄 원문은 여기서 형에
맞지 않는 $h\in\Gamma(Y,\mathscr{O}_X)$라고 쓴다. 그러나
$\mathscr{O}_X$는 $X$ 위의 층이고 이어지는 $D(h)$는 $X$ 안에서
$x$의 근방이므로, $h\in\Gamma(X,\mathscr{O}_X)$로 바로잡았다.}

\begin{proposition}[10.9.5]
\label{I.10.9.5-ko}
(\hyperref[I.10.9.1-ko]{10.9.1})의 가정 아래에서, 모든 연접
$\mathscr{O}_Y$-가군 $\mathscr{G}$에 대하여 다음과 같은
$(\mathscr{O}_X)_{/X'}$-가군의 함자적인 표준 동형이 존재한다.
\[
  (f^*(\mathscr{G}))_{/X'}\xrightarrow{\sim}
  \widehat f^*(\mathscr{G}_{/Y'})
\]
\end{proposition}

$(f^*(\mathscr{G}))_{/X'}$를 $i_X^*(f^*(\mathscr{G}))$와,
$\widehat f^*(\mathscr{G}_{/Y'})$를
$\widehat f^*(i_Y^*(\mathscr{G}))$와 표준적으로 동일시하면
(\hyperref[I.10.8.8-ko]{10.8.8}), 명제는
(\hyperref[I.10.9.2-ko]{10.9.2})의 도식이 가환이라는 사실에서 곧바로
따른다.

\begin{env}[10.9.6]
\label{I.10.9.6-ko}
이제 $\mathscr{F}$를 연접 $\mathscr{O}_X$-가군,
$\mathscr{G}$를 연접 $\mathscr{O}_Y$-가군이라 하자.
$u:\mathscr{G}\to\mathscr{F}$가 $\mathscr{G}$에서 $\mathscr{F}$로
가는 $f$-사상이면, 이에 대응하는 $\mathscr{O}_X$-준동형
$u^\sharp:f^*(\mathscr{G})\to\mathscr{F}$가 있다. 따라서 이 준동형을
완비화하면 연속 $(\mathscr{O}_X)_{/X'}$-준동형
$(u^\sharp)_{/X'}:(f^*(\mathscr{G}))_{/X'}\to\mathscr{F}_{/X'}$를
얻고, (\hyperref[I.10.9.5-ko]{10.9.5})에 의해
$\widehat f$-사상 $v:\mathscr{G}_{/Y'}\to\mathscr{F}_{/X'}$가 존재하고,
\oldpage[I]{200}

유일하며, $v^\sharp=(u^\sharp)_{/X'}$를 만족한다. 삼중항
$(\mathscr{F},X,X')$들($\mathscr{F}$는 연접 $\mathscr{O}_X$-가군이고
$X'$은 $X$의 닫힌 부분집합이다)을 하나의 \emph{범주}의 대상으로
생각하고, 사상
$(\mathscr{F},X,X')\to(\mathscr{G},Y,Y')$들이
$f(X')\subset Y'$인 준스킴의 사상 $f:X\to Y$와
$f$-사상 $u:\mathscr{G}\to\mathscr{F}$로 이루어진다고 하자. 그러면
$(X_{/X'},\mathscr{F}_{/X'})$는 $(\mathscr{F},X,X')$에 관한
\emph{함자}라고 말할 수 있다. 이 함자는 국소 뇌터 형식적 준스킴
$\mathfrak{Z}$와 $\mathscr{O}_{\mathfrak{Z}}$-가군 $\mathscr{H}$로
이루어진 쌍 $(\mathfrak{Z},\mathscr{H})$들의 범주에 값을 갖는다.
후자의 범주에서 사상은 형식적 준스킴의 사상 $g$와 $g$-사상으로
이루어진 쌍이다.
\end{env}

\begin{proposition}[10.9.7]
\label{I.10.9.7-ko}
$S$, $X$, $Y$를 세 국소 뇌터 준스킴,
$g:X\to S$, $h:Y\to S$를 두 사상, $S'$을 $S$의 닫힌 부분집합,
$X'$ (각각 $Y'$)을 $g(X')\subset S'$ (각각
$h(Y')\subset S'$)을 만족하는 $X$ (각각 $Y$)의 닫힌 부분집합이라
하자~; $Z=X\times_S Y$라 하고, $Z$가 국소 뇌터라고 가정하며,
$Z'=p^{-1}(X')\cap q^{-1}(Y')$라 하자. 여기서 $p$와 $q$는
$X\times_S Y$의 사영이다. 이 조건에서 완비화 $Z_{/Z'}$는
$S_{/S'}$-형식적 준스킴들의 곱
$(X_{/X'})\times_{S_{/S'}}(Y_{/Y'})$와 동일시된다. 구조 사상들은
각각 $\widehat g$, $\widehat h$와 동일시되고, 사영들은 각각
$\widehat p$, $\widehat q$와 동일시된다.
\end{proposition}

문제가 $S$, $X$, $Y$에 관하여 국소적임은 명백하므로,
$S=\operatorname{Spec}(A)$,
$X=\operatorname{Spec}(B)$, $Y=\operatorname{Spec}(C)$,
$S'=V(\mathfrak{J})$, $X'=V(\mathfrak{K})$,
$Y'=V(\mathfrak{L})$인 경우로 환원된다. 여기서
$\mathfrak{J}$, $\mathfrak{K}$, $\mathfrak{L}$은
$\varphi(\mathfrak{J})\subset\mathfrak{K}$ 및
$\psi(\mathfrak{J})\subset\mathfrak{L}$을 만족하는 세 아이디얼이고,
$\varphi:A\to B$와 $\psi:A\to C$는 각각 $g$와 $h$에 대응하는
준동형이다. 이때 $Z=\operatorname{Spec}(B\otimes_A C)$이고
$Z'=V(\mathfrak{M})$임을 알고 있다. 여기서 $\mathfrak{M}$은
아이디얼
$\operatorname{Im}(\mathfrak{K}\otimes_A C)
+\operatorname{Im}(B\otimes_A\mathfrak{L})$이다.
(\hyperref[I.10.7.2-ko]{10.7.2})에 의해, 결론은 완비 텐서곱
$(\widehat B\otimes_{\widehat A}\widehat C)^\wedge$
(여기서 $\widehat A$, $\widehat B$, $\widehat C$는 각각
$A$, $B$, $C$에 $\mathfrak{J}$-, $\mathfrak{K}$-,
$\mathfrak{L}$-준아딕 위상을 주어 얻는 분리 완비화이다)이
텐서곱 $B\otimes_A C$의 $\mathfrak{M}$-준아딕 위상에 관한
분리 완비화라는 사실에서 따른다
(\hyperref[0.7.7.2-ko]{0, 7.7.2}).

또한 $T$가 국소 뇌터 $S$-준스킴이고,
$u:T\to X$, $v:T\to Y$가 두 $S$-사상이며, $T'$이
$u(T')\subset X'$, $v(T')\subset Y'$을 만족하는 $T$의 닫힌
부분집합이면, 완비화한 것들 사이로의 연장
$((u,v)_S)^\wedge$는
$(\widehat u,\widehat v)_{S_{/S'}}$와 동일시됨을 유의하자.

\begin{corollary}[10.9.8]
\label{I.10.9.8-ko}
$X$, $Y$를 두 국소 뇌터 $S$-준스킴이라 하고,
$X\times_S Y$도 국소 뇌터라고 하자~; $S'$을 $S$의 닫힌
부분집합이라 하고, $X'$ (각각 $Y'$)을 $X$ (각각 $Y$)의 닫힌
부분집합이라 하되 그 구조 사상에 의한 상이 $S'$에 포함된다고 하자.
$f(X')\subset Y'$을 만족하는 모든 $S$-사상 $f:X\to Y$에 대하여,
그래프 사상 $\Gamma_{\widehat f}$는 $f$의 그래프 사상의 연장
$(\Gamma_f)^\wedge$와 동일시된다.
\end{corollary}

\begin{corollary}[10.9.9]
\label{I.10.9.9-ko}
$X$, $Y$를 두 국소 뇌터 준스킴, $f:X\to Y$를 사상,
$Y'$을 $Y$의 닫힌 부분집합이라 하고, $X'=f^{-1}(Y')$라 하자.
그러면 형식적 준스킴\footnote{\textbf{번역 주.} 인쇄 원문은 여기서
\emph{préschéma}라고 쓰지만, $X_{/X'}$는 정의상 형식적 준스킴이고
우변도 형식적 준스킴들의 곱이므로 \emph{préschéma formel}에서
\emph{formel}이 빠진 것으로 보아 바로잡았다.} $X_{/X'}$는 다음 가환도식
\[
  \label{I.10.9.9.1-ko}
  \xymatrix{
    X\ar[d]_f &
    X_{/X'}\ar[l]\ar[d]^{\widehat f}\\
    Y &
    Y_{/Y'}\ar[l]
  }
\]
에 의해 형식적 준스킴들의 곱
$X\times_Y(Y_{/Y'})$와 동일시된다.
\end{corollary}

$S$와 $S'$을 모두 $Y$로, $X$와 $X'$을 모두 $X$로 바꾸어
(\hyperref[I.10.9.7-ko]{10.9.7})을 적용하면 충분하다.
\oldpage[I]{201}

\begin{remark}[10.9.10]
\label{I.10.9.10-ko}
$X$가 합 $X_1\amalg X_2$
(\hyperref[subsection:I.3.1-ko]{3.1})이고, $X'$이 합집합
$X'_1\cup X'_2$이며, 여기서 $X'_i$가 $X_i$의 닫힌 부분집합이면
($i=1,2$), 곧바로
$X_{/X'}={X_1}_{/X'_1}\amalg {X_2}_{/X'_2}$임을 알 수 있다.
\end{remark}

\subsection{아핀 형식적 스킴 위의 연접층에 대한 응용.}
\label{subsection:I.10.10-ko}

\begin{env}[10.10.1]
\label{I.10.10.1-ko}
이 절 전체에서 $A$는 \emph{뇌터 아딕환}, $\mathfrak{J}$는 $A$의
정의 아이디얼을 나타낸다. $X=\operatorname{Spec}(A)$,
$\mathfrak{X}=\operatorname{Spf}(A)$라 하자. 후자의 바탕공간은
$V(\mathfrak{J})$와 동일시되며, 이는 $X$의 닫힌 부분집합이다
(\hyperref[I.10.1.2-ko]{10.1.2}). 또한 정의
(\hyperref[I.10.1.2-ko]{10.1.2})와 정의
(\hyperref[I.10.8.4-ko]{10.8.4})에서 \emph{아핀 형식적 스킴
$\mathfrak{X}$}는 \emph{완비화 $X_{/\mathfrak{X}}$}, 곧 아핀 스킴
$X$를 그 바탕공간의 닫힌 부분집합 $\mathfrak{X}$을 따라 완비화한
것과 동일함을 알 수 있다. 따라서 각 연접 $\mathscr{O}_X$-가군 $\mathscr{F}$에는
유한 생성 $\mathscr{O}_{\mathfrak{X}}$-가군
$\mathscr{F}_{/\mathfrak{X}}$가 대응하며, 이는 또한 위상환의 층
$\mathscr{O}_{\mathfrak{X}}$ 위의 위상 가군의 층이다. 한편 모든 연접
$\mathscr{O}_X$-가군 $\mathscr{F}$는 $\widetilde M$ 꼴이다. 여기서
$M$은 유한 생성 $A$-가군이다
(\hyperref[I.1.5.1-ko]{1.5.1})~; 우리는
$(\widetilde M)_{/\mathfrak{X}}=M^\Delta$로 두겠다. 또한
$u:M\to N$이 유한 생성 $A$-가군 사이의 $A$-준동형이면, 이에
준동형 $\widetilde u:\widetilde M\to\widetilde N$이 대응하고, 따라서
연속 준동형
$\widetilde u_{/\mathfrak{X}}:(\widetilde M)_{/\mathfrak{X}}\to
(\widetilde N)_{/\mathfrak{X}}$도 대응한다. 이를 $u^\Delta$로
나타내겠다. 곧바로
$(v\circ u)^\Delta=v^\Delta\circ u^\Delta$임을 알 수 있다~; 이로써
\emph{$M^\Delta$를 주는 가법적 공변 함자}를 정의하였다. 이 함자는
유한 생성 $A$-가군의 범주에서 유한 생성
$\mathscr{O}_{\mathfrak{X}}$-가군의 범주로 간다. $A$가 \emph{이산}환이면
$M^\Delta=\widetilde M$이다.
\end{env}

\begin{proposition}[10.10.2]
\label{I.10.10.2-ko}
\medskip\noindent
\begin{enumerate}
  \item[\rm{(i)}] $M^\Delta$는 $M$에 관한 완전 함자이고, $A$-가군의
    함자적인 표준 동형
    $\Gamma(\mathfrak{X},M^\Delta)\xrightarrow{\sim}M$이 존재한다.
  \item[\rm{(ii)}] $M$과 $N$이 두 유한 생성 $A$-가군이면, 다음
    함자적인 표준 동형들이 존재한다.
    \[
      (M\otimes_A N)^\Delta
      \xrightarrow{\sim}
      M^\Delta\otimes_{\mathscr{O}_{\mathfrak{X}}}N^\Delta
      \tag{10.10.2.1}\label{I.10.10.2.1-ko}
    \]
    \[
      (\operatorname{Hom}_A(M,N))^\Delta
      \xrightarrow{\sim}
      \mathscr{H}\!om_{\mathscr{O}_{\mathfrak{X}}}
      (M^\Delta,N^\Delta)
      \tag{10.10.2.2}\label{I.10.10.2.2-ko}
    \]
  \item[\rm{(iii)}] 사상 $u\to u^\Delta$는 다음 함자적인 동형이다.
    \[
      \operatorname{Hom}_A(M,N)
      \xrightarrow{\sim}
      \operatorname{Hom}_{\mathscr{O}_{\mathfrak{X}}}
      (M^\Delta,N^\Delta)
      \tag{10.10.2.3}\label{I.10.10.2.3-ko}
    \]
\end{enumerate}
\end{proposition}

$M^\Delta$를 주는 함자의 완전성은 $\widetilde M$을 주는 함자
(\hyperref[I.1.3.5-ko]{1.3.5})와 $\mathscr{F}_{/\mathfrak{X}}$
(\hyperref[I.10.8.8-ko]{10.8.8})의 완전성에서 따른다. 정의에 따라
$\Gamma(\mathfrak{X},M^\Delta)$는 $A$-가군
$\Gamma(X,\widetilde M)=M$의 $\mathfrak{J}$-준아딕 위상에 관한 분리
완비화이다~; 그러나 $A$가 완비이고 $M$이 유한 생성인 까닭에,
$M$이 분리이고 완비임을 알고 있다
(\hyperref[0.7.3.6-ko]{0, 7.3.6}). 이로써 (i)의 증명이 끝난다. 동형
(\hyperref[I.10.10.2.1-ko]{10.10.2.1}) (각각
\hyperref[I.10.10.2.2-ko]{10.10.2.2})는 동형
(\hyperref[I.1.3.12-ko]{1.3.12, (i)})와
(\hyperref[I.10.8.10.1-ko]{10.8.10.1}) (각각
(\hyperref[I.1.3.12-ko]{1.3.12, (ii)})와
(\hyperref[I.10.8.10.2-ko]{10.8.10.2}))의 합성에서 얻어진다. 끝으로
$\operatorname{Hom}_A(M,N)$은 유한 생성 $A$-가군이므로 여기에 (i)를
적용할 수 있다. 그러면
$\Gamma(\mathfrak{X},(\operatorname{Hom}_A(M,N))^\Delta)$가
$\operatorname{Hom}_A(M,N)$과 동일시되고,
(\hyperref[I.10.10.2.2-ko]{10.10.2.2})를 이용하면 준동형
(\hyperref[I.10.10.2.3-ko]{10.10.2.3})이 동형임을 알 수 있다.

(\hyperref[I.10.10.2-ko]{10.10.2})에서
(\hyperref[I.1.3.7-ko]{1.3.7})과
(\hyperref[I.1.3.12-ko]{1.3.12})에서 얻은 것들과 유사한 일련의
귀결들을 얻으며, 그 정식화는 독자에게 맡긴다.
\oldpage[I]{202}

$M^\Delta$의 완전성을 완전열
$0\to\mathfrak{J}\to A\to A/\mathfrak{J}\to0$에 적용하면, 여기서
$\mathscr{O}_{\mathfrak{X}}$의 아이디얼층 중
$\mathfrak{J}^\Delta$로 나타낸 것이
(\hyperref[I.10.3.1-ko]{10.3.1})에서 같은 기호로 나타낸
아이디얼층과 일치함을 알 수 있다는 점에 유의하자. 이는
(\hyperref[I.10.3.2-ko]{10.3.2})에 따른다.

\begin{proposition}[10.10.3]
\label{I.10.10.3-ko}
(\hyperref[I.10.10.1-ko]{10.10.1})의 가정 아래에서
$\mathscr{O}_{\mathfrak{X}}$는 연접인 환의 층이다.
\end{proposition}

$f\in A$이면 $A_{\{f\}}$가 뇌터 아딕환임을 알고 있다
(\hyperref[0.7.6.11-ko]{0, 7.6.11}). 문제는 국소적이므로
(\hyperref[I.10.1.4-ko]{10.1.4})에 따라 준동형
$v:\mathscr{O}_{\mathfrak{X}}^n\to\mathscr{O}_{\mathfrak{X}}$의 핵이
유한 생성 $\mathscr{O}_{\mathfrak{X}}$-가군임을 보이면 충분하다.
이때 $v=u^\Delta$이다. 여기서 $u$는 $A$-준동형 $A^n\to A$이다
(\hyperref[I.10.10.2-ko]{10.10.2}). $A$가 뇌터이므로 $u$의 핵은
유한 생성이다. 다시 말해, 어떤 준동형
$A^m\xrightarrow{w}A^n$이 있어서 열
$A^m\xrightarrow{w}A^n\xrightarrow{u}A$가 완전하다. 따라서
(\hyperref[I.10.10.2-ko]{10.10.2})에 의해 열
$\mathscr{O}_{\mathfrak{X}}^m\xrightarrow{w^\Delta}
\mathscr{O}_{\mathfrak{X}}^n\xrightarrow{v}
\mathscr{O}_{\mathfrak{X}}$도 완전하다. 이로써 $v$의 핵이 유한
생성임을 보였다.

\begin{env}[10.10.4]
\label{I.10.10.4-ko}
앞의 표기 아래에서 $A_n=A/\mathfrak{J}^{n+1}$로 두고, $X_n$을 아핀
스킴
$\operatorname{Spec}(A_n)=(\mathfrak{X},
\mathscr{O}_{\mathfrak{X}}/\mathscr{J}^{n+1})$이라 하자. 여기서
$\mathscr{J}=\mathfrak{J}^\Delta$는 아이디얼 $\mathfrak{J}$에 대응하는
$\mathscr{O}_{\mathfrak{X}}$의 정의 아이디얼층이다. $m\leq n$일 때
표준 준동형 $A_n\to A_m$에 대응하는 준스킴의 사상 $X_m\to X_n$을
$u_{mn}$이라 하자~; 형식적 스킴 $\mathfrak{X}$는 $u_{mn}$들에 관한
$X_n$들의 귀납극한이다
(\hyperref[I.10.6.3-ko]{10.6.3}).
\end{env}

\begin{proposition}[10.10.5]
\label{I.10.10.5-ko}
(\hyperref[I.10.10.1-ko]{10.10.1})의 가정 아래에서
$\mathscr{F}$를 $\mathscr{O}_{\mathfrak{X}}$-가군이라 하자. 다음 조건들은
서로 동치이다~:
\begin{enumerate}
  \item[a)] $\mathscr{F}$는 연접 $\mathscr{O}_{\mathfrak{X}}$-가군이다.
  \item[b)] $\mathscr{F}$는 열 $(\mathscr{F}_n)$의 사영극한
    (\hyperref[I.10.6.6-ko]{10.6.6})과 동형이다. 이 열은 연접
    $\mathscr{O}_{X_n}$-가군들로 이루어지고
    $u_{mn}^*(\mathscr{F}_n)=\mathscr{F}_m$을 만족한다.
  \item[c)] 유한 생성 $A$-가군 $M$이 존재한다. 그 가군은
    (\hyperref[I.10.10.2-ko]{10.10.2, (i)})에 의해 표준 동형을
    제외하고 결정되며, $\mathscr{F}$는 $M^\Delta$와 동형이다.
\end{enumerate}
\end{proposition}

먼저 b)가 c)를 함의함을 보이자. $\mathscr{F}_n=\widetilde{M_n}$이고,
여기서 $M_n$은 유한 생성 $A_n$-가군이다. 가정에 의해
$M_m=M_n\otimes_{A_n}A_m$이 $m\leq n$일 때 성립한다
(\hyperref[I.1.6.5-ko]{1.6.5})~; 따라서 $M_n$들은 표준 디-준동형
$M_n\to M_m$ ($m\leq n$)에 관하여 사영계를 이룬다. $A_n$들의
정의로부터 이 사영계가
(\hyperref[0.7.2.9-ko]{0, 7.2.9})의 조건들을 만족함이 곧바로 따른다~;
그러므로 그 사영극한 $M$은 유한 생성 $A$-가군이고,
$M_n=M\otimes_A A_n$이 모든 $n$에 대하여 성립한다. 따라서
$\mathscr{F}_n$은 $X_n$ 위에서
$\widetilde M\otimes_{\mathscr{O}_X}
(\mathscr{O}_X/\widetilde{\mathfrak{J}}^{n+1})$에 의해 유도되는
가군이고, 정의 (\hyperref[I.10.8.4-ko]{10.8.4})에 의해
$\mathscr{F}=M^\Delta$이다.

역으로 c)는 b)를 함의한다. 실제로 $u_n$이 몰입 사상 $X_n\to X$이면,
$u_n^*(\widetilde M)=(M\otimes_A A_n)^{\sim}$은 $X_n$ 위에서
$\widetilde M\otimes_{\mathscr{O}_X}
(\mathscr{O}_X/\widetilde{\mathfrak{J}}^{n+1})$에 의해 유도되는
가군이고, 정의 (\hyperref[I.10.8.4-ko]{10.8.4})에 의해
$M^\Delta=\varprojlim u_n^*(\widetilde M)$이다~;
$u_m=u_n\circ u_{mn}$이 $m\leq n$일 때 성립하므로,
$\mathscr{F}_n=u_n^*(\widetilde M)$들은 b)의 조건들을 만족한다. 이로써
주장이 따른다.

이제 c)가 a)를 함의함을 보이자. 정의에 의해
$\mathscr{O}_{\mathfrak{X}}=A^\Delta$이다~; $M$은 어떤 준동형
$A^m\to A^n$의 여핵이므로, (\hyperref[I.10.10.2-ko]{10.10.2})에 의해
$M^\Delta$는 준동형
$\mathscr{O}_{\mathfrak{X}}^m\to\mathscr{O}_{\mathfrak{X}}^n$의
여핵이다. $\mathscr{O}_{\mathfrak{X}}$가 연접인 환의 층이므로
(\hyperref[I.10.10.3-ko]{10.10.3}), $M^\Delta$도 연접이다
(\hyperref[0.5.3.4-ko]{0, 5.3.4}).
\oldpage[I]{203}

끝으로 a)는 b)를 함의한다. $\mathscr{O}_{\mathfrak{X}}$-가군으로
보면,
$\mathscr{O}_{X_n}=\mathscr{O}_{\mathfrak{X}}/
\mathscr{J}^{n+1}=A_n^\Delta$이다~; $\mathscr{F}_n=\mathscr{F}
\otimes_{\mathscr{O}_{\mathfrak{X}}}\mathscr{O}_{X_n}$은 연접
$\mathscr{O}_{\mathfrak{X}}$-가군이다
(\hyperref[0.5.3.5-ko]{0, 5.3.5}). 또한 이것은
$\mathscr{O}_{X_n}$-가군이고 $\mathscr{J}^{n+1}$이 연접이므로,
$\mathscr{F}_n$은 연접 $\mathscr{O}_{X_n}$-가군이다
(\hyperref[0.5.3.10-ko]{0, 5.3.10}).
$u_{mn}^*(\mathscr{F}_n)=\mathscr{F}_m$이 $m\leq n$일 때 성립함도
곧바로 알 수 있다
($X_m\to X_n$에서 얻는 바탕공간 사이의 연속사상은
$\mathfrak{X}$의 항등사상임을 기억하면 된다). 따라서 층
$\mathscr{G}=\varprojlim\mathscr{F}_n$은 연접
$\mathscr{O}_{\mathfrak{X}}$-가군이다. 이는 b)가 a)를 함의함을 이미
보였기 때문이다. 표준 준동형 $\mathscr{F}\to\mathscr{F}_n$들은
사영계를 이루며, 사영극한을 취하면 표준 준동형
$w:\mathscr{F}\to\mathscr{G}$를 얻는다. 이제 남은 것은 $w$가 전단사임을
보이는 일뿐이다. 이 문제는 국소적이므로 $\mathscr{F}$가 어떤 준동형
$\mathscr{O}_{\mathfrak{X}}^p\to\mathscr{O}_{\mathfrak{X}}^q$의
여핵인 경우로 국한할 수 있다~; 이 준동형은 $v^\Delta$ 꼴이고, 여기서
$v$는 준동형 $A^p\to A^q$이다\footnote{\textbf{번역 주.} 인쇄 원문은
앞의 층 준동형을
$\mathscr{O}_{\mathfrak{X}}^p\to\mathscr{O}_{\mathfrak{X}}^q$로
둔 뒤, 이를 유도하는 $v$를 $A^m\to A^n$으로 쓴다. 그러나
$\Delta$ 대응에서는 $(A^r)^\Delta=\mathscr{O}_{\mathfrak{X}}^r$이므로
유한 자유 가군의 계수(階數)가 보존된다. 따라서 $v:A^p\to A^q$로
바로잡았다.}
(\hyperref[I.10.10.2-ko]{10.10.2}). 그러므로 $\mathscr{F}$는
$M^\Delta$와 동형이다. 여기서 $M=\operatorname{Coker}v$이다
(\hyperref[I.10.10.2-ko]{10.10.2}). 이제
(\hyperref[I.10.10.2-ko]{10.10.2})에 의해
$\mathscr{F}_n=M^\Delta\otimes_{\mathscr{O}_{\mathfrak{X}}}A_n^\Delta
=(M\otimes_A A_n)^\Delta$이다. 또한 $\mathfrak{J}$-아딕 위상은
$M\otimes_A A_n$ 위에서 이산이므로,
$(M\otimes_A A_n)^\Delta=(M\otimes_A A_n)^{\sim}$이다
($\mathscr{O}_{X_n}$-가군으로서). 위에서
$M^\Delta=\varprojlim\mathscr{F}_n$임을 보았으므로, 이 경우 $w$는
실제로 항등사상이다.
\par\hfill C.Q.F.D.\par

\begin{corollary}[10.10.6]
\label{I.10.10.6-ko}
$\mathscr{F}$가 (\hyperref[I.10.10.5-ko]{10.10.5})의 조건 b)를
만족하면, 사영계 $(\mathscr{F}_n)$은
$\mathscr{F}\otimes_{\mathscr{O}_{\mathfrak{X}}}\mathscr{O}_{X_n}$들로
이루어진 계와 동형이다.
\end{corollary}

\begin{env}[10.10.7]
\label{I.10.10.7-ko}
이제 $A$, $B$를 두 뇌터 아딕환이라 하고,
$\varphi:B\to A$를 연속 준동형이라 하자. $\mathfrak{J}$ (각각
$\mathfrak{K}$)를 $A$ (각각 $B$)의 정의 아이디얼이라 하되
$\varphi(\mathfrak{K})\subset\mathfrak{J}$라고 하고,
$X=\operatorname{Spec}(A)$, $Y=\operatorname{Spec}(B)$,
$\mathfrak{X}=\operatorname{Spf}(A)$,
$\mathfrak{Y}=\operatorname{Spf}(B)$라 하자. $f:X\to Y$를
$\varphi$에 대응하는 준스킴의 사상이라 하고
(\hyperref[I.1.6.1-ko]{1.6.1}), 완비화한 것들 사이로의 연장을
$\widehat f:\mathfrak{X}\to\mathfrak{Y}$라 하자
(\hyperref[I.10.9.1-ko]{10.9.1}). 이는 또한 $\varphi$에 대응하는
형식적 준스킴의 사상이다 (\hyperref[I.10.2.2-ko]{10.2.2}).
\end{env}

\begin{proposition}[10.10.8]
\label{I.10.10.8-ko}
모든 유한 생성 $B$-가군 $N$에 대하여 다음과 같은
$\mathscr{O}_{\mathfrak{X}}$-가군의 함자적인 표준 동형이 존재한다~:
\[
  \widehat f^*(N^\Delta)
  \xrightarrow{\sim}
  (N\otimes_B A)^\Delta.
\]
\end{proposition}

실제로 표준 사상들을 $i_X:\mathfrak{X}\to X$와
$i_Y:\mathfrak{Y}\to Y$로 나타내면, 함자적인 표준 동형을 통해
동일시할 때 (\hyperref[I.10.8.8-ko]{10.8.8})
$N^\Delta=i_Y^*(\widetilde N)$이고
\[
  (N\otimes_B A)^\Delta
  =i_X^*((N\otimes_B A)^{\sim})
  =i_X^*(f^*(\widetilde N))
\]
이다 (\hyperref[I.1.6.5-ko]{1.6.5})~; 따라서 명제는 도표
(\hyperref[I.10.9.2-ko]{10.9.2})의 가환성에서 따른다.

\begin{corollary}[10.10.9]
\label{I.10.10.9-ko}
$\mathfrak b$를 $B$의 임의의 아이디얼이라 하면,
$\widehat f^*(\mathfrak b^\Delta)\mathscr{O}_{\mathfrak{X}}
=(\mathfrak bA)^\Delta$이다.
\end{corollary}

실제로 $j$를 표준 단사 준동형 $\mathfrak b\to B$라 하자. 이에 대응하는
$j^\Delta:\mathfrak b^\Delta\to\mathscr{O}_{\mathfrak{Y}}$는
$\mathscr{O}_{\mathfrak{Y}}$-가군의 표준 단사 준동형이다~;
정의에 따라
$\widehat f^*(\mathfrak b^\Delta)\mathscr{O}_{\mathfrak{X}}$는 준동형
$\widehat f^*(j^\Delta):\widehat f^*(\mathfrak b^\Delta)\to
\mathscr{O}_{\mathfrak{X}}=\widehat f^*(\mathscr{O}_{\mathfrak{Y}})$의
상이다~; 그런데 (\hyperref[I.10.10.8-ko]{10.10.8})에 의해 이 준동형은
$(j\otimes1)^\Delta:(\mathfrak b\otimes_B A)^\Delta\to
\mathscr{O}_{\mathfrak{X}}=(B\otimes_B A)^\Delta$와 동일시된다.
$j\otimes1$의 상은 아이디얼 $\mathfrak bA$이고 이는 $A$의 아이디얼이므로,
(\hyperref[I.10.10.2-ko]{10.10.2})에 따라 $(j\otimes1)^\Delta$의 상은
$(\mathfrak bA)^\Delta$이다. 이로써 결론이 따른다.
\oldpage[I]{204}

\subsection{형식적 준스킴 위의 연접층.}
\label{subsection:I.10.11-ko}

\begin{proposition}[10.11.1]
\label{I.10.11.1-ko}
$\mathfrak{X}$가 국소 뇌터 형식적 준스킴이면,
$\mathscr{O}_{\mathfrak{X}}$는 연접인 환의 층이고 $\mathfrak{X}$의
모든 정의 아이디얼층은 연접이다.
\end{proposition}

문제는 국소적이므로 뇌터 아핀 형식적 스킴인 경우로 환원되며, 따라서
명제는 (\hyperref[I.10.10.3-ko]{10.10.3})과
(\hyperref[I.10.10.5-ko]{10.10.5})에서 따른다.

\begin{env}[10.11.2]
\label{I.10.11.2-ko}
$\mathfrak{X}$를 국소 뇌터 형식적 준스킴,
$\mathscr{J}$를 $\mathfrak{X}$의 정의 아이디얼층이라 하자. $X_n$을
(보통 의미의) 국소 뇌터 준스킴
$(\mathfrak{X},\mathscr{O}_{\mathfrak{X}}/\mathscr{J}^{n+1})$이라 하자.
그러면 $\mathfrak{X}$는 열 $(X_n)$의 귀납극한이고, 그 전이사상들은
표준 사상 $u_{mn}:X_m\to X_n$이다
(\hyperref[I.10.6.3-ko]{10.6.3}). 이 표기 아래에서 다음이 성립한다~:
\end{env}

\begin{theorem}[10.11.3]
\label{I.10.11.3-ko}
$\mathscr{O}_{\mathfrak{X}}$-가군 $\mathscr{F}$가 연접이기 위한
필요충분조건은 그것이 열 $(\mathscr{F}_n)$의 사영극한과 동형인 것이다.
여기서 각 $\mathscr{F}_n$은 연접
$\mathscr{O}_{X_n}$-가군이며,
$u_{mn}^*(\mathscr{F}_n)=\mathscr{F}_m$이 $m\leq n$일 때 성립한다
(\hyperref[I.10.6.6-ko]{10.6.6}). 이때 사영계 $(\mathscr{F}_n)$은
$u_n^*(\mathscr{F})=\mathscr{F}
\otimes_{\mathscr{O}_{\mathfrak{X}}}\mathscr{O}_{X_n}$들로 이루어진
계와 동형이다. 여기서 $u_n$은 표준 사상 $X_n\to\mathfrak{X}$이다.
\end{theorem}

문제는 국소적이므로 $\mathfrak{X}$가 뇌터 아핀 형식적 스킴인 경우로
환원되며, 이 경우 정리는 (\hyperref[I.10.10.5-ko]{10.10.5})와
(\hyperref[I.10.10.6-ko]{10.10.6})의 귀결이다.

따라서 \emph{연접 $\mathscr{O}_{\mathfrak{X}}$-가군을 주는 것은
사영계 $(\mathscr{F}_n)$를 주는 것과 동치이다. 여기서 각 항은 연접
$\mathscr{O}_{X_n}$-가군이고,
$u_{mn}^*(\mathscr{F}_n)=\mathscr{F}_m$이 $m\leq n$일 때 성립한다}.

\begin{corollary}[10.11.4]
\label{I.10.11.4-ko}
$\mathscr{F}$와 $\mathscr{G}$가 두 연접
$\mathscr{O}_{\mathfrak{X}}$-가군이면,
(\hyperref[I.10.11.3-ko]{10.11.3})의 표기를 쓰면 다음 함자적인 표준
동형을 정의할 수 있다~:
\[
  \label{I.10.11.4.1-ko}
  \operatorname{Hom}_{\mathscr{O}_{\mathfrak{X}}}
  (\mathscr{F},\mathscr{G})
  \xrightarrow{\sim}
  \varprojlim_n
  \operatorname{Hom}_{\mathscr{O}_{X_n}}
  (\mathscr{F}_n,\mathscr{G}_n).
  \tag{10.11.4.1}
\]
\end{corollary}

우변의 사영극한은 다음 사상들에 관하여 취한 것으로 이해해야 한다~:
$\theta_n\mapsto u_{mn}^*(\theta_n)$ ($m\leq n$). 이들은
$\operatorname{Hom}_{\mathscr{O}_{X_n}}
(\mathscr{F}_n,\mathscr{G}_n)$에서
$\operatorname{Hom}_{\mathscr{O}_{X_m}}
(\mathscr{F}_m,\mathscr{G}_m)$으로 가는 사상이다. 준동형
(\hyperref[I.10.11.4.1-ko]{10.11.4.1})은 원소
$\theta\in\operatorname{Hom}_{\mathscr{O}_{\mathfrak{X}}}
(\mathscr{F},\mathscr{G})$에 열 $(u_n^*(\theta))$을 대응시킨다~;
(\hyperref[I.10.11.3-ko]{10.11.3})을 고려하면, 사영계
$(\theta_n)\in\varprojlim_n
\operatorname{Hom}_{\mathscr{O}_{X_n}}
(\mathscr{F}_n,\mathscr{G}_n)$에
$\operatorname{Hom}_{\mathscr{O}_{\mathfrak{X}}}
(\mathscr{F},\mathscr{G})$ 안에서의 그 사영극한을 대응시킴으로써
앞의 준동형의 역준동형을 정의할 수 있음을 곧 알 수 있다.

\begin{corollary}[10.11.5]
\label{I.10.11.5-ko}
준동형 $\theta:\mathscr{F}\to\mathscr{G}$가 전사일 필요충분조건은
이에 대응하는 준동형
$\theta_0=u_0^*(\theta):\mathscr{F}_0\to\mathscr{G}_0$가 전사인 것이다.
\end{corollary}

문제는 국소적이므로
$\mathfrak{X}=\operatorname{Spf}(A)$이고 $A$가 뇌터 아딕환이며,
$\mathscr{F}=M^\Delta$, $\mathscr{G}=N^\Delta$, $\theta=u^\Delta$인
경우로 환원된다. 여기서 $M$과 $N$은 유한 생성 $A$-가군이고
$u$는 준동형 $M\to N$이다~; 또한 이때
$\theta_0=\widetilde{u_0}$이고, 여기서 $u_0$는 준동형
$u\otimes1:M\otimes_A A/\mathfrak{J}\to
N\otimes_A A/\mathfrak{J}$이다~;
결론은 $\theta$와 $u$가 동시에 전사이고, 또 $\theta_0$와
$u_0$가 동시에 전사라는 사실
(\hyperref[I.1.3.9-ko]{1.3.9} 및
\hyperref[I.10.10.2-ko]{10.10.2})과 $u$와 $u_0$가 동시에 전사라는
사실 (\hyperref[0.7.1.14-ko]{0, 7.1.14})에서 따른다.

\begin{env}[10.11.6]
\label{I.10.11.6-ko}
정리 (\hyperref[I.10.11.3-ko]{10.11.3})은 모든 연접
$\mathscr{O}_{\mathfrak{X}}$-가군 $\mathscr{F}$를
\emph{위상} $\mathscr{O}_{\mathfrak{X}}$-가군으로 볼 수 있음을
보여 준다. 즉 이를 \emph{의사이산} 군의 층 $\mathscr{F}_n$들의
사영극한으로 본다
(\hyperref[0.3.8.1-ko]{0, 3.8.1}). 따라서
(\hyperref[I.10.11.4-ko]{10.11.4})에 의해 준동형
$u:\mathscr{F}\to\mathscr{G}$가 연접
$\mathscr{O}_{\mathfrak{X}}$-가군 사이의 준동형이면 자동으로
\emph{연속}이다

\oldpage[I]{205}

(\hyperref[0.3.8.2-ko]{0, 3.8.2}). 또한 $\mathscr{H}$가 연접
$\mathscr{O}_{\mathfrak{X}}$-부분가군으로서 연접
$\mathscr{O}_{\mathfrak{X}}$-가군 $\mathscr{F}$에 들어 있으면, 모든 열린집합
$U\subset\mathfrak{X}$에 대하여 $\Gamma(U,\mathscr{H})$는 위상군
$\Gamma(U,\mathscr{F})$의 \emph{닫힌} 부분군이다. 실제로 함자
$\Gamma$가 왼쪽 완전이므로 $\Gamma(U,\mathscr{H})$는 준동형
$\Gamma(U,\mathscr{F})\to\Gamma(U,\mathscr{F}/\mathscr{H})$의
핵이다. 이 준동형은 앞서 본 바에 따라 \emph{연속}인데,
$\mathscr{F}/\mathscr{H}$가 연접이기 때문이다
(\hyperref[0.5.3.4-ko]{0, 5.3.4})~; 따라서 이 주장은
$\Gamma(U,\mathscr{F}/\mathscr{H})$가 분리 위상군이라는 사실에서
따른다.
\end{env}

\begin{proposition}[10.11.7]
\label{I.10.11.7-ko}
$\mathscr{F}$와 $\mathscr{G}$를 두 연접
$\mathscr{O}_{\mathfrak{X}}$-가군이라 하자.
(\hyperref[I.10.11.3-ko]{10.11.3})의 표기를 쓰면 다음 위상
$\mathscr{O}_{\mathfrak{X}}$-가군의 함자적인 표준 동형들을 정의할 수 있다
(\hyperref[I.10.11.6-ko]{10.11.6})~:
\[
  \label{I.10.11.7.1-ko}
  \mathscr{F}\otimes_{\mathscr{O}_{\mathfrak{X}}}\mathscr{G}
  \xrightarrow{\sim}
  \varprojlim_n
  (\mathscr{F}_n\otimes_{\mathscr{O}_{X_n}}\mathscr{G}_n)
  \tag{10.11.7.1}
\]
\[
  \label{I.10.11.7.2-ko}
  \mathscr{H}\!om_{\mathscr{O}_{\mathfrak{X}}}
  (\mathscr{F},\mathscr{G})
  \xrightarrow{\sim}
  \varprojlim_n
  \mathscr{H}\!om_{\mathscr{O}_{X_n}}(\mathscr{F}_n,\mathscr{G}_n)
  \tag{10.11.7.2}
\]
\end{proposition}

동형 (\hyperref[I.10.11.7.1-ko]{10.11.7.1})의 존재는 공식
\[
  \mathscr{F}_n\otimes_{\mathscr{O}_{X_n}}\mathscr{G}_n
  =
  (\mathscr{F}\otimes_{\mathscr{O}_{\mathfrak{X}}}
  \mathscr{O}_{X_n})
  \otimes_{\mathscr{O}_{X_n}}
  (\mathscr{G}\otimes_{\mathscr{O}_{\mathfrak{X}}}
  \mathscr{O}_{X_n})
  =
  (\mathscr{F}\otimes_{\mathscr{O}_{\mathfrak{X}}}\mathscr{G})
  \otimes_{\mathscr{O}_{\mathfrak{X}}}\mathscr{O}_{X_n}
\]
과 (\hyperref[I.10.11.3-ko]{10.11.3})에서 따른다. 두 변을 위상을
고려하지 않은 가군의 층으로 보았을 때, 동형
(\hyperref[I.10.11.7.2-ko]{10.11.7.2})은
$\mathscr{H}\!om_{\mathscr{O}_{\mathfrak{X}}}
(\mathscr{F},\mathscr{G})$와
$\mathscr{H}\!om_{\mathscr{O}_{X_n}}(\mathscr{F}_n,\mathscr{G}_n)$의
단면의 정의와 $\mathfrak{X}$의 임의의 뇌터 아핀 형식적 열린집합 위에
유도된 준스킴에 (\hyperref[I.10.11.4.1-ko]{10.11.4.1})을 적용하여 얻는
동형의 존재에서 따른다. 이제 준콤팩트 집합 위에서 동형
(\hyperref[I.10.11.7.2-ko]{10.11.7.2})이 쌍연속임을 보이는 것만 남는다.
따라서 $\mathfrak{X}=\operatorname{Spf}(A)$이고 $A$가 뇌터 아딕환인
경우로 환원되며, 이때 (\hyperref[I.10.10.5-ko]{10.10.5})에 따라
$\mathscr{F}=M^\Delta$, $\mathscr{G}=N^\Delta$이고 $M$, $N$은 유한 생성
$A$-가군이다~; (\hyperref[I.10.10.2.1-ko]{10.10.2.1}),
(\hyperref[I.10.10.2.3-ko]{10.10.2.3}) 및
(\hyperref[I.1.3.12-ko]{1.3.12, (ii)})을 고려하면, 표준 동형
$\operatorname{Hom}_A(M,N)\xrightarrow{\sim}
\varprojlim_n\operatorname{Hom}_{A_n}(M_n,N_n)$
($M_n=M\otimes_A A_n$, $N_n=N\otimes_A A_n$)이 연속임을 보이는 문제로
환원되는데, 이는 (\hyperref[0.7.8.2-ko]{0, 7.8.2})에서 증명되었다.

\begin{env}[10.11.8]
\label{I.10.11.8-ko}
$\operatorname{Hom}_{\mathscr{O}_{\mathfrak{X}}}
(\mathscr{F},\mathscr{G})$는 위상군의 층
$\mathscr{H}\!om_{\mathscr{O}_{\mathfrak{X}}}
(\mathscr{F},\mathscr{G})$의 단면군이므로 군 위상을 갖는다.
$\mathfrak{X}$가 \emph{뇌터}이면,
(\hyperref[I.10.11.7.2-ko]{10.11.7.2})에 따라 이 군에서 $0$의 근방들의
기본계는 부분군
$\operatorname{Hom}_{\mathscr{O}_{\mathfrak{X}}}
(\mathscr{F},\mathscr{J}^n\mathscr{G})$ ($n$은 임의)을 취하여 얻는다.
\end{env}

\begin{proposition}[10.11.9]
\label{I.10.11.9-ko}
$\mathfrak{X}$를 뇌터 형식적 준스킴,
$\mathscr{F}$와 $\mathscr{G}$를 두 연접
$\mathscr{O}_{\mathfrak{X}}$-가군이라 하자. 위상군
$\operatorname{Hom}_{\mathscr{O}_{\mathfrak{X}}}
(\mathscr{F},\mathscr{G})$에서 전사인 준동형들(각각 단사인 준동형들,
전단사인 준동형들)은 열린 부분집합을 이룬다.
\end{proposition}

(\hyperref[I.10.11.5-ko]{10.11.5})에 따라
$\operatorname{Hom}_{\mathscr{O}_{\mathfrak{X}}}
(\mathscr{F},\mathscr{G})$ 안의 전사 준동형들의 집합은 다음 연속사상
$\operatorname{Hom}_{\mathscr{O}_{\mathfrak{X}}}
(\mathscr{F},\mathscr{G})\to\allowbreak
\operatorname{Hom}_{\mathscr{O}_{X_0}}(\mathscr{F}_0,\mathscr{G}_0)$에
의한 이산군
$\operatorname{Hom}_{\mathscr{O}_{X_0}}(\mathscr{F}_0,\mathscr{G}_0)$의
어떤 부분집합의 역상이므로 첫 번째 주장이 따른다. 두 번째 주장을
증명하기 위해 $\mathfrak{X}$를 유한 개의 뇌터 아핀 형식적 열린집합
$U_i$로 덮자. $\theta\in
\operatorname{Hom}_{\mathscr{O}_{\mathfrak{X}}}
(\mathscr{F},\mathscr{G})$가 단사일 필요충분조건은 각 연속 제한사상
$\operatorname{Hom}_{\mathscr{O}_{\mathfrak{X}}}
(\mathscr{F},\mathscr{G})\to
\operatorname{Hom}_{\mathscr{O}_{\mathfrak{X}}|U_i}
(\mathscr{F}|U_i,\mathscr{G}|U_i)$ 아래에서 얻는 상이 모두 단사인 것이다~;
따라서 아핀인 경우로 환원되며, 이 경우에는 이미
(\hyperref[0.7.8.3-ko]{0, 7.8.3})에서 증명되었다.

\oldpage[I]{206}

\subsection{형식적 준스킴의 아딕 사상.}
\label{subsection:I.10.12-ko}

\begin{env}[10.12.1]
\label{I.10.12.1-ko}
$\mathfrak{X}$, $\mathfrak{S}$를 두 \emph{국소 뇌터} 형식적
준스킴이라 하자~; 사상 $f:\mathfrak{X}\to\mathfrak{S}$가
\emph{아딕}이라 함은 아이디얼 $\mathscr{J}$가 존재하여 그것이
$\mathfrak{S}$의 정의 아이디얼이고
$\mathscr{K}=f^*(\mathscr{J})\mathscr{O}_{\mathfrak{X}}$가
$\mathfrak{X}$의 정의 아이디얼인 것을 뜻한다~; 이때
$\mathfrak{X}$를
\emph{$\mathfrak{S}$-아딕 형식적 준스킴}이라고도 한다
($f$에 관하여). 이 경우 $\mathscr{J}_1$을 $\mathfrak{S}$의
\emph{임의의} 정의 아이디얼이라 하면,
$\mathscr{K}_1=f^*(\mathscr{J}_1)\mathscr{O}_{\mathfrak{X}}$는
$\mathfrak{X}$의 정의 아이디얼이다. 실제로 문제는 국소적이므로
$\mathfrak{X}$와 $\mathfrak{S}$가 뇌터 아핀 형식적 준스킴이라고
가정해도 된다~; 따라서 정수 $n$이 존재하여
$\mathscr{J}^n\subset\mathscr{J}_1$이고
$\mathscr{J}_1^n\subset\mathscr{J}$이다
(\hyperref[I.10.3.6-ko]{10.3.6} 및
\hyperref[0.7.1.4-ko]{0, 7.1.4}). 그러므로
$\mathscr{K}^n\subset\mathscr{K}_1$이고
$\mathscr{K}_1^n\subset\mathscr{K}$이다. 첫 번째 관계는
$\mathscr{K}_1=\mathfrak{K}_1^\Delta$임을 보인다. 여기서
$\mathfrak{K}_1$은
$A=\Gamma(\mathfrak{X},\mathscr{O}_{\mathfrak{X}})$의 열린
아이디얼이다. 두 번째 관계는 $\mathfrak{K}_1$이 $A$의 정의
아이디얼임을 보이므로
(\hyperref[0.7.1.4-ko]{0, 7.1.4}), 이로써 주장이 따른다.
\end{env}

앞의 내용에서 $\mathfrak{X}$와 $\mathfrak{Y}$가 두
$\mathfrak{S}$-아딕 형식적 준스킴이면,
\emph{모든 $\mathfrak{S}$-사상
$u:\mathfrak{X}\to\mathfrak{Y}$은 아딕이다}라는 것이 곧바로 따른다~:
실제로 $f:\mathfrak{X}\to\mathfrak{S}$,
$g:\mathfrak{Y}\to\mathfrak{S}$를 각각 구조 사상이라 하고,
$\mathscr{J}$를 $\mathfrak{S}$의 정의 아이디얼이라 하자. 그러면
$f=g\circ u$이므로
$u^*(g^*(\mathscr{J})\mathscr{O}_{\mathfrak{Y}})
\mathscr{O}_{\mathfrak{X}}=f^*(\mathscr{J})
\mathscr{O}_{\mathfrak{X}}$는 $\mathfrak{X}$의 정의 아이디얼이고,
가정에 따라 $g^*(\mathscr{J})\mathscr{O}_{\mathfrak{Y}}$는
$\mathfrak{Y}$의 정의 아이디얼이다.

\begin{env}[10.12.2]
\label{I.10.12.2-ko}
이하에서는 국소 뇌터 형식적 준스킴 $\mathfrak{S}$와 아이디얼
$\mathscr{J}$를 고정하되, 후자는 $\mathfrak{S}$의 정의 아이디얼이라고
가정하자~; 또한
$S_n=(\mathfrak{S},\mathscr{O}_{\mathfrak{S}}/\mathscr{J}^{n+1})$로
두겠다. $\mathfrak{S}$-아딕 형식적 준스킴들(국소 뇌터인 것들)은
자명하게 하나의 \emph{범주}를 이룬다. 귀납계 $(X_n)$이 (보통 의미의)
국소 뇌터 $S_n$-준스킴들로 이루어져 있을 때, 다음 조건을 만족하면
이를 \emph{아딕 $(S_n)$-귀납계}라 하겠다. 구조 사상
$f_n:X_n\to S_n$들에 대하여 $m\leq n$이면 도식
\[
  \label{I.10.12.2.1-ko}
  \xymatrix{
    X_n \ar[d]_{f_n}
    & X_m \ar[l] \ar[d]^{f_m}\\
    S_n
    & S_m \ar[l]
  }
  \tag{10.12.2.1}
\]
은 가환이고, \emph{$X_m$을 곱
$X_n\times_{S_n}S_m=(X_n)_{(S_m)}$과 동일시한다}. 아딕 귀납계들도
하나의 \emph{범주}를 이룬다~: 실제로 이러한 계들의 사상
$(X_n)\to(Y_n)$을 \emph{$S_n$-사상 $u_n:X_n\to Y_n$들로 이루어진
귀납계}로 정의하되, $u_m$이 $(u_n)_{(S_m)}$과 동일시되는 조건을
모든 $m\leq n$에 대하여 만족하도록 하면 충분하다. 그러면~:
\end{env}

\begin{theorem}[10.12.3]
\label{I.10.12.3-ko}
$\mathfrak{S}$-아딕 형식적 준스킴들의 범주와 아딕
$(S_n)$-귀납계들의 범주 사이에는 표준 동치가 있다.
\end{theorem}

이 동치는 다음과 같이 얻는다~: $\mathfrak{X}$가
$\mathfrak{S}$-아딕 형식적 준스킴이고,
$f:\mathfrak{X}\to\mathfrak{S}$가 구조 사상이라 하면,
$\mathscr{K}=f^*(\mathscr{J})\mathscr{O}_{\mathfrak{X}}$는
$\mathfrak{X}$의 정의 아이디얼이며, $\mathfrak{X}$에 귀납계
$X_n=(\mathfrak{X},\mathscr{O}_{\mathfrak{X}}/\mathscr{K}^{n+1})$을
대응시킨다. 여기서 $f_n:X_n\to S_n$은 $f$에 대응하는 구조 사상이다
(\hyperref[I.10.5.6-ko]{10.5.6}). 먼저 $(X_n)$이
\emph{아딕 귀납계}임을 보이자~: $f=(\psi,\theta)$라 쓰면
$\psi^*(\mathscr{J})\mathscr{O}_{\mathfrak{X}}=\mathscr{K}$이고, 따라서
$\psi^*(\mathscr{J}^n)\mathscr{O}_{\mathfrak{X}}=\mathscr{K}^n$이며, 이는
모든 $n$에 대하여 성립한다.
또한 함자 $\psi^*$의 완전성에 의해
$\mathscr{K}^{m+1}/\mathscr{K}^{n+1}
=\psi^*(\mathscr{J}^{m+1}/\mathscr{J}^{n+1})
(\mathscr{O}_{\mathfrak{X}}/\mathscr{K}^{n+1})$이 $m\leq n$일 때
성립한다. 그러므로 원하는 결론은
(\hyperref[I.4.4.5-ko]{4.4.5})에서 따른다. 더욱이
$\mathfrak{S}$-사상 $u:\mathfrak{X}\to\mathfrak{Y}$이 두
$\mathfrak{S}$-아딕 형식적 준스킴 사이의 사상일 때에는, 명백한 표기를 쓰면,

\oldpage[I]{207}

$S_n$-사상 $u_n:X_n\to Y_n$들의 귀납계로서 $u_m$이
$(u_n)_{(S_m)}$와 동일시되는 조건이 $m\leq n$일 때 성립하는 것이
대응함도 곧바로 확인된다.

이렇게 동치를 올바르게 정의했음은 다음의 더 정밀한 명제에서 따른다~:

\begin{proposition}[10.12.3.1]
\label{I.10.12.3.1-ko}
$(X_n)$을 $S_n$-준스킴들의 귀납계라 하자. 구조 사상
$f_n:X_n\to S_n$들이 도식
(\hyperref[I.10.12.2.1-ko]{10.12.2.1})을 가환이게 하고, $X_m$을
$X_n\times_{S_n}S_m$와 동일시하는 조건이 $m\leq n$일 때 성립한다고
가정하자.
그러면 귀납계 $(X_n)$은 (\hyperref[I.10.6.3-ko]{10.6.3})의 조건
b)와 c)를 만족한다. 그 귀납극한을 $\mathfrak{X}$라 하고,
$f:\mathfrak{X}\to\mathfrak{S}$를 귀납계 $(f_n)$의 귀납극한 사상이라
하자. 이때 $X_0$이 국소 뇌터이면
$\mathfrak{X}$도 국소 뇌터이고 $f$는 아딕 사상이다.
\end{proposition}

$\mathscr{O}_{S_n}$에서 부분준스킴 $S_m$을 $S_n$ 안에 정의하는
아이디얼층은 멱영이다. 따라서 (\hyperref[I.4.4.5-ko]{4.4.5})에
의해 $\mathscr{O}_{X_n}$에서 부분준스킴 $X_m$을 $X_n$ 안에 정의하는
아이디얼층도 멱영이므로, (\hyperref[I.10.6.3-ko]{10.6.3})의 조건들은
실제로 만족된다. 그러므로 문제는 $\mathfrak{X}$와 $\mathfrak{S}$에서
국소적이며, $\mathfrak{S}=\operatorname{Spf}(A)$,
$\mathscr{J}=\mathfrak{J}^\Delta$라 가정해도 된다. 여기서 $A$는
뇌터 $\mathfrak{J}$-아딕환이고, $X_n=\operatorname{Spec}(B_n)$이다.
$A_n=A/\mathfrak{J}^{n+1}$이라 하면,
가정에 의해 $B_0$은 뇌터이고,
$\mathfrak{J}_n=\mathfrak{J}/\mathfrak{J}^{n+1}$이라 놓을 때
$B_m=B_n/\mathfrak{J}_n^{m+1}B_n$이다. 따라서 $B_n\to B_0$의 핵은
$\mathfrak{K}_n=\mathfrak{J}_nB_n$이고, $B_n\to B_m$의 핵은
$\mathfrak{K}_n^{m+1}$이며, 이는 $m\leq n$일 때 성립한다. 또한 $A_1$이
뇌터이므로 $\mathfrak{J}_1$은 유한 생성 $A_1$-가군이다. 따라서
$\overline{\mathfrak{K}}_1=\mathfrak{K}_1/\mathfrak{K}_1^2$은
유한 생성 $B_1$-가군이고, \emph{따라서 특히}
$B_0=B_1/\mathfrak{K}_1$-가군으로도 유한 생성이다. 이제
$\mathfrak{X}$가 뇌터임은 (\hyperref[I.10.6.4-ko]{10.6.4})에서
따른다. $B=\varprojlim B_n$이라 하면
$\mathfrak{X}=\operatorname{Spf}(B)$이고, $\mathfrak{K}$를
$B\to B_0$의 핵이라 하면 $B_n=B/\mathfrak{K}^{n+1}$이다.
$\rho_n:A/\mathfrak{J}^{n+1}\to B/\mathfrak{K}^{n+1}$을 $f_n$에
대응하는 준동형이라 하면
\[
  \mathfrak{K}/\mathfrak{K}^{n+1}
  =(B/\mathfrak{K}^{n+1})\rho_n(\mathfrak{J}/\mathfrak{J}^{n+1})
\]
이다. 한편 준동형 $\rho:A\to B$가 $f$에 대응하며
$\varprojlim\rho_n$과 같으므로, 아이디얼 $\mathfrak{J}B$는 $B$의
아이디얼로서 $\mathfrak{K}$ 안에서 조밀하다. 그런데 $B$의 모든
아이디얼은 닫혀 있으므로 (\hyperref[0.7.3.5-ko]{0, 7.3.5}),
$\mathfrak{K}=\mathfrak{J}B$이다.
$\mathscr{K}=\mathfrak{K}^\Delta$라 하면 관계
$f^*(\mathscr{J})\mathscr{O}_{\mathfrak{X}}=\mathscr{K}$는
(\hyperref[I.10.10.9-ko]{10.10.9})에서 따르며, 이것으로 증명이 끝난다.

\begin{env}[10.12.3.2]
\label{I.10.12.3.2-ko}
앞의 동치는 두 $\mathfrak{S}$-아딕 형식적 준스킴
$\mathfrak{X}$, $\mathfrak{Y}$에 대하여 다음 \emph{표준 전단사}를 준다~:
\[
  \operatorname{Hom}_{\mathfrak{S}}(\mathfrak{X},\mathfrak{Y})
  \xrightarrow{\sim}
  \varprojlim_n\operatorname{Hom}_{S_n}(X_n,Y_n)
\]
여기서 사영극한은 사상 $u_n\to(u_n)_{(S_m)}$들에 관하여 취하며,
$m\leq n$이다.
\end{env}

\subsection{유한형 사상.}
\label{subsection:I.10.13-ko}

\begin{proposition}[10.13.1]
\label{I.10.13.1-ko}
$\mathfrak{Y}$를 국소 뇌터 형식적 준스킴,
$\mathscr{K}$를 $\mathfrak{Y}$의 정의 아이디얼,
$f:\mathfrak{X}\to\mathfrak{Y}$를 형식적 준스킴의 사상이라 하자. 다음
조건들은 동치이다~:
\begin{enumerate}
  \item[a)] $\mathfrak{X}$는 국소 뇌터이고, $f$는 아딕 사상이며
    (\hyperref[I.10.12.1-ko]{10.12.1}),
    $\mathscr{J}=f^*(\mathscr{K})\mathscr{O}_{\mathfrak{X}}$라 놓을 때
    사상 $f_0:(\mathfrak{X},\mathscr{O}_{\mathfrak{X}}/\mathscr{J})
    \to(\mathfrak{Y},\mathscr{O}_{\mathfrak{Y}}/\mathscr{K})$는 $f$로부터
    유도되며 유한형이다.
  \item[b)] $\mathfrak{X}$는 국소 뇌터이고, 아딕
    $(Y_n)$-귀납계 $(X_n)$의 귀납극한이며, 사상 $X_0\to Y_0$는
    유한형이다.

\oldpage[I]{208}

  \item[c)] $\mathfrak{Y}$의 모든 점은 다음 성질을 갖는 뇌터 아핀
    형식적 열린집합인 근방 $V$를 가진다~:
    \begin{enumerate}
      \item[(\textbf{Q})] $f^{-1}(V)$는 뇌터 아핀 형식적 열린집합
        $U_i$들의 유한 합집합이고, 각각에 대하여 뇌터 아딕환
        $\Gamma(U_i,\mathscr{O}_{\mathfrak{X}})$는
        $\Gamma(V,\mathscr{O}_{\mathfrak{Y}})$ 위의 제한 형식적 멱급수
        대수를 어떤 (필연적으로 닫힌) 아이디얼로 나눈 몫과 위상적으로
        동형이다 (\hyperref[0.7.5.1-ko]{0, 7.5.1}).
    \end{enumerate}
\end{enumerate}
\end{proposition}

a)가 b)를 함의하는 것은 (\hyperref[I.10.12.3-ko]{10.12.3})에 따라
명백하다. b)가 c)를 함의함을 보이기 위해, 문제는 $\mathfrak{Y}$ 위에서
국소적이므로 $\mathfrak{Y}=\operatorname{Spf}(B)$이고 $B$가 뇌터
아딕환이라고 가정할 수 있다. $\mathscr{K}=\mathfrak{K}^\Delta$라
하되, $\mathfrak{K}$는 $B$의 정의 아이디얼이라 하자. 가정에 의해 $X_0$는
$Y_0$ 위에서 유한형이므로, $X_0$는 아핀 열린집합 $U_i$들의 유한
합집합이고, 환 $A_{i0}$, 즉 $X_0$가 $U_i$ 위에 유도하는 아핀
스킴의 환은 $B/\mathfrak{K}$, 즉 $Y_0$의 환 위의 유한 생성 대수이다
(\hyperref[I.6.3.2-ko]{6.3.2}).
(\hyperref[I.5.1.9-ko]{5.1.9})에 따라 $U_i$는 각 뇌터 준스킴
$X_n$에서도 아핀 열린집합이며, 환 $A_{in}$을 $X_n$이 $U_i$ 위에
유도하는 아핀 스킴의 환이라 하면 가정 b)에 의해 $m\leq n$일 때
$A_{im}$은 $A_{in}/\mathfrak{K}^{m+1}A_{in}$과 동형이다. 따라서
$U_i$ 위에서 $\mathfrak{X}$가 유도하는 형식적 준스킴은
$\operatorname{Spf}(A_i)$와 동형이며, 여기서
$A_i=\varprojlim_n A_{in}$이다 (\hyperref[I.10.6.4-ko]{10.6.4}).
$A_i$는 $\mathfrak{K}A_i$-아딕환이고, $A_i/\mathfrak{K}A_i$는
$A_{i0}$과 동형이며 $B/\mathfrak{K}$ 위의 유한 생성 대수이다.
그러므로 (\hyperref[0.7.5.5-ko]{0, 7.5.5})에 따라 $A_i$는 $B$ 위의
제한 형식적 멱급수 대수를 어떤 아이디얼로 나눈 몫과 위상적으로
동형이다. 이 아이디얼은 반드시 닫혀 있는데, 그러한 대수는 뇌터이기 때문이다
(\hyperref[0.7.5.4-ko]{0, 7.5.4}).

c)가 a)를 함의함을 증명하기 위해, $\mathfrak{X}=\operatorname{Spf}(A)$
도 아핀인 경우로 한정할 수 있다. 여기서 $A$는 뇌터 아딕환이고, $B$
위의 제한 형식적 멱급수 대수를 어떤 닫힌 아이디얼로 나눈 몫과
동형이다. 그러면 (\hyperref[0.7.5.5-ko]{0, 7.5.5})에 따라
$A/\mathfrak{K}A$는 $B/\mathfrak{K}$ 위의 유한 생성 대수이고,
$\mathfrak{K}A=\mathfrak{J}$는 $A$의 정의 아이디얼이다. 따라서
(\hyperref[I.10.10.9-ko]{10.10.9})에 의해 a)의 조건들이 성립한다.

명제 (\hyperref[I.10.13.1-ko]{10.13.1})의 조건들이 성립하면, a)는
\emph{모든 정의 아이디얼 $\mathscr{K}$}에 대하여 성립한다. 여기서
정의 아이디얼은 $\mathfrak{Y}$의 것이며, 이는 c)에 따른다. 따라서
b)에서는 \emph{모든 $f_n$이 유한형 사상}임에 유의하자.

\begin{corollary}[10.13.2]
\label{I.10.13.2-ko}
(\hyperref[I.10.13.1-ko]{10.13.1})의 조건들이 성립하면, 모든 뇌터
아핀 형식적 열린집합 $V$는 $\mathfrak{Y}$ 안에서 성질
(\textbf{Q})를 가지며, $\mathfrak{Y}$가 뇌터이면 $\mathfrak{X}$도
뇌터이다.
\end{corollary}

이는 (\hyperref[I.10.13.1-ko]{10.13.1})과
(\hyperref[I.6.3.2-ko]{6.3.2})에서 곧바로 따른다.

\begin{definition}[10.13.3]
\label{I.10.13.3-ko}
(\hyperref[I.10.13.1-ko]{10.13.1})의 동치인 성질 a), b), c)가
성립하면, 사상 $f$를 \emph{유한형}이라고 하거나 $\mathfrak{X}$를
\emph{유한형 $\mathfrak{Y}$-형식적 준스킴} 또는
\emph{$\mathfrak{Y}$ 위의 유한형 형식적 준스킴}이라고 한다.
\end{definition}

\begin{corollary}[10.13.4]
\label{I.10.13.4-ko}
$\mathfrak{X}=\operatorname{Spf}(A)$,
$\mathfrak{Y}=\operatorname{Spf}(B)$를 두 뇌터 아핀 형식적 스킴이라
하자. $\mathfrak{X}$가 $\mathfrak{Y}$ 위에서 유한형이기 위한
필요충분조건은 뇌터 아딕환 $A$가 $B$ 위의 제한 형식적 멱급수 대수를
어떤 닫힌 아이디얼로 나눈 몫과 동형인 것이다.
\end{corollary}

실제로 (\hyperref[I.10.13.1-ko]{10.13.1})의 표기를 사용하자.
$\mathfrak{X}$가 $\mathfrak{Y}$ 위에서 유한형이면
(\hyperref[I.6.3.3-ko]{6.3.3})에 따라 $A/\mathfrak{K}A$는
$(B/\mathfrak{K})$ 위의 유한 생성 대수이고,
$\mathfrak{K}A$는 $A$의 정의 아이디얼이다
(\hyperref[I.10.10.9-ko]{10.10.9}). 그러므로
(\hyperref[0.7.5.5-ko]{0, 7.5.5})에서 결론이 따른다.

\begin{proposition}[10.13.5]
\label{I.10.13.5-ko}
\medskip\noindent
\begin{enumerate}
  \item[\rm{(i)}] 형식적 준스킴의 두 유한형 사상의 합성은
    유한형이다.

\oldpage[I]{209}

  \item[\rm{(ii)}] $\mathfrak{X}$, $\mathfrak{S}$,
    $\mathfrak{S}'$을 세 국소 뇌터(각각 뇌터) 형식적 준스킴이라 하고,
    $f:\mathfrak{X}\to\mathfrak{S}$,
    $g:\mathfrak{S}'\to\mathfrak{S}$를 두 사상이라 하자. $f$가
    유한형이면 $\mathfrak{X}\times_{\mathfrak{S}}\mathfrak{S}'$은
    국소 뇌터(각각 뇌터)이고 $\mathfrak{S}'$ 위에서 유한형이다.
  \item[\rm{(iii)}] $\mathfrak{S}$를 국소 뇌터 형식적 준스킴,
    $\mathfrak{X}'$, $\mathfrak{Y}'$을 두 국소 뇌터
    $\mathfrak{S}$-형식적 준스킴이라 하되,
    $\mathfrak{X}'\times_{\mathfrak{S}}\mathfrak{Y}'$은 국소 뇌터라
    하자.
    $\mathfrak{X}$, $\mathfrak{Y}$가 국소 뇌터
    $\mathfrak{S}$-형식적 준스킴이고,
    $f:\mathfrak{X}\to\mathfrak{X}'$,
    $g:\mathfrak{Y}\to\mathfrak{Y}'$이 두 유한형
    $\mathfrak{S}$-사상이면,
    $\mathfrak{X}\times_{\mathfrak{S}}\mathfrak{Y}$은 국소 뇌터이고
    $f\times_{\mathfrak{S}}g$는 유한형 $\mathfrak{S}$-사상이다.
\end{enumerate}
\end{proposition}

(iii)은 (\hyperref[I.3.5.1-ko]{3.5.1})의 형식적 논증에 의해 (i)과
(ii)에서 따르므로, (i)과 (ii)만 증명하면 충분하다.

$\mathfrak{X}$, $\mathfrak{Y}$, $\mathfrak{Z}$를 세 국소 뇌터 형식적
준스킴이라 하고, $f:\mathfrak{X}\to\mathfrak{Y}$,
$g:\mathfrak{Y}\to\mathfrak{Z}$를 두 유한형 사상이라 하자.
$\mathscr{L}$이 $\mathfrak{Z}$의 정의 아이디얼이면,
$\mathscr{K}=g^*(\mathscr{L})\mathscr{O}_{\mathfrak{Y}}$는
$\mathfrak{Y}$의 정의 아이디얼이고,
$\mathscr{J}=f^*(g^*(\mathscr{L}))\mathscr{O}_{\mathfrak{X}}$는
$\mathfrak{X}$의 정의 아이디얼이다. 다음과 같이 놓자~:
$X_0=(\mathfrak{X},\mathscr{O}_{\mathfrak{X}}/\mathscr{J})$,
$Y_0=(\mathfrak{Y},\mathscr{O}_{\mathfrak{Y}}/\mathscr{K})$,
$Z_0=(\mathfrak{Z},\mathscr{O}_{\mathfrak{Z}}/\mathscr{L})$.
또 $f_0:X_0\to Y_0$, $g_0:Y_0\to Z_0$를 각각 $f$, $g$에 대응하는
사상이라 하자. 가정에 의해 $f_0$와 $g_0$가 유한형이므로,
$g_0\circ f_0$도 유한형이다 (\hyperref[I.6.3.4-ko]{6.3.4}). 이
사상은 $g\circ f$에 대응하므로, (\hyperref[I.10.13.1-ko]{10.13.1})에
따라 $g\circ f$는 유한형이다.

(ii)의 조건 아래에서 $\mathfrak{S}$ (각각 $\mathfrak{X}$,
$\mathfrak{S}'$)는 국소 뇌터 준스킴들의 열 $(S_n)$ (각각 $(X_n)$,
$(S'_n)$)의 귀납극한이고, (\hyperref[I.10.13.1-ko]{10.13.1})에 따라
$X_m=X_n\times_{S_n}S_m$라고 가정할 수 있으며, 이는 모든
$m\leq n$에 대해 성립한다. 그러면
형식적 준스킴
$\mathfrak{X}\times_{\mathfrak{S}}\mathfrak{S}'$은 준스킴
$X_n\times_{S_n}S'_n$들의 귀납극한이고
(\hyperref[I.10.7.4-ko]{10.7.4}), 다음이 성립한다~:
\[
  X_m\times_{S_m}S'_m
  =(X_n\times_{S_n}S_m)\times_{S_m}S'_m
  =(X_n\times_{S_n}S'_n)\times_{S'_n}S'_m.
\]
또한 $X_0\times_{S_0}S'_0$은 $X_0$가 $S_0$ 위에서 유한형이므로
국소 뇌터이다
(\hyperref[I.6.3.8-ko]{6.3.8}). 따라서 먼저
(\hyperref[I.10.12.3.1-ko]{10.12.3.1})에 의해
$\mathfrak{X}\times_{\mathfrak{S}}\mathfrak{S}'$은 국소 뇌터이다.
더 나아가 $X_0\times_{S_0}S'_0$은 $S'_0$ 위에서 유한형이므로
(\hyperref[I.6.3.8-ko]{6.3.8}),
(\hyperref[I.10.12.3.1-ko]{10.12.3.1})과
(\hyperref[I.10.13.1-ko]{10.13.1})에 따라
$\mathfrak{X}\times_{\mathfrak{S}}\mathfrak{S}'$은
$\mathfrak{S}'$ 위에서 유한형이다. 이것으로 (ii)의 증명이 끝난다.
뇌터 준스킴에 관한 주장은 (\hyperref[I.6.3.8-ko]{6.3.8})의 즉각적인
귀결이다.

\begin{corollary}[10.13.6]
\label{I.10.13.6-ko}
(\hyperref[I.10.9.9-ko]{10.9.9})의 가정 아래에서 $f$가 유한형
사상이면, 완비화한 것들 사이로의 연장 $\widehat f$도 유한형이다.
\end{corollary}

\subsection{형식적 준스킴의 닫힌 부분준스킴.}
\label{subsection:I.10.14-ko}

\begin{proposition}[10.14.1]
\label{I.10.14.1-ko}
$\mathfrak{X}$를 국소 뇌터 형식적 준스킴,
$\mathscr{A}$를 $\mathscr{O}_{\mathfrak{X}}$의 연접 아이디얼층이라
하자. $\mathfrak{Y}$가 $\mathscr{O}_{\mathfrak{X}}/\mathscr{A}$의
(닫힌) 지지집합이면, 위상환 달린 공간
$(\mathfrak{Y},(\mathscr{O}_{\mathfrak{X}}/\mathscr{A})|\mathfrak{Y})$은
국소 뇌터 형식적 준스킴이며, $\mathfrak{X}$가 뇌터이면 이 형식적
준스킴도 뇌터이다.
\end{proposition}

(\hyperref[I.10.10.3-ko]{10.10.3})과
(\hyperref[0.5.3.4-ko]{0, 5.3.4})에 따라
$\mathscr{O}_{\mathfrak{X}}/\mathscr{A}$는 구조층 위의 연접 가군이고, 따라서 그
지지집합 $\mathfrak{Y}$는 닫혀 있다
(\hyperref[0.5.2.2-ko]{0, 5.2.2}). $\mathscr{J}$를
$\mathfrak{X}$의 정의 아이디얼이라 하고,
$X_n=(\mathfrak{X},\mathscr{O}_{\mathfrak{X}}/\mathscr{J}^{n+1})$이라
하자~; 환의 층 $\mathscr{O}_{\mathfrak{X}}/\mathscr{A}$는 층들
$\mathscr{O}_{\mathfrak{X}}/(\mathscr{A}+\mathscr{J}^{n+1})
=(\mathscr{O}_{\mathfrak{X}}/\mathscr{A})
\otimes_{\mathscr{O}_{\mathfrak{X}}}
(\mathscr{O}_{\mathfrak{X}}/\mathscr{J}^{n+1})$의 사영극한이다
(\hyperref[I.10.11.3-ko]{10.11.3}). 이 층들은 모두 지지집합
$\mathfrak{Y}$를 갖는다. 층
$(\mathscr{A}+\mathscr{J}^{n+1})/\mathscr{J}^{n+1}$은
$\mathscr{O}_{\mathfrak{X}}$-가군으로서 연접이다. 실제로
$\mathscr{J}^{n+1}$이 연접이므로,
$(\mathscr{A}+\mathscr{J}^{n+1})/\mathscr{J}^{n+1}$은 또한 연접
$(\mathscr{O}_{\mathfrak{X}}/\mathscr{J}^{n+1})$-가군이다
(\hyperref[0.5.3.10-ko]{0, 5.3.10})~; $Y_n$을 이 아이디얼층으로
정의되는 $X_n$의 닫힌 부분준스킴이라 하면, 위상환 달린 공간
$(\mathfrak{Y},(\mathscr{O}_{\mathfrak{X}}/\mathscr{A})|\mathfrak{Y})$

\oldpage[I]{210}
은 $Y_n$들의 귀납극한인 형식적 준스킴임이 곧바로 확인된다.
(\hyperref[I.10.6.4-ko]{10.6.4})의 조건들이 만족되므로, 이 형식적
준스킴은 국소 뇌터이고, $\mathfrak{X}$가 뇌터이면 뇌터이다. 실제로
이 경우에는 (\hyperref[I.6.1.4-ko]{6.1.4})에 따라 $Y_0$가 뇌터이다.

\begin{definition}[10.14.2]
\label{I.10.14.2-ko}
형식적 준스킴 $\mathfrak{X}$의 \emph{닫힌 부분준스킴}이라 함은
$(\mathfrak{Y},(\mathscr{O}_{\mathfrak{X}}/\mathscr{A})|\mathfrak{Y})$
꼴의 모든 형식적 준스킴을 말한다. 여기서 $\mathscr{A}$는
$\mathscr{O}_{\mathfrak{X}}$의 연접 아이디얼층이다~; 이 형식적
준스킴을 $\mathscr{A}$가 정의하는 닫힌 부분준스킴이라고 한다.
\footnote{\textbf{번역 주.} 인쇄 원문은 여기와 바로 다음 문장에서
단지 연접 가군이라고 쓰지만, 몫환의 구성과 뒤따르는 대응에는 연접
아이디얼층이 필요하므로 두 곳 모두 이를 교정했다.}
\end{definition}

이와 같이 정의된 $\mathscr{O}_{\mathfrak{X}}$의 연접 아이디얼층들과
$\mathfrak{X}$의 닫힌 부분준스킴들 사이의 대응은 전단사임이 명백하다.

위상환 달린 공간의 사상
$j=(\psi,\theta):\mathfrak{Y}\to\mathfrak{X}$에서 $\psi$는 단사사상
$\mathfrak{Y}\to\mathfrak{X}$이고 $\theta^\sharp$는 표준 준동형
$\mathscr{O}_{\mathfrak{X}}\to\mathscr{O}_{\mathfrak{X}}/\mathscr{A}$이다.
이는 명백히 (\hyperref[I.10.4.5-ko]{10.4.5}) 형식적 준스킴의
사상이며, 이를 $\mathfrak{Y}$에서 $\mathfrak{X}$로 가는
\emph{표준 포함 사상}이라 한다. $\mathfrak{X}=\operatorname{Spf}(A)$이고
$A$가 뇌터 아딕환이면, $\mathscr{A}=\mathfrak{a}^\Delta$이고 여기서
$\mathfrak{a}$는 $A$의 아이디얼이다
(\hyperref[I.10.10.5-ko]{10.10.5}). 앞의 내용에서 곧바로
$\mathfrak{Y}=\operatorname{Spf}(A/\mathfrak{a})$임이 동형을 제외하면
따르고, $j$는 (\hyperref[I.10.2.2-ko]{10.2.2}) 표준 준동형
$A\to A/\mathfrak{a}$에 대응한다.

국소 뇌터 형식적 준스킴의 사상
$f:\mathfrak{Z}\to\mathfrak{X}$가 \emph{닫힌 몰입 사상}이라 함은
이 사상이 $\mathfrak{Z}\xrightarrow{g}\mathfrak{Y}\xrightarrow{j}\mathfrak{X}$로
인수분해되고, 여기서 $g$는 $\mathfrak{Z}$에서 닫힌 부분준스킴
$\mathfrak{Y}$로 가는 동형사상이며, 이 부분준스킴은 $\mathfrak{X}$의
닫힌 부분준스킴이고 $j$는 표준 포함
사상인 것을 뜻한다. $j$는 환 달린 공간의 단사사상이므로 $g$와
$\mathfrak{Y}$는 반드시 \emph{유일}하다.

\begin{proposition}[10.14.3]
\label{I.10.14.3-ko}
닫힌 몰입 사상은 유한형 사상이다.
\end{proposition}

$\mathfrak{X}$가 아핀 형식적 스킴 $\operatorname{Spf}(A)$이고
$\mathfrak{Y}=\operatorname{Spf}(A/\mathfrak{a})$인 경우로 곧바로
환원된다~; 명제는 (\hyperref[I.10.13.1-ko]{10.13.1, c)})에서 따른다.

\begin{lemma}[10.14.4]
\label{I.10.14.4-ko}
$f:\mathfrak{Y}\to\mathfrak{X}$를 국소 뇌터 형식적 준스킴의
사상이라 하고, $(U_\alpha)$를 $f(\mathfrak{Y})$의 덮개라 하자. 이
덮개의 원소들은 $\mathfrak{X}$의 뇌터 아핀 형식적 열린집합이고, 각
$f^{-1}(U_\alpha)$는
$\mathfrak{Y}$의 뇌터 아핀 형식적 열린집합이라고 하자. $f$가 닫힌
몰입 사상이기 위한 필요충분조건은 $f(\mathfrak{Y})$가
$\mathfrak{X}$의 닫힌 부분집합이고, 모든 $\alpha$에 대하여 $f$를
$f^{-1}(U_\alpha)$에 제한한 사상이
(\hyperref[I.10.4.6-ko]{10.4.6})에 따라 전사 준동형
$\Gamma(U_\alpha,\mathscr{O}_{\mathfrak{X}})\to
\Gamma(f^{-1}(U_\alpha),\mathscr{O}_{\mathfrak{Y}})$에 대응하는 것이다.
\end{lemma}

이 조건들이 필요함은 명백하다. 역으로 이 조건들이 성립한다고 하고,
$\mathfrak{a}_\alpha$를 준동형
$\Gamma(U_\alpha,\mathscr{O}_{\mathfrak{X}})\to
\Gamma(f^{-1}(U_\alpha),\mathscr{O}_{\mathfrak{Y}})$의 핵이라 하자.
연접 아이디얼층 $\mathscr{A}$를 $\mathscr{O}_{\mathfrak{X}}$ 위에
다음과 같이 정의한다. $\mathscr{A}|U_\alpha=\mathfrak{a}_\alpha^\Delta$로
두고, $\mathscr{A}$는 $U_\alpha$들의 합집합의 여집합에서 단위
아이디얼로 둔다. 실제로 $f(\mathfrak{Y})$가 닫혀 있고
$(\mathscr{O}_{\mathfrak{X}}|U_\alpha)/\mathfrak{a}_\alpha^\Delta$의
지지집합이 $U_\alpha\cap f(\mathfrak{Y})$이므로,
$\mathfrak{a}_\alpha^\Delta$와 $\mathfrak{a}_\beta^\Delta$가 뇌터
아핀 형식적 열린집합 $V\subset U_\alpha\cap U_\beta$ 위에서 같은
층을 유도함을 확인하면 충분하다.\footnote{\textbf{번역 주.} 인쇄
원문은 여집합에서 영 아이디얼로 둔다고 쓰지만, 닫힌 부분집합인
사상의 상 밖에서 몫층이 영층이 되려면 단위 아이디얼이어야 한다.
또한 그 닫힌 부분집합은 핵 아이디얼층 자체의 지지집합이 아니라 그
몫층의 지지집합이므로 두 곳을 함께 교정했다.}
$f$를 $f^{-1}(U_\alpha)$에 제한한 사상은 이 형식적 준스킴에서
$U_\alpha$로 가는 닫힌 몰입 사상이므로, $f^{-1}(V)$는
$f^{-1}(U_\alpha)$의 뇌터 아핀 형식적 열린집합이고 $f$를
$f^{-1}(V)$에 제한한 사상도 닫힌 몰입 사상이다~; 이 제한에 대응하는
전사 준동형의 핵을 $\mathfrak{b}$라 하자. 그 준동형은
$\Gamma(V,\mathscr{O}_{\mathfrak{X}})\to
\Gamma(f^{-1}(V),\mathscr{O}_{\mathfrak{Y}})$이다. 그러면
(\hyperref[I.10.10.2-ko]{10.10.2})에 따라
$\mathfrak{a}_\alpha^\Delta$가 $\mathfrak{b}^\Delta$를 $V$ 위에
유도함은 명백하다. 이렇게 아이디얼층 $\mathscr{A}$를 정의하면
$f=j\circ g$임이 명백하다. 여기서
$j:\mathfrak{Z}\to\mathfrak{X}$는 닫힌 부분준스킴 $\mathfrak{Z}$에서
$\mathfrak{X}$로 가는 표준 포함 사상이고, 이 부분준스킴은
$\mathscr{A}$가 정의하며,
$g$는 $\mathfrak{Y}$에서 $\mathfrak{Z}$로 가는 동형사상이다.
\footnote{\textbf{번역 주.} 인쇄 원문은 합성의 순서를 거꾸로 쓰지만,
선언된 정의역과 공역에 맞추어 바로잡았다.}

\begin{proposition}[10.14.5]
\label{I.10.14.5-ko}
\begin{enumerate}
  \item[\rm{(i)}] $f:\mathfrak{Z}\to\mathfrak{Y}$,
  $g:\mathfrak{Y}\to\mathfrak{X}$가 국소 뇌터 형식적 준스킴의
  닫힌 몰입 사상이면, $g\circ f$는 닫힌 몰입 사상이다.

\oldpage[I]{211}
  \item[\rm{(ii)}] $\mathfrak{X}$, $\mathfrak{Y}$,
  $\mathfrak{S}$를 세 국소 뇌터 형식적 준스킴이라 하고,
  $f:\mathfrak{X}\to\mathfrak{S}$를 닫힌 몰입 사상,
  $g:\mathfrak{Y}\to\mathfrak{S}$를 사상이라 하자. 그러면 사상
  $\mathfrak{X}\times_{\mathfrak{S}}\mathfrak{Y}\to\mathfrak{Y}$는
  닫힌 몰입 사상이다.
  \item[\rm{(iii)}] $\mathfrak{S}$를 국소 뇌터 형식적 준스킴,
  $\mathfrak{X}'$, $\mathfrak{Y}'$을 두 국소 뇌터
  $\mathfrak{S}$-형식적 준스킴이라 하되,
  $\mathfrak{X}'\times_{\mathfrak{S}}\mathfrak{Y}'$은 국소 뇌터라고
  하자. $\mathfrak{X}$, $\mathfrak{Y}$가 두 국소 뇌터
  $\mathfrak{S}$-형식적 준스킴이고,
  $f:\mathfrak{X}\to\mathfrak{X}'$, $g:\mathfrak{Y}\to\mathfrak{Y}'$이
  닫힌 몰입 사상인 두 $\mathfrak{S}$-사상이면,
  $f\times_{\mathfrak{S}}g$는 닫힌 몰입 사상이다.
\end{enumerate}
\end{proposition}

(\hyperref[I.3.5.1-ko]{3.5.1})에 따라 여기서도 (i)과 (ii)만 증명하면
충분하다.

(i)을 증명하기 위해 $\mathfrak{Y}$ (각각 $\mathfrak{Z}$)가
$\mathfrak{X}$ (각각 $\mathfrak{Y}$)의 닫힌 부분준스킴이고, 이들이
각각 연접 아이디얼층 $\mathscr{J}$ (각각 $\mathscr{K}$)로 정의된다고
가정할 수 있다. 이 두 층은 각각 $\mathscr{O}_{\mathfrak{X}}$ (각각
$\mathscr{O}_{\mathfrak{Y}}$)의 아이디얼층이다~; $\psi$가 바탕공간의 단사사상
$\mathfrak{Y}\to\mathfrak{X}$이면, $\psi_*(\mathscr{K})$는
$\psi_*(\mathscr{O}_{\mathfrak{Y}})
=\mathscr{O}_{\mathfrak{X}}/\mathscr{J}$의 연접 아이디얼층이다
(\hyperref[0.5.3.12-ko]{0, 5.3.12}). 따라서 이는 또한 연접
$\mathscr{O}_{\mathfrak{X}}$-가군이다
(\hyperref[0.5.3.10-ko]{0, 5.3.10})~; 핵 $\mathscr{K}_1$, 즉 준동형
$\mathscr{O}_{\mathfrak{X}}\to
(\mathscr{O}_{\mathfrak{X}}/\mathscr{J})/\psi_*(\mathscr{K})$의 핵은
$\mathscr{O}_{\mathfrak{X}}$의 연접 아이디얼층이고
(\hyperref[0.5.3.4-ko]{0, 5.3.4}),
$\mathscr{O}_{\mathfrak{X}}/\mathscr{K}_1$은
$\psi_*(\mathscr{O}_{\mathfrak{Y}}/\mathscr{K})$와 동형이다. 따라서
$\mathfrak{Z}$는 $\mathfrak{X}$의 닫힌 부분준스킴과 동형이다.

(ii)를 증명하기 위해 $\mathfrak{S}=\operatorname{Spf}(A)$,
$\mathfrak{X}=\operatorname{Spf}(B)$,
$\mathfrak{Y}=\operatorname{Spf}(C)$인 경우로 한정할 수 있음은
명백하다. 여기서 $A$는 뇌터 $\mathfrak{J}$-아딕환이고,
$B=A/\mathfrak{a}$이며 $\mathfrak{a}$는 $A$의 아이디얼이고,
$C$는 아딕이고 뇌터인 위상 $A$-대수이다. 이제 준동형
$C\to C\widehat{\otimes}_A(A/\mathfrak{a})$이 \emph{전사}임을
보이면 충분하다~: 실제로 $A/\mathfrak{a}$는 유한 생성 $A$-가군이고
그 위상은 $\mathfrak{J}$-아딕 위상이다~; 따라서
(\hyperref[0.7.7.8-ko]{0, 7.7.8})에 의해
$C\widehat{\otimes}_A(A/\mathfrak{a})$는
$C\otimes_A(A/\mathfrak{a})=C/\mathfrak{a}C$와 동일시되므로 주장이
따른다.

\begin{corollary}[10.14.6]
\label{I.10.14.6-ko}
(\hyperref[I.10.14.5-ko]{10.14.5, (ii)})의 가정 아래에서
$p:\mathfrak{X}\times_{\mathfrak{S}}\mathfrak{Y}\to\mathfrak{X}$,
$q:\mathfrak{X}\times_{\mathfrak{S}}\mathfrak{Y}\to\mathfrak{Y}$를
표준 사영들이라 하자. 그러면 도식
\[
  \xymatrix{
    \mathfrak{X} \ar[d]_{f}
    & \mathfrak{X}\times_{\mathfrak{S}}\mathfrak{Y}
      \ar[l]_{p} \ar[d]^{q}\\
    \mathfrak{S}
    & \mathfrak{Y} \ar[l]^{g}
  }
\]
은 가환이다. 모든 연접 $\mathscr{O}_{\mathfrak{X}}$-가군
$\mathscr{F}$에 대하여 다음 표준 $\mathscr{O}_{\mathfrak{Y}}$-가군
동형사상이 존재한다~:
\[
  \label{I.10.14.6.1-ko}
  u:g^*(f_*(\mathscr{F}))
  \xrightarrow{\sim}
  q_*(p^*(\mathscr{F})).
  \tag{10.14.6.1}
\]
\end{corollary}

준동형 $g^*(f_*(\mathscr{F}))\to q_*(p^*(\mathscr{F}))$을 정의하는
것은 준동형
$f_*(\mathscr{F})\to g_*(q_*(p^*(\mathscr{F})))
=f_*(p_*(p^*(\mathscr{F})))$을 정의하는 것과 같다
(\hyperref[0.4.4.3-ko]{0, 4.4.3})~: 앞의 수반 대응 아래에서 원문의
표기를 따라 $u=f_*(\rho)$로 쓰겠다. 여기서 $\rho$는 표준 준동형
$\mathscr{F}\to p_*(p^*(\mathscr{F}))$이다
(\hyperref[0.4.4.3-ko]{0, 4.4.3}).\footnote{\textbf{번역 주.} 이
등식은 서로 다른 준동형 집합의 원소들을 수반 전단사로 대응시켜 쓴
표기이며, 두 사상이 문자 그대로 같은 형을 갖는다는 뜻이 아니다.}
$u$가 동형사상임을 보이기 위해 $\mathfrak{S}$, $\mathfrak{X}$,
$\mathfrak{Y}$가 각각 뇌터 아딕환 $A$, $B$, $C$의 형식적
스펙트럼이고, (\hyperref[I.10.14.5-ko]{10.14.5, (ii)})에서 본
조건들이 성립하는 경우로 곧바로 환원된다~; 이때
$\mathscr{F}=M^\Delta$이고, $M$은 유한 생성
$(A/\mathfrak{a})$-가군이다
(\hyperref[I.10.10.5-ko]{10.10.5}).
(\hyperref[I.10.14.6.1-ko]{10.14.6.1})의 양변은
(\hyperref[I.10.10.8-ko]{10.10.8})에 따라 각각
$(C\otimes_A M)^\Delta$와
$((C/\mathfrak{a}C)\otimes_{A/\mathfrak{a}}M)^\Delta$로 동일시된다.
그런데
$(C/\mathfrak{a}C)\otimes_{A/\mathfrak{a}}M
=(C\otimes_A(A/\mathfrak{a}))\otimes_{A/\mathfrak{a}}M$은
$C\otimes_A M$과 표준적으로 동일시되므로 따름정리가 증명된다.

\oldpage[I]{212}
\begin{corollary}[10.14.7]
\label{I.10.14.7-ko}
$X$를 (보통 의미의) 국소 뇌터 준스킴, $Y$를 $X$의 닫힌 부분준스킴,
$j$를 표준 포함 사상 $Y\to X$, $X'$을 $X$의 닫힌 부분집합이라 하고
$Y'=Y\cap X'$이라 하자~; 그러면
$\widehat{j}:Y_{/Y'}\to X_{/X'}$은 닫힌 몰입 사상이고, 모든 연접
$\mathscr{O}_Y$-가군 $\mathscr{F}$에 대하여
\[
  \widehat{j}_*(\mathscr{F}_{/Y'})=(j_*(\mathscr{F}))_{/X'}.
\]
\end{corollary}

$Y'=j^{-1}(X')$이므로 (\hyperref[I.10.9.9-ko]{10.9.9})를 사용하고
(\hyperref[I.10.14.5-ko]{10.14.5})와
(\hyperref[I.10.14.6-ko]{10.14.6})을 적용하면 충분하다.
\subsection{분리된 형식적 준스킴.}
\label{subsection:I.10.15-ko}

\begin{definition}[10.15.1]
\label{I.10.15.1-ko}
형식적 준스킴 $\mathfrak{S}$, $\mathfrak{S}$-형식적 준스킴
$\mathfrak{X}$, 구조 사상 $f:\mathfrak{X}\to\mathfrak{S}$가 주어졌다고
하자. 대각 사상
$\Delta_{\mathfrak{X}|\mathfrak{S}}:\mathfrak{X}\to
\mathfrak{X}\times_{\mathfrak{S}}\mathfrak{X}$ (또한
$\Delta_{\mathfrak{X}}$라고도 쓴다)을 사상
$(1_{\mathfrak{X}},1_{\mathfrak{X}})_{\mathfrak{S}}$라 한다. 대각 사상
$\Delta_{\mathfrak{X}}$에 의한 바탕공간 $\mathfrak{X}$의 상이
$\mathfrak{X}\times_{\mathfrak{S}}\mathfrak{X}$의 바탕공간의 닫힌 부분집합이면,
$\mathfrak{X}$가 $\mathfrak{S}$ 위에서 분리되어 있다고, 또는
$\mathfrak{S}$-형식적 스킴이라고, 또는 $f$가 분리 사상이라고 한다. 형식적
준스킴 $\mathfrak{X}$가 분리되어 있다고, 또는 형식적 스킴이라고 하는 것은
$\mathbf{Z}$ 위에서 분리되어 있다는 뜻이다.
\end{definition}

\begin{proposition}[10.15.2]
\label{I.10.15.2-ko}
형식적 준스킴 $\mathfrak{S}$, $\mathfrak{X}$가 각각 보통 준스킴들의 열
$(S_n)$, $(X_n)$의 귀납극한이고, 사상
$f:\mathfrak{X}\to\mathfrak{S}$가 사상들의 열
$f_n:X_n\to S_n$의 귀납극한이라고 하자. $f$가 분리 사상이기 위한
필요충분조건은 $f_0:X_0\to S_0$가 분리 사상인 것이다.
\end{proposition}

실제로, $\Delta_{\mathfrak{X}|\mathfrak{S}}$는 사상들의 열
$\Delta_{X_n|S_n}$의 귀납극한이다
(\hyperref[I.10.7.4-ko]{10.7.4}). 또한 바탕공간 $\mathfrak{X}$의
$\Delta_{\mathfrak{X}|\mathfrak{S}}$에 의한 상(각각 바탕공간
$\mathfrak{X}\times_{\mathfrak{S}}\mathfrak{X}$)은 바탕공간 $X_0$의
$\Delta_{X_0|S_0}$에 의한 상(각각 바탕공간
$X_0\times_{S_0}X_0$)과 동일하다~; 따라서 결론이 따른다.

\begin{proposition}[10.15.3]
\label{I.10.15.3-ko}
이하에서는 고려하는 모든 형식적 준스킴(각각 형식적 준스킴의 모든 사상)이
보통 준스킴들의 열(각각 보통 준스킴의 사상들의 열)의 귀납극한이라고
가정한다.
\begin{enumerate}
  \item[\rm{(i)}] 두 분리 사상의 합성은 분리 사상이다.
  \item[\rm{(ii)}] $f:\mathfrak{X}\to\mathfrak{X}'$,
  $g:\mathfrak{Y}\to\mathfrak{Y}'$가 두 분리
  $\mathfrak{S}$-사상이면 $f\times_{\mathfrak{S}}g$도 분리 사상이다.
  \item[\rm{(iii)}] $f:\mathfrak{X}\to\mathfrak{Y}$가 분리
  $\mathfrak{S}$-사상이면, 밑 형식적 준스킴의 임의의 확장
  $\mathfrak{S}'\to\mathfrak{S}$에 대하여 $\mathfrak{S}'$-사상
  $f_{(\mathfrak{S}')}$도 분리 사상이다.
  \item[\rm{(iv)}] 두 사상의 합성 $g\circ f$가 분리 사상이면 $f$도 분리
  사상이다.
\end{enumerate}
\emph{(이 명제에서 동일한 형식적 준스킴 $\mathfrak{Z}$가 하나의 명제 안에
여러 번 나타나는 경우에는, 나타나는 모든 곳에서 그것을 보통 준스킴들의
동일한 열 $(Z_n)$의 귀납극한으로 간주한다는 것을 함의한다. 또한
형식적 준스킴 $\mathfrak{Z}$에서 형식적 준스킴으로 가는 사상(각각 형식적
준스킴에서 $\mathfrak{Z}$로 가는 사상)은 보통 준스킴의 사상들
 $Z_n$에서 보통 준스킴으로 가는 사상(각각 보통 준스킴에서 $Z_n$으로
가는 사상)의 귀납극한으로 간주한다.)}
\end{proposition}

실제로 (\hyperref[I.10.15.2-ko]{10.15.2})의 표기를 사용하면
$(g\circ f)_0=g_0\circ f_0$이고,
$(f\times_{\mathfrak{S}}g)_0=f_0\times_{S_0}g_0$이다~; 따라서
(\hyperref[I.10.15.3-ko]{10.15.3})의 주장들은 보통 준스킴에 대한
(\hyperref[I.10.15.2-ko]{10.15.2})와
(\hyperref[I.5.5.1-ko]{5.5.1})의 대응하는 주장들의 즉각적인
귀결이다.

독자에게는 (\hyperref[I.10.15.3-ko]{10.15.3})과 같은 종류의 형식적
준스킴 및 사상에 대하여
(\hyperref[I.5.5.5-ko]{5.5.5}),

\oldpage[I]{213}
(\hyperref[I.5.5.9-ko]{5.5.9}) 및
(\hyperref[I.5.5.10-ko]{5.5.10})에 대응하는 명제들을 같은 방식으로
서술하는 일을 맡긴다(여기서 «~아핀 열린집합~»을
«~(\hyperref[I.10.6.3-ko]{10.6.3})의 조건 b)를 만족하는 아핀 형식적
열린집합~»으로 바꾼다).

유사한 논증으로 모든 뇌터 아핀 형식적 스킴이 분리되어 있음도 보인다.
따라서 이 용어가 정당화된다.

\begin{proposition}[10.15.4]
\label{I.10.15.4-ko}
형식적 준스킴 $\mathfrak{S}$를 국소 뇌터라고 하고,
$\mathfrak{X}$, $\mathfrak{Y}$를 두 국소 뇌터
$\mathfrak{S}$-형식적 준스킴이라 하자. $\mathfrak{X}$ 또는
$\mathfrak{Y}$가 $\mathfrak{S}$ 위에서 유한형이라고 가정하자(따라서
$\mathfrak{X}\times_{\mathfrak{S}}\mathfrak{Y}$는 국소 뇌터이다).
또 $\mathfrak{Y}$가 $\mathfrak{S}$ 위에서 분리되어 있다고 하자. 그러면
$\mathfrak{S}$-사상 $f:\mathfrak{X}\to\mathfrak{Y}$의 그래프 사상
$\Gamma_f=(1_{\mathfrak{X}},f)_{\mathfrak{S}}:\mathfrak{X}\to
\mathfrak{X}\times_{\mathfrak{S}}\mathfrak{Y}$는 닫힌 몰입 사상이다.
\end{proposition}

형식적 준스킴 $\mathfrak{S}$가 국소 뇌터 준스킴들의 열
$(S_n)$의 귀납극한이고, $\mathfrak{X}$(각각 $\mathfrak{Y}$)가 열
$(X_n)$(각각 $(Y_n)$)의 $S_n$-준스킴들의 귀납극한이며, $f$가 $S_n$-사상들의 열
$f_n:X_n\to Y_n$의 귀납극한이라고 가정할 수 있다. 그러면
$\mathfrak{X}\times_{\mathfrak{S}}\mathfrak{Y}$는 열
$(X_n\times_{S_n}Y_n)$의 귀납극한이고, $\Gamma_f$는 열
$(\Gamma_{f_n})$의 귀납극한이다
(\hyperref[I.10.7.4-ko]{10.7.4}). 가정에 따라 $Y_0$는 $S_0$ 위에서
분리되어 있다(\hyperref[I.10.15.2-ko]{10.15.2}). 따라서 공간
$\Gamma_{f_0}(X_0)$는 $X_0\times_{S_0}Y_0$의 닫힌 부분공간이다.
$\mathfrak{X}\times_{\mathfrak{S}}\mathfrak{Y}$의 바탕공간(각각
$\Gamma_f(\mathfrak{X})$)과 $X_0\times_{S_0}Y_0$(각각
$\Gamma_{f_0}(X_0)$)의 바탕공간이 서로 같으므로, 이미
$\Gamma_f(\mathfrak{X})$가
$\mathfrak{X}\times_{\mathfrak{S}}\mathfrak{Y}$의 닫힌
부분공간임을 알 수 있다. 이제 $(U,V)$가 $\mathfrak{X}$의 뇌터 아핀
형식적 열린집합 $U$(각각 $\mathfrak{Y}$의 뇌터 아핀 형식적 열린집합
$V$)로 이루어진 순서쌍들 중 $f(U)\subset V$를 만족하는 것들을
취하면, 열린집합 $U\times_{\mathfrak{S}}V$들은
$\mathfrak{X}\times_{\mathfrak{S}}\mathfrak{Y}$ 안에서
$\Gamma_f(\mathfrak{X})$의 덮개를 이룬다. $f_U:U\to V$를 $f$의
$U$로의 제한이라 하면, $\Gamma_{f_U}:U\to
U\times_{\mathfrak{S}}V$는 $U$ 위에서 $\Gamma_f$를 제한한 것이다.
따라서 $\Gamma_{f_U}$가 닫힌 몰입 사상임을 보이면
(\hyperref[I.10.14.4-ko]{10.14.4})에 따라 $\Gamma_f$도 닫힌 몰입
사상이 되며, 문제는 다음 경우로 환원된다:
$\mathfrak{S}=\operatorname{Spf}(A)$,
$\mathfrak{X}=\operatorname{Spf}(B)$,
$\mathfrak{Y}=\operatorname{Spf}(C)$가 아핀이고
($A$, $B$, $C$는 뇌터 아딕환), $f$가 연속 $A$-준동형
$\varphi:C\to B$에 대응하는 경우이다. 이때 $\Gamma_f$는 유일한 연속
준동형
$\omega:B\widehat{\otimes}_A C\to B$에 대응한다. 표준 준동형
$B\to B\widehat{\otimes}_A C$ 및
$C\to B\widehat{\otimes}_A C$를 이 유일한 준동형과 각각 합성하면 항등사상과
$\varphi$를 각각 얻는다. 그런데 $\omega$가
\emph{전사}임은 명백하므로 주장이 따른다.

\begin{corollary}[10.15.5]
\label{I.10.15.5-ko}
형식적 준스킴 $\mathfrak{S}$를 국소 뇌터라고 하고,
$\mathfrak{X}$를 $\mathfrak{S}$ 위의 유한형 형식적 준스킴이라 하자.
$\mathfrak{X}$가 $\mathfrak{S}$ 위에서 분리되어 있기 위한 필요충분조건은
대각 사상
$\mathfrak{X}\to\mathfrak{X}\times_{\mathfrak{S}}\mathfrak{X}$가
닫힌 몰입 사상인 것이다.
\end{corollary}

\begin{proposition}[10.15.6]
\label{I.10.15.6-ko}
국소 뇌터 형식적 준스킴들의 닫힌 몰입 사상
$j:\mathfrak{Y}\to\mathfrak{X}$는 분리 사상이다.
\end{proposition}

(\hyperref[I.10.14.2-ko]{10.14.2})의 표기를 사용하면
$j_0:Y_0\to X_0$는 닫힌 몰입 사상이고, 따라서 분리 사상이다. 이제
(\hyperref[I.10.15.2-ko]{10.15.2})를 적용하면 충분하다.

\begin{proposition}[10.15.7]
\label{I.10.15.7-ko}
$X$를 (보통 의미의) 국소 뇌터 준스킴, $X'$을 $X$의 닫힌 부분집합,
$\widehat{X}=X_{/X'}$라 하자. $\widehat{X}$가 분리되기 위한
필요충분조건은 $\widehat{X}_{\mathrm{red}}$가 분리되는 것이며,
특히 $X$가 분리되어 있으면 이 조건이 충족된다.
\end{proposition}

실제로, (\hyperref[I.10.8.5-ko]{10.8.5})의 표기를 사용하면
$\widehat{X}$가 분리되기 위한 필요충분조건은 $X'_0$가 분리되는
것이다(\hyperref[I.10.15.2-ko]{10.15.2}). 그런데
$\widehat{X}_{\mathrm{red}}=(X'_0)_{\mathrm{red}}$이므로,
$\widehat{X}_{\mathrm{red}}$가 분리된다고 말하는 것과 동치이다
(\hyperref[I.5.5.1-ko]{5.5.1, (vi)}).
